% ======================================================================
% Literal transcription of the 1994 manuscript
%   David Joyner, "On the transfer of Shalika germs for G_2" (summer, 1994)
% from its AMS-TeX source files into a single amsart LaTeX file.
%
% Claude: Transcription made by Claude (Anthropic), November 13, 2026.
% Claude: Source files, in the order used here:
%   g2new0.tex (Section 0), g2new1.tex (Section 1), g2new2.tex (Section 2),
%   g2new3.tex (Section 3), g2new4.tex (Section 4), g2new5.tex (Section 5),
%   g2new6.tex (Section 6), g2new7.tex (Section 7), g2new8.tex (Section 8),
%   g2new9.tex (Section 9), g2newapp.tex (Appendix), g2newref.tex (References).
% Claude: The files g2new5d.tex and g2new9d.tex are variant versions of
%   Sections 5 and 9. The body follows the unsuffixed files. Every passage in
%   which a "d" file differs is reproduced in an \iffalse ... \fi block or a
%   comment at the corresponding place, marked "VARIANT".
% Claude: Several source files contain text after \enddocument, which TeX never
%   typeset. That text is transcribed in \iffalse ... \fi blocks marked
%   "UNTYPESET". Change \iffalse to \iftrue to typeset a block.
% Claude: No mathematical or linguistic corrections have been made. Misprints,
%   spelling, hand-typed cross-references and section numbering are as in the
%   source. Additions by Claude appear only in comments that begin "Claude:".
% Claude: Macro collision. The source macro \off means \overline{F} in
%   g2new0, g2new1, g2new2 and g2new6, and {\Cal O}_F in g2new3, g2new4,
%   g2new5, g2new7, g2new8 and g2new9. Each occurrence below is rendered as
%   the original file printed it (\ovF or \OF). Places where the printed
%   meaning looks unintended are flagged in comments.
% Claude: Printed citation labels are those of the original ([Bor], [H1], ...).
%   Citation keys follow the house style (initials plus two-digit year).
%   Citations to labels that are missing from the original bibliography
%   ([Sp3], [Sp], [A], [H4], [D]) are typeset literally and flagged.
% Claude: newtx is not installed in the build environment, so Times is
%   provided by mathptmx, the fallback allowed by the style guide.
% ======================================================================
\documentclass[12pt,letterpaper]{amsart}
\usepackage[T1]{fontenc}
\usepackage{mathptmx}
\usepackage[letterpaper,margin=1in]{geometry}
\usepackage{amsmath,amssymb,amsthm,amscd}
\usepackage{tikz}
\usepackage{microtype}
\usepackage[hidelinks,bookmarks=false,pdftitle={On the transfer of Shalika germs for G2},pdfauthor={David Joyner}]{hyperref}

\newcommand{\ovF}{\overline{F}}
\newcommand{\OF}{\mathcal{O}_F}
\newcommand{\ooo}{\mathcal{O}}
\newcommand{\ovo}{\overline{\mathcal{O}}}
\newcommand{\aaa}{\mathbb{A}}
\newcommand{\zzz}{\mathbb{Z}}
\newcommand{\rrr}{\mathbb{R}}
\newcommand{\qqq}{\mathbb{Q}}
\newcommand{\gga}{\mathbb{G}}
\newcommand{\eps}{\epsilon}
\newcommand{\unx}{\underline{x}}
\newcommand{\iddots}{\mathinner{\mkern1mu\raise1pt\hbox{.}\mkern2mu%
  \raise4pt\hbox{.}\mkern2mu\raise7pt\vbox{\kern7pt\hbox{.}}\mkern1mu}}
\DeclareMathOperator{\Ind}{Ind}
% Claude: emergencystretch lets TeX reflow paragraphs that contain long
% inline formulas; it changes no content.
\setlength{\emergencystretch}{3em}

% Claude: AMS-TeX \proclaim{Name} ... \endproclaim is rendered by the
% environment "proclaim", which prints the name exactly as in the source.
\newcommand{\proclaimname}{}
\newtheorem*{proclaimx}{\proclaimname}
\newenvironment{proclaim}[1]%
  {\renewcommand{\proclaimname}{#1}\begin{proclaimx}}{\end{proclaimx}}

\title{On the transfer of Shalika germs for $G_2$}
\author{David Joyner}
\date{summer, 1994}

\begin{document}
\maketitle
% ======================================================================
% Claude: Section 0, source file g2new0.tex. Here \off = \overline{F}.
% ======================================================================
\section*{\S\,0 Introduction}

\subsection*{Motivation and rough outline}

To apply the trace formula to the transfer of automorphic representations
of two reductive groups $G$, $H$ over a number field $K$, one must,
among other things, have a comparison of local orbital integrals, for each
local completion $F$ of $K$. The comparison in the real and complex case
was solved by Shelstad \cite{She82}.
In the $p$-adic case, such a local comparison - the smooth matching conjecture
stated below -
boils down, thanks to the work of R. Langlands and D. Shelstad
(\cite{LS87}, \cite{LS90}), to a certain matching of Shalika germs
conjecture, also stated below.
Shalika germs are mysteriously behaving functions
associated to a semisimple element $s\in G(F)$ and a unipotent
conjugacy class ${\mathcal O}$ of $M=Cent_{G(F)}(s)$, defined in a punctured
compact-open neighborhood of $s$ in the regular set \cite{Sh72}. They play a
role in the analysis of $p$-adic orbital integrals on reductive groups
roughly analogous to the role that the power functions play in the Taylor
series of a function on the real line. By a well-known reduction known
as Harish-Chandra descent, it usually suffices to assume that $s$ is central
and $M=G$.
Langlands \cite{La83} introduced a geometric approach to the computation,
or at least comparison, of Shalika germs.
The comparison of Shalika germs associated to the
$0$-regular unipotent orbit was solved by Shelstad \cite{She89},
using results from \cite{LS87}, \cite{LS90}.
Langland's technique was further developed by
T. Hales in his thesis \cite{Ha92a}. Hales concentrated on the comparison of
$1$-regular germs, matching the stable subregular germs in general and
matching the $\kappa$-stable subregular germs for even rank unitary
groups. He also worked out a number of the computations necessary
to deal with the rank two case. This technique of Langlands goes by
the name Igusa theory since it was Igusa \cite{Ig78} who developed a crucial
asymptotic formula for an integral over fibers of a morphism
$Y_C\rightarrow C$, where $C$ is a curve and $Y_C$ is a non-singular
variety. This technique also relies on a formula for the coefficients
of this expansion, involving principal value integrals over $p$-adic
analytic manifolds, worked out in \cite{Ig78}, \cite{La83} (see also
\cite{LS83}). In \cite{La83}, \cite{LS89}, Langlands and Shelstad
matched all the germs for groups of absolute type $A_2$
(see also \cite{Ha92b}, where a descent argument is used instead of
the direct comparison in \cite{LS89}).
In \cite{Ha89}, Hales worked out the comparison of germs for the forms of
$B_2=C_2$ using Igusa theory and a reduction/descent argument.
In the present article, we verify the matching of Shalika germs conjecture
mentioned above for the $1$-, $2$- and $3$-regular germs for the
exceptional group $G_2$,
which, together with the result of \cite{She89}, proves the smooth
matching theorem for $G_2$. In other words, our main result
is the proof of the following

\begin{proclaim}{Smooth Matching Conjecture for $G_2$}
Let $T_G$ denote an elliptic Cartan of $G=G_2$
which has an $F$-embedding into $H$ and let $\kappa$ denote
a character of $H^1(Gal(\ovF/F),T_G)$ associated to a
cuspidal endoscopic group $H=H(T_G,\kappa)$
\cite{La83} (thus $H$ is either $PGL(3)$ or $(SL(2)\times SL(2))/\{\pm 1\}$).
Let $T_H$ denote the image of an $F$-embedding of $T_G$ into $H$.
For each $f\in C_c^\infty(G(F))$ there is an $f^H\in C_c^\infty(H(F))$
for which
\[
\Delta_{G/H}^{T_G,\kappa}(\gamma_H,\gamma_\lambda)
\Phi^{T_G,\kappa}(\gamma_G,f)
=\Phi^{T_H,st}(\gamma_H,f^H),
\]
whenever $\gamma_G$ is the image of a regular $\gamma_H\in T_H(F)$.
(The notation is explained further below.)
\end{proclaim}

\medskip
Before sketching the ideas in the paper, we describe what is new
or different from \cite{Ha89}, \cite{Ha92b}, which this paper
is comparable to. The notation will be explained later.

In \S 1, we carefully recall
in detail the construction of the Igusa variety and the
determination of local coordinates \cite{La83}, \cite{Ha92a}, \cite{Ha89}.

Section 2 gives the calculations for determining the action of the Galois
group on the local coordinates on the
Igusa variety for $G_2$ as in \cite{Ha89}.

Section 3 describes some new results: for any reductive group $G$,
we produce specific patches generalizing those used in \cite{Ha89}
and show for $G_2$
that they cover the base Igusa variety $Y''$ (described below) at least up to
codimension two. These specific patches are needed for doing certain
local calculations.
It is likely that these patches are regular under more
general conditions, which would imply $Y''$ is regular in codimension one
(conjectured in \cite{Ha92a}), but this is only proven in our
special case. Anyway, by another method
we show that, as a consequence of the theory of toroidal embeddings \cite{Ke72},
$Y''$ is normal and Cohen-Macauley (see Lemma 3.7 below).
This implies Hales' conjecture that
$Y''$ is regular in codimension one in general.

In section 4 we compute the Igusa constants of all
the base divisors on $Y''$ by an easy direct calculation
and give an ``endoscopic interpretation'' of some of the
calculations.
We also compute the Igusa forms using local coordinates of $Y_1$.

In section 5, we first record some basic facts about
the Igusa expansion over $Y_C$. Then we
prove the geometric ``no spurious divisors lemma'' for $G_2$
which is basically an inequality $e(\theta,\beta)\leq 1$,
where $e(\theta,\beta)$ is as in \cite{La83}, which allows one to work
on $Y''$ rather than its desingularization $Y_C$.
Finally, we recall the statement of Hales' important ``reduction lemma''.

In sections 6, 7, and 8
we (a) compute coordinates
on the $1$-, $2$- and $3$-regular unipotent orbits (the latter two
are not Richardson which is something new, though they are what I call
``endoscopically Richardson''), (b) compute
coordinates on the fiber $\phi^{-1}(\gamma_\lambda)^\Gamma$
(explained below), and (c) compute the Igusa form
on the fiber and its explicit dependence on a tangent vector $X$
of $C$. Not much in these
sections is substantially different from \cite{Ha89}, \cite{Ha92b}. One
new idea is that, roughly speaking, the subgroup generated by the
one-parameter root-groups associated to roots ``outside $H$''
may be used to determine local coordinates for the
``endoscopically Richardson'' orbit $\ooo=\Ind_H^G(1)$.

The last section presents a conjecture which our method of proof
seems to support. It seems to be the ``ultimate'' theorem which one
might prove by this method without introducing any new ideas.

For the most part, the appendices collect explicit
group-theoretic computations
needed to justify a number of claims made in the main body of the paper.
For example, we verify that the cuspidal endoscopic groups of $G_2$ are
$PGL(3)$ and $(SL(2)\times SL(2))/\{\pm 1\}$.

Turning now to a brief discussion of the constructions we use,
here is a rough sketch of Langlands' idea in its simplest form:
Let $\gamma\in G(F)$ be an elliptic regular element and let
$T_G=Cent_G(\gamma)^o$. The goal is to compute,
or at least find some manageable expression for, the stable orbital
integral $\int_{(T_G\setminus G)(F)}f(g\gamma g^{-1})dg$.
(More generally, one can and does consider $\kappa$-stable orbital
integrals.) The variety $T_G\setminus G$
may be identified, over the algebraic closure
$\overline{F}$, with the variety of regular stars: let
$\Omega$ denote the absolute Weyl group of $T_G$,
$B_G$ a Borel subgroup of $G$
containing $T_G$, and $S_0$ denote the $G(\ovF)$-orbit of the
$|\Omega|$-tuple $(B_G^\omega)_{\omega\in \Omega}$, where
$G$ acts by conjugation on each coordinate. This orbit is called the
variety of regular stars and since
$\cap_{\omega\in \Omega}N_G(B_G^\omega)=T_G$,
we have $S_0\cong T_G\setminus G$ over $\ovF$. Let
${\mathcal B}\cong B_G\setminus G$
denote the Borel variety. Although ${\mathcal B}$ is complete,
$S_0$ is not but
conjugation induces a principal $N_G$-bundle structure on
$S_0\rightarrow {\mathcal B}$, where $N_G$ denotes
the unipotent radical of $B_G$,
so $S_0$ may be roughly regarded as
the base space ${\mathcal B}$ with fibers at each point isomorphic to
$N_G$. The Zariski closure of
$S_0$ in ${\mathcal B}^{|\Omega|}$ is denoted $S$ and is called the variety of
stars. It is projective, complete, but contains singularities.
It turns out that $S_0$ is Zariski open in $S$.
Roughly speaking, the idea is to regard the above stable orbital integral
as an integral over $S_0(F)\subset S(\ovF)$. Let $T_G'$
denote the subvariety of regular elements of $T_G$. We
restrict $\gamma$ to a regular curve $C$ in $T'_G\cap \{1\}$
in a neighborhood of the identity and then construct a smooth variety $Y_C$
with a fibering $\phi:Y_C\rightarrow C$, ``the''
smooth Igusa variety, whose fibers
$\phi^{-1}(\gamma)$, $\gamma\in C,\ \gamma\not= 1$,
are isomorphic to $T_G\setminus G$. In fact, the above stable orbital integral
may be regarded now as a fiber integral parameterized by a
local coordinate $\lambda$ in a neighborhood of the singular point $1\in C$.
Taking the limit as $\lambda$ tends to this singular point,
one may then use techniques of Igusa \cite{Ig78}, as developed
by Langlands \cite{La83}, to express the stable orbital integral
as a sum involving principal value integrals over certain divisors
contained in $\phi^{-1}(1) \subset Y_C$. By the homogeneity
of the Shalika germs and the
homogeneity of the coefficients in the Igusa expansion, only certain
divisors can contribute to a given Shalika germ.

Here is a rough sketch of Igusa theory \cite{La83}, \cite{Ha92a}:
Assume $G$ is $F$-quasisplit. Let $C$ denote a $1$-parameter subgroup of $G$
in $T_G$, defined over $F$, containing the identity
and no other singular point and whose tangent over $\ovF$ is not
contained in any root hyperplane. Let $C'=C-\{1\}$. We sketch the
construction of a non-singular algebraic variety $Y_C$, defined over $F$,
for which
\[
\begin{CD}
C'\times T_G\setminus G @> i >> Y_C @> \pi >> G\\
@. @VV \phi V @. \\
@. C, @.
\end{CD}
\]
with the following properties:

\noindent
(a) if $T_G$ and $C$ are defined over $F$ then so are
$Y_C$, $\phi$, $\pi$, and $i$;

\noindent
(b) $\pi : Y_C\rightarrow G$ is proper and $G$-equivariant;

\noindent
(c) $\phi : Y_C\rightarrow C$ is proper;

\noindent
(d) $\pi(\phi^{-1}(1))=V$, the unipotent variety of $G$;

\noindent
(e) if ${\mathcal O}\subset G$ denotes a
unipotent orbit and if $E=\pi^{-1}(u)$,
for some $u\in \pi(Y_C)\cap {\mathcal O}$,
then $\pi|_{E}:E\cap \pi^{-1}({\mathcal O})\rightarrow {\mathcal O}$
is faithfully flat
\cite[Ch II, Proposition 3.1]{DG80} and, in particular, surjective;

\noindent
(f) $i(C'\times T_G\setminus G
)\subset Y_C$ is open;

\noindent
(g) $\pi\circ i(\gamma,g)=\gamma^g$ and $\phi\circ i(\gamma,g)=\gamma$.

\noindent
(h)  $Y_C$ is $2p+1$-dimensional, where $p=|\Sigma^+|$
denotes the number of positive roots.

\noindent
A divisor $E\subset Y_C$ for which $\pi(E)=\ovo$ is called an
\textbf{$\ooo$-divisor}. In general, although $E$ is smooth,
$\ovo$ need not even be normal \cite[\S 0]{KP79}.

\noindent
Step 1: Let $X$ denote the Zariski closure of $T'\times T_G\setminus G$
in $G\times S$ and let $Y$ denote the Zariski closure of
$C'\times T_G\setminus G$
in $X$. Let $\chi :G\rightarrow T_G/\Omega$ denote the adjoint
quotient \cite{Sl80}, where $\Omega$ denotes the Weyl group of $T_G$.
We will see below that $X$ is covered by patches
$X(B_+ ,B_{-})\subset
B_{-}\times S\times N_+$ and $Y$ is covered by patches
$Y(B_+ ,B_-)\times N_+\hookrightarrow \chi^{-1}(C) N_-\times S\times N_+$,
where $N_-,N_+$ denotes the unipotent radical of $B_-,B_+$, resp..
Here $B_+,B_-$ are opposite Borels for which
$T_+=B_+\cap B_-$ is defined over $F$. The completeness
of $S$ implies that the projections $Y\rightarrow G$
and $X\rightarrow G$ are proper.

\noindent
Step 2: We construct below a resolution $S_1$ of $S$ with a proper
morphism $S_1\rightarrow S$ and an injective morphism
$S_0\rightarrow S_1$. The image of $S_0$ in $S_1$ is denoted $S''$.
Both $S$ and $S_1$ are complete but $S''$ is not.
Let $X''$, $X_1$ denote the Zariski closure of $T'\times T_G\setminus G$
in $G\times S''$, $G\times S_1$, respectively.
Let $Y''$, $Y_1$ denote the Zariski closure of $C'\times T_G\setminus G$
in $X''$, $X_1$, respectively. Again, we have proper maps
$Y_1\rightarrow G$ and $X_1\rightarrow G$. Though $Y''\rightarrow G$
may not be proper, it is especially useful because
of its relatively simple local coordinate relations. We can explicitly
desingularize $Y''$ for $G=G_2$ by blowing up four
times in succession.
(The derivation of this is lengthy and not really needed for the
proof of our main result so is omitted.)
Recall that the blow down morphisms
are proper and defined over $F$. They are also $G$-equivariant
by the universal property of blowing up \cite{Har77}.

\noindent
Step 3: Resolve the singularities of $Y_1$ contained in
$Y_1-Y''$ by using Mumford's semi-stable reduction theorem \cite{Ke72}.
This yields a variety $\tilde{Y_1}$ and a blow down
$\xi:\tilde{Y_1}\rightarrow Y_1$ such that all the singularities
of $\tilde{Y_1}$ map down to $Y_1- Y''$ under $\xi$. Now
perform an embedded resolution of singularities on
$Y''\hookrightarrow \tilde{Y_1}$ as in step 2. The resulting
resolution of $\tilde{Y_1}$
is denoted $Y_C$. The blow down map $Y_C\rightarrow Y_1$
is proper, so $\pi$ and $\phi$ are proper.

We remark that the properness of $\pi$ and
$\phi$ are needed to guarantee that the appropriate
principal value integrals make sense. Properness implies that
the principal value integral of the
pull-back of the function $f_\kappa$ is well-defined
and that some of the hypotheses to Lemma 5.7 below hold.

\subsection*{Notation and detailed outline}

Let $\gamma_\lambda=exp(\lambda X)\in C'(F)\subset T_G(F)$ be the image of
$\gamma_{\lambda,H}\in T_H(F)\subset H(F)$ in the sense of \cite{LS90}.
After some general background, we will examine the $1$-, $2$-,
and $3$-regular germs for $G_2$ using the formula (\cite{Sh72}, \cite{LS90})
\[
\begin{split}
\Delta_{G/H}^{T_G,\kappa}(\gamma_{\lambda,H},\gamma_{\lambda})
\Phi^{T_G,\kappa}(\gamma_{\lambda},f)
&=\kappa(inv(\gamma^0_{G},\gamma_{\lambda}))
D_H(\gamma_{\lambda,H})^{-1}
\int_{\phi^{-1}(\gamma_{\lambda})^\Gamma}\pi^*(f)
\kappa(t_*)\frac{|\omega|}{|\phi^*(d\lambda)|}\\
&=\sum_{{\mathcal O}}\mu_{\mathcal O}(f)
\Gamma_{{\mathcal O}}^{T_G,\kappa}(\gamma_{\lambda}),
\ \ \ f\in C_c^\infty(G),
\end{split}
\tag{0.1}
\]
where $\Delta_{G/H}^{T_G,\kappa}$ is the local transfer factor given in a
neighborhood of the identity by
\[
\Delta_{G/H}^{T_G,\kappa}(\gamma_{\lambda,H},\gamma_G)
=\kappa(inv(\gamma^0_G,\gamma_G))D_G(\gamma_G)/D_H(\gamma_H)
\]
as in \cite{LS90}, $\gamma_G^0$ being a fixed regular element of $T_G(F)$.
The sum over ${\mathcal O}$ in (0.1) runs over all stable unipotent orbits
${\mathcal O}^{st}$ in $G(F)$, ${\mathcal O}\subset {\mathcal O}^{st}$
denoting an arbitrarily chosen $G(F)$-orbit within ${\mathcal O}^{st}$.
Also, in (0.1),
$\gamma_\lambda=\gamma_{\lambda,G}\in C\cap V_f\subset T_G$
has parameter $\lambda$,
so that $\gamma_0=1$, where $V_f\subset G_{reg}$
is a punctured neighborhood of $1$,
$t_*$ is a cocycle-valued function which will be defined following
0.6 below, and
$\phi^{-1}(\gamma_\lambda)^\Gamma$ denotes the fixed points of the fiber
under the action of $\Gamma:=Gal(\overline{F}/F)$ on $Y_C$ determined
by $T_G$ as described in \S 2 below.
The measure
$\frac{|\omega|}{|\phi^*(d\lambda)|}$ incorporates the factor
$D_G(\gamma_{\lambda,G})$.
If $\gamma_G'\in G(F)$ is stably conjugate to $\gamma_G$ then
$\Delta_{G/H}^{T_G,\kappa}(\gamma_H,\gamma'_G)
=\kappa(inv(\gamma_G,\gamma'_G))
\Delta_{G/H}^{T_G,\kappa}(\gamma_H,\gamma_G)$.
We will normalize
the measures and use the same notation as in \cite{LS89}. In particular,
\[
\Gamma_{{\mathcal O}}^{T,\kappa}(\gamma)=
\Delta_{G/H}^{T,\kappa}(\gamma_H,\gamma)\sum_\delta \kappa(\delta)
\Gamma_{\mathcal O}(\gamma^\delta),
\]
if $\gamma$ is an image of some $\gamma_H$. Here $\delta$ runs over
classes in $H^1(\Gamma,T_G)$ associated to the distinct stable
conjugacy classes $\{\gamma^\delta\}$ of $\gamma$. These germs
$\Gamma_{\mathcal O}(\gamma)$ are slightly different from those in \cite{Sh72}
since the germs here do not include the discriminant factor
in their definition. By a result of Harish-Chandra, if
${\mathcal O}$ is $r$-regular then $\Gamma_{\mathcal O}$ satisfies
\[
\Gamma_{\mathcal O}(exp(\lambda^2 X))=|\lambda|^{-dim\ {\mathcal O}}
\Gamma_{\mathcal O}(exp( X)),
\]
where for any $u\in {\mathcal O}$,
$dim\ {\mathcal O}=dim\ G_u\setminus G=dim\ G -rk(G)-2r=2|\Sigma^+_G|-2r$ with
$G_u$ denoting the centralizer of $u$ in $G$,
$r=dim\ {\mathcal B}_u$, ${\mathcal B}$ denoting the Borel variety,
${\mathcal B}_u$ its subvariety fixed by the conjugation action of $u$
and $\Sigma^+_G$ denoting the positive roots with respect to some Borel.
Note that since
\[
D_G(exp(\lambda X))=\prod_{\sigma >0}(1-\sigma(exp(\lambda X))^{-1})
=\lambda^{|\Sigma^+_G|}\prod_{\sigma >0}\sigma(X)+\text{ higher order terms},
\]
the normalized germ $\Gamma^*_{\mathcal O}(\gamma)
:=|D_G(\gamma)|\Gamma_{\mathcal O}(\gamma)$ satisfies
$\Gamma^*_{\mathcal O}(exp(\lambda^2 X))
=|\lambda|^{2r}\Gamma^*_{\mathcal O}(exp( X))$.

The ultimate object of our desires is the following

\begin{proclaim}{Matching of Germs Conjecture \cite{LS89}}
Let $H$ be an endoscopic group
associated to $T,\kappa$. Given an $r$-regular
unipotent orbit ${\mathcal O}$ in $G(F)$ there is a finite collection
of $r$-regular unipotent classes ${\mathcal O}_i$ in $H(F)$ and complex
numbers $a_i$ such that, for every Cartan $T_H$ of $H$ over $F$ and every
regular ray $\{\lambda X|\lambda\not=0\}$ in the Lie algebra ${\mathfrak t}_H$,
we have
\[
\Gamma_{{\mathcal O}}^{T,\kappa}(\gamma_\lambda)=
\sum_i a_i\Gamma_{{\mathcal O}_i}^{T_H,st}(\gamma_\lambda ),
\]
for $\lambda$ sufficiently small.
\end{proclaim}

Our goal is to try to verify this for $G=G_2$ and $r=1,2,3$.

The expression in (0.1) may also be written in
the form $1/D_H(\gamma_{\lambda,H})$
times
\[
F(\lambda)=\int_{\phi^{-1}(\lambda)(F)}f_\kappa(y)|\omega_\lambda|,
\tag{0.2}
\]
where $f_\kappa$, as defined by (0.1) and (0.2),
satisfies the conditions of being an Igusa integrand
in the sense of \cite{La83}. Igusa's expansion
\cite{La83} is an asymptotic expression of the form
\[
F(\lambda) =\sum_{\theta,\beta}\theta(\lambda)|\lambda|^{\beta -1}
\sum_{r=1}^{e(\theta,\beta)}(-\log_q|\lambda|)^{r-1}F_r(\theta,\beta,f_\kappa),
\tag{0.3}
\]
where $\theta$ is a unitary character of $F^{\times}$ and $\beta >0$
is a rational number for which ${\mathcal E}(\theta,\beta)\not= \emptyset$.
Here ${\mathcal E}(\theta,\beta)$ denotes the set of irreducible components
(called \textbf{branches}) in $\phi^{-1}(1)\subset Y_C$ such that
$\theta^{a(E)}=\kappa_E$ and $\beta=b(E)/a(E)$, where $a(E)$, $b(E)$,
$\kappa_E$, and $e(\theta,\beta)$ are as in \cite{La83}. As T. Hales
has observed, since there are no logarithmic terms in the
Shalika germ expansion, we have

\begin{proclaim}{Lemma 0.4} If $f_\kappa$ is as in (0.2) then
$F_r(\theta,\beta,f_\kappa)=0$
for all $r>1$.
\end{proclaim}

By Howe, Harish-Chandra, and
Rogawski \cite{Ro80}, the trivial germ is a non-zero constant $c_H$.
We describe how we will be investigating the other
terms in (0.1) associated to the
$1$-regular, $2$-regular,
and $3$-regular orbits ${\mathcal O}$.
These will match with the trivial orbit in the germ expansion on $H$,
\[
\Phi^{T_H,st}(\gamma_{\lambda,H},f^H)
=\int_{\phi_H^{-1}(\gamma_{\lambda,H})^\Gamma}\pi^*(f^H)
\frac{|\omega|}{|\phi^*(d\lambda)|}
=\sum_{{\mathcal O}}\mu_{\mathcal O}(f^H)
\Gamma_{{\mathcal O}}^{T_H,st}(\gamma_{\lambda,H}),
\tag{0.5}
\]
where $H=(SL(2)\times SL(2))/\{\pm 1\}$ or $H=PGL(3)$, and where
$f^H\in C_c^\infty(H)$ exists and
$\Delta^{T_G,\kappa}_{G/H}$-matches with $f$.
These terms are the $\beta =2,3,4$ terms in the Igusa expansions
for $G_2$, $H_1=PGL(3)$ and $H_2=(SL(2)\times SL(2)))/\{\pm 1\}$.
It ultimately turns out that, in these expansions, we have
$e(\theta,\beta)\leq 1$ (this is proved in \S\S 4-5 below).
Therefore, formula 0.6 below applies to our situation.

Putting together results of \cite{La83}, \cite[\S 2]{Ha89},
using the Shalika germ expansion, we
have the following

\begin{proclaim}{Formula 0.6(a)} Assume that we have a Cartan $T_G$, a character
$\kappa$ of $H^1(\Gamma,T_G)$, and
a $(\beta-1)$-regular
stable unipotent class ${\mathcal O}^{st}$ for which there is only
one $F$-rational irreducible branch $E_\beta$ of $\{\lambda =0\}$ in
$Y_C$ with $\pi(E_\beta)=\overline{{\mathcal O}}^{st}$ and
$\beta(E_\beta)=\beta$. Then we have
\[
D_H(\gamma_{H,\lambda})
\sum_{{\mathcal O}\subset {\mathcal O}^{st}}
\mu_{{\mathcal O}}(f)\Gamma_{{\mathcal O}}^{T_G,\kappa}(\gamma_\lambda)
=\kappa(inv(\gamma_G^0,\gamma_\lambda))
|\lambda|^{\beta -1}\sum_{\theta}\theta(\lambda)
\int_{E_\beta(F)} \pi^*(f)\kappa_\theta(t_*)|\omega_{E_\beta}|,
\]
for $\gamma_\lambda\in C\cap V_f$. Here $\kappa_\theta(t_*)$
is an extension of $\kappa(t_*)/\theta(\lambda)$ to
$E_\beta$ and $\overline{{\mathcal O}}^{st}$ denotes the $F$-Zariski closure
of the stable orbit and $t_*$ is a cocycle-valued function
described more below, and the sum on the left runs over all
unipotent $G(F)$-orbits ${\mathcal O}$ contained in the
stable orbit ${\mathcal O}^{st}$.

\noindent
0.6(b): Same assumptions. Then for each $\overline{\mathcal O}^{st}$
there is a measure $\omega_{\overline{\mathcal O}^{st}}$ and for each
$F$-rational fiber
$E_\beta(u)=E_\beta(F)\cap \pi^{-1}(u)$,
$u\in \overline{\mathcal O}(\overline{F})$,
there is a measure
$\omega_{E_\beta(u)}$ such that
\[
\int_{E_\beta(F)} \pi^*(f)\kappa_\theta(t_*)|\omega_{E_\beta}|
=\int_{\overline{\mathcal O}^{st}}f\int_{E_\beta(u)} \kappa_\theta(t_*)
|\omega_{\overline{\mathcal O}^{st}}||\omega_{E_\beta(u)}|,
\]
for each $\theta$ occuring above.
\end{proclaim}

This will follow from the results of section 5 below.
By \cite[ch 6, \S 2]{Ha92a}, only the $\theta =1$
term on the left side of formula 0.6(a) contributes a non-zero
term to the sum.
It turns out that either $\kappa$ is
quadratic (and $H=(SL(2)\times SL(2))/\{\pm 1\}$)
or $\kappa$ is a cubic character (and $H=PGL(3)$).

Now we turn to recalling some of the properties of $t_*$ \cite{La83},
\cite{Ha92a}. Recall $\Gamma=Gal(\overline{F}/F)$.
The cocycle-valued function $t_*$ is a map
$Y_C\rightarrow Z^1(\Gamma, T_G(\overline{F}))$ given by
$P\longmapsto (\sigma\longmapsto t_\sigma(P))$
which has the following properties:

\noindent
(a) it is locally constant on the $p$-adic manifold $Y_C(F)$ \cite{La83};

\noindent
(b) it factors through $Z^1(Gal(E/F), T_G(E))$,
where $E$ denotes the Galois closure of the splitting
field of $T_G$;

\noindent
(c) for fixed $\sigma\in \Gamma$, $t_\sigma$ is a regular on the variety
$Y_C$;

\noindent
(d) for $y=i(\gamma_\lambda ,g)\in Y_C$, such that for all
$\sigma\in \Gamma$,
$\sigma(g)g^{-1}$ belongs to $T_G(\overline{F})$, we have
$\kappa (t_*(y))=\kappa(\sigma(g)g^{-1})$, where
$\sigma(g)g^{-1}$ represents
the cocycle $\sigma\longmapsto \sigma(g)g^{-1}$.

After $Y_C$ has been constructed, the idea is to determine the sets
${\mathcal E}(\theta, \beta)$ with $\beta =2,3,4$. First,
one determines the $\ovF$-branches in $\phi^{-1}(1)\subset Y_C(\ovF)$ which
map down to $Y''$ under $Y_C\rightarrow Y_1$ which have
$\beta =2,3,4$. These are called \textbf{fundamental divisors}.
In fact, in the case of $G=G_2$, we can construct on $Y''$ itself
divisors occurring in $\phi^{-1}(1)\subset Y''$.
Such divisors will be called \textbf{base divisors}.
Of course, a base divisor is fundamental.
We sometimes identify a base divisor with its closure in $Y_C$.
Then one determines which of these fundamental divisors
has $F$-rational
coordinates with respect to the Galois action determined by $T_G$.
These branches will contribute to Igusa's
expansion. Next, one must consider the $\ovF$-branches in
$\phi^{-1}(1)\subset Y_C(\ovF)$ which do not
map down to $Y''$ under $Y_C\rightarrow Y_1$. These are
called \textbf{spurious divisors} and are relatively complicated to deal with.
It turns
out that the computations in \cite[Ch III]{Ha92a} show that there are
no $F$-rational spurious divisors contributing to the subregular germ
for $G_2$. For other reasons, the spurious divisors can be avoided
altogether for the $2$-regular and $3$-regular germs.
In fact, we have the following result, proven in \S\S 4-5 below.

\begin{proclaim}{Lemma 0.8} Let $E$ be an $F$-rational
base ${\mathcal O}$-divisor of the
variety $Y''$ and let $\tilde{E}$ denote its strict transform
in $Y_C$. The principal
value integrals $\int_{{E}(F)}\pi^*(f)\kappa_1(t_*) |\omega_{{E}}|$ and
$\int_{\tilde{E}(F)}\tilde{\pi}^*(f)\kappa_1(\tilde{t}_*) |\omega_{\tilde{E}}|$
exist. Here $\kappa_1$ is as in formula 0.6(a), $t_*$ is the
cocycle function on $Y''$, and $\tilde{t}_*$ is the cocycle function
on $Y_C$. Moreover $e(1,\beta_E)=1$ and
$F_1(1,\beta_E,f_\kappa)$ is given
by $\kappa(inv(\gamma_G^0,\gamma_G))$ times
\[
\begin{split}
PV\int_{{E}(F)}\pi^*(f)\kappa_1(t_*) |\omega_{{E}}|
&=PV\int_{\tilde{E}(F)}\tilde{\pi}^*(f)
\kappa_1(\tilde{t}_*) |\omega_{\tilde{E}}|\\
&=\int_{\overline{\mathcal O}^{st}(F)}
f\int_{{E}_u(F)}\kappa_1(t_*)|\omega_{{E}_u}|
|\omega_{\overline{\mathcal O}^{st}}|,
\end{split}
\]
where the measures are as in (0.6).
\end{proclaim}

As a consequence of this, it will follow that
spurious divisors do not contribute to the Igusa
expansion of (0.1) or (0.2).

The proof that $e(\theta,\beta)\leq 1$ in the case of
$G_2$ will now be sketched. From (0.1) we know that the
germs associated to the
$(\beta -1)$-regular unipotent orbits ${\mathcal O}$ get
``matched'' with the terms in the Igusa expansion of
homogeniety degree $\beta -1$ in $\lambda$. Since both sides depend
on $X$ and $\lambda$ we might write such a comparison briefly
as
$
\Gamma_\beta (\lambda,X,f)=I_\beta(\lambda,X,f),
$
the left side denoting the sum of the normalized $\kappa$-stable
germs of homogeniety degree $\beta -1$ and the right side
denoting the corresponding term in the Igusa expansion
(0.2), (0.3). If we replace $X$ in both sides by $cX$ then both sides
should have the same degree of homogeniety in $c$ as well. However,
by computing explicitly the dependence of the term $I_\beta(\lambda,X,f)$
in the Igusa expansion on $X$ we will see, using the case-by-case
computation in \S 5, that $e(\theta,\beta)\geq 2$ implies that the two sides
have different degrees of homogeniety in $c$, a contradiction.

It remains then to analyze the integrals
over the fundamental divisors. We cannot compute these integrals
but, in most cases, we can determine their dependence on $\lambda$
and the tangent direction of $C$, and this is enough for our purposes.
On the right-hand side of
formula 0.6(a) we find what is essentially the factor
\[
\Delta_{0}(X)=\lim_{\lambda\rightarrow 0}
\frac{{\Delta_{loc}(\gamma_{\lambda,G},\gamma_{\lambda,H})}}{D_{G/H}(\gamma_{\lambda,G},\gamma_{\lambda,H})},
\]
in the notation of \cite[\S 2]{Ha92b}.
The local translation invariance property in the hypotheses to
Hales' reduction lemma 5.19 below holds for
$\kappa(t_*)\Delta_0(X)$ as we will see (this is carried out in \S\S6-8
below, which we compute a quantity which is essentially
$\Delta_0(X)$).
Therefore, we leave the $\kappa$-stable ${\mathcal O}$-germ
unchanged if we replace $X$ by $X+Y$, where $Y$ is sufficiently small
and as in lemma 5.19.
We may now apply lemma 5.18 to obtain the matching for functions
whose support meets the stable orbit ${\mathcal O}^{st}(F)$ but no
other unipotent orbit.
(This is where \cite{LS90} is implicitly used.)
This descent argument, remarkably, allows us to
avoid doing the difficult, more detailed calculations necessary for a
direct determination of the dependence of
$\Delta_{0}(X)\kappa(t_\sigma)$ on $X$.

As we've already mentioned, the $0$-regular orbit has already been matched
\cite{She89}.
The $1$-regular $\kappa$-stable germ of $G_2$ matches with the
$1$-regular stable germ of $PGL(3)$ on cubic Cartans and with the
$1$-regular stable germ of $(SL(2)\times SL(2))/\{\pm 1\}$
on biquadratic Cartans.
The stable germ associated to the
trivial unipotent orbit of $(SL(2)\times SL(2))/\{\pm 1\}$
matches with the 2-regular $\kappa$-stable
germ of $G_2$ on $T$, where $T$ and $\kappa$ determine
$H_{T,\kappa}=(SL(2)\times SL(2))/\{\pm 1\}$. The stable germ associated to the
trivial unipotent orbit of $PGL(3)$ matches with the 3-regular $\kappa$-stable
germ of $G_2$ on $T$, where $T$ and $\kappa$ determine $H_{T,\kappa}=PGL(3)$.
These matchings are in accordance with the fact that the $2-$regular
and $3-$regular unipotent orbits of $G_2(\overline{F})$
are endoscopically induced from the trivial orbits of $A_1\times A_1$
and $A_2$, respectively \cite[ch. III, table on page 221]{Spa82}.
Finally, the matching for the  trivial orbit follows from the work
of Rogawski, Howe, and Harish-Chandra \cite{Ro80}.

Though I've tried to make this paper self-contained (outside of
the usual background in group theory \cite{Bor91} and
algebraic geometry \cite{Mu88}), familiarity
with the ideas in \cite{La83}, \cite{LS89}, \cite{Ha92a}, and \cite{Ha92b} is assumed.

\medskip
Acknowledgement: I'd like to thank Magdy Assem, Dan Bump,
Caroline Grant, Roland Martin, and especially Tom Hales for many
enlightening conversations on the subject of this paper.
% ======================================================================
% Claude: Section 1, source file g2new1.tex (\pageno 10). Here \off = \overline{F}.
% ======================================================================
\section*{\S\,1 Some Background}

Let $G$ denote a connected semi-simple
quasi-split algebraic group defined over
$F$ and let $T_G$ denote a maximal torus of $G$ defined
over $F$.
In fact, we will assume that $T_G$ is an elliptic Cartan
but we will not make use of this immediately.
We shall work over a separable closure $\overline{F}$
most of this section.
A \textbf{root} is a character
$\rho\in X^*(T)$ of a torus $T$ such that there exists an isomorphism
$\theta_\rho :\gga_a\rightarrow U_\rho$, $U_\rho\subset G$
a closed subgroup, for which $t^{-1}\theta_\rho(\xi)t=
\theta_\rho(\rho(t)\xi)$, for $t\in T$ and $\xi\in\gga_a$.
Here $\gga_a$ denotes the 1-dimensional additive affine group.
Let $B_G$ denote a Borel containing $T_G$,
$\Sigma_G$ a set of roots of $T_G$, $\Sigma_G^+$ the set of
positive roots with respect to $B_G$,
and $\Delta$ the set of simple roots in $\Sigma_G^+$.
When $G=G_2$, we usually write
$\Delta=\{\alpha,\beta\}$ and
\[
\Sigma_G^+=\{\alpha,\beta,\gamma=\alpha+\beta,
\delta=2\alpha+\beta,\epsilon=3\alpha+\beta,
\zeta=3\alpha+2\beta\}.
\]

Fix a Borel $\underline{B}$ of $G$ defined over ${F}$, a Cartan
$\underline{T}$ defined over $F$ contained in $\underline{B}$,
and a positive Weyl chamber $W_0$ of $T_G$ with respect to
$\Sigma_G^+$.

Let $\Omega$ denote the absolute Weyl group of $T_G$.
We may occasionally identify the roots of $G$ with respect to $T_G$
with those of $G$ with respect to $\underline{T}$ by means of
a fixed isomorphism $T_G\rightarrow \underline{T}$, and, thus, identify
the absolute Weyl group of $\underline{T}$ with that of $T_G$.
Representing $\omega\in \Omega$ by
$w\in N_{G(\overline{F})}({T}_G(\overline{F}))$,
define
\[
{B}_G(W):=w^{-1}{B}_Gw={B}_G^w,
\]
for $W=\omega^{-1}W_0$.
The Borels
$\{{B}_G(W))\ |\ W\in {\mathcal W}\}$ are precisely the Borel
subgroups of $G$ which contain ${T}_G$, where ${\mathcal W}$ denotes
the set of Weyl chambers of $T_G$.

Recall that a \textbf{regular star} is a $|\Omega|$-tuple of Borels
$(B(W))_{W\in {\mathcal W}}$ such that, for all $W\in {\mathcal W}$,
$B(W)={B}_G(W)^h=h^{-1}{B}_G(W)h$, for some
$h\in G(\overline{F})$ (independent of $W$).
The variety of regular stars is denoted by $S_0$.
There is a natural bijection between $T_G(\ovF)\setminus G(\ovF)$
and $S_0(\ovF)$ given by $g\longmapsto (B(W)^g)_{W\in {\mathcal W}}$.

Let
\[
X_0:=\{(g,(B(W))_{W\in {\mathcal W}})\in G\times S_0\ |\
g\in\cap_{W\in {\mathcal W}}B(W)\text{ is regular}\},
\tag{1.1}
\]
so that we have an $\overline{F}$-isomorphism
\[
T_G'\times T_G\setminus G
\rightarrow X_0
\tag{1.2}
\]
\[
(t,gT_G)\longmapsto (g^{-1}tg,{B}_G(W)^g).
\]

\begin{proclaim}{Definition} Let
$\phi:X_0\rightarrow T_G$ be the map which sends
$(g,s)\in {\mathcal C}(T_G)\times S_0$ to the point in $T_G$ which meets the
conjugacy class of $g$, where ${\mathcal C}(T_G)$ denotes the
union of the conjugates of $T_G$.
Let $\pi :X_0\rightarrow G$ denote the projection map onto
the ``$G$-coordinate''.
\end{proclaim}

The map $\pi$ restricts to an isomorphism onto the
conjugacy class of $\gamma$,
\[
\pi:\phi^{-1}(\gamma)\rightarrow \{g\gamma g^{-1}\ |\ g\in G/T_G\}.
\]

\subsection*{Construction of $Y_1$}

Let $C\subset T_G$ be a one-parameter subgroup of $G$
containing $\{1\}$ but no other singular point.
Let $C':=C-\{1\}$
and let $Y_0:=\phi^{-1}(C')\subset X_0$. The isomorphism
(1.2) restricts to an $\overline{F}$-isomorphism
\[
C'\times T_G\setminus G\rightarrow Y_0.
\tag{1.4}
\]
It can be shown that $Y_0$ is a regular algebraic subvariety of $G\times S$
\cite[p.\ 480]{La83} but is not complete.
Let $Y$ denote the closure of $Y_0$ in $G\times S$. $Y$ is
singular in general.

Let ${\mathcal B}\cong B_G\setminus G$ denote the Borel variety.
It will be helpful
later to record here local coordinates for ${\mathcal B}$.
First we recall the following
well-known result \cite[ch IV, Corollary 14.13]{Bor91}:

\begin{proclaim}{Lemma 1.5} If $B,B',B''\in {\mathcal B}$ and $B', B''$ are opposite
to $B$ then they are conjugate by an element of $B$.
\end{proclaim}

As a consequence of the Levi decomposition,
it follows that, for $B, B',B''$ as in the lemma,
$B'$ and $B''$ are conjugate by an element of $N$, the unipotent radical of
$B$. Let ${\mathcal B}(B)\subset {\mathcal B}$
denote the open subvariety of all Borels
opposite to $B$.

\begin{proclaim}{Corollary 1.6} For each $B\in {\mathcal B}$, ${\mathcal B}(B)$
is isomorphic
as an affine variety to $N$ and the coordinates on $N$ (i.e., the matrix
coefficients) serve as coordinates on ${\mathcal B}(B)$.
\end{proclaim}

A point
$(B(W))_{W\in{\mathcal W}}\in {\mathcal B}^{|\Omega|}$ is called a \textbf{star}
if for every $\omega\in\Omega$ and every simple root $\alpha$
of $T_G$ in $B_G$ there is an $h\in G(\overline{F})$
such that
\[
B(\omega^{-1}W_0)=h^{-1}B_Gh
\]
and
\[
B((s_\alpha\omega)^{-1}W_0)\subset h^{-1}P_{G,\alpha} h
\]
where ${P}_{G,\alpha}$ is the smallest parabolic subgroup
of $G$ containing $B_G$ and $B_G^{s_\alpha}$.
Here $s_\alpha$ denotes the reflection associated to
the simple root $\alpha$ of $T_G$ in $B_G$. The variety of stars,
denoted $S$, is closed in ${\mathcal B}^{|\Omega |}$, projective,
but not in general regular as a variety.
It is relatively easy to see that a regular star is a star.
In fact, the variety of stars is the Zariski closure of $S_0$
in ${\mathcal B}^{|\Omega|}$. (One way to see this is by means of
a more explicit realization of $S_0$ as a subset of $S$ using
coordinate patches (Lemma 1.2 below). That $S_0$ is
open in its Zariski closure is a
well-known result in group theory, e.g.,
% Claude: [Sp3] is cited here but does not appear in the bibliography of the
% source; the citation is typeset literally.
[Sp3, ch.4, Lemma 4.3.1], \cite[II.6, Prop. 6.7]{Bor91}.)

We now define some open patches on $S$ and regular
coordinates on these patches.
Fix a Borel $B_+$ of $G$ and another Borel $B_-$ which is
opposite to $B_+$; let $T_+:=B_+\cap B_-$. Let
\[
S(B_+):=\{(B(W))_{W\in{\mathcal W}}\in S\ |\ \forall W\in{\mathcal W},\
B(W) \text{ is opposite to }B_+\}
\]
and let
\[
S(B_+,B_-):=\{(B(W))_{W\in{\mathcal W}}\in S(B_+)\
|\ B(W_0)=B_-\}.
\]
Let $N_+$ denote the unipotent radical of $B_+$. The map
\[
((B(W))_{W\in{\mathcal W}},n)\longmapsto (B(W)^n)_{W\in{\mathcal W}}
\]
defines an isomorphism
\[
S(B_+,B_-)\times N_+ \rightarrow S(B_+).
\tag{1.7}
\]
This is because for $B_-$ and $B_+$ as above, the map
$n\longmapsto B_-^n$, $n\in N_+$, defines an isomorphism from
$N_+$ to the open subvariety ${\mathcal B}(B_+)\subset {\mathcal B}$
of all Borels opposite to $B_+$.

Following \cite{La83}, we now describe local coordinates for
$S(B_+,B_-)$.

We identify the absolute Weyl group of $T_G$ with the absolute
Weyl group of
$T_+$.
We identify the set of Weyl chambers of $T_G$ with those of $T_+$.
If $(W,\rho)$ denotes a Weyl chamber and a
root defining a wall
of it (which we refer to below as a \textbf{wall-chamber pair}
or often simply as a \textbf{pair}),
we may find a $\omega\in\Omega$ for which
$W=\omega^{-1}W_0$ and $\rho=\omega \alpha$, for some
simple root $\alpha$ of $T_+$ in $B_-$;
we sometimes write $\alpha(W,\rho)$ for
the $\alpha$ obtained in this way.

Given a star $(B(W))_{W\in{\mathcal W}}\in S(B_+,B_-)$,
to each pair $(W,\rho)$ one defines
the value at $(B(W))_{W\in{\mathcal W}}$ of the coordinate
function $z(W,\rho)$ as follows. Let $ker(\rho)
\subset X^*(T_G)\otimes {\rrr}$
denote the wall separating two adjacent Weyl chambers
$W$, $W'$ and let $h\in N_+(\ovF)$ denote the unique element
satisfying
\[
B(W)=h^{-1}B_-h
\]
and
\[
B(W')\subset h^{-1}P_{-,\alpha} h
\]
where $P_{-,\alpha}$ is the smallest parabolic subgroup
of $G$ containing $B_-$ and $B_-^{s_\alpha}$, $\alpha
=\alpha (W,\rho)$.
Recall that
\[
\gga_a\cong U_{-\alpha}\cong P_{-,\alpha}/B_-,
\]
via $z\longmapsto u=exp(X_{-\alpha})\longmapsto B_-^u$, which we may
regard as an element of $P_{-,\alpha}/B_-\hookrightarrow G/B_-\cong {\mathcal B}$.
Let ${\mathbb A}$ denote the affine line.
There is a coordinate function $z(W,\rho)$,
\[
z(W,\rho):S(B_+,B_-)\rightarrow \aaa,
\]
such that
\[
hB(W')h^{-1}=\theta_{-\alpha} (z(W,\rho))B_-\theta_{-\alpha} (-z(W,\rho)),
\ \ \alpha=\alpha(W,\rho),
\tag{1.8}
\]
where $X_\alpha$, $X_{-\alpha}$ are root vectors satisfying
\[
[X_\alpha, X_{-\alpha}]=H_\alpha,\ \
\beta(H_{\alpha})=2\frac{(\beta,\alpha)}{(\alpha,\alpha)},\ \
\theta_{-\alpha}(z)=\exp(zX_{-\alpha})\in U_{-\alpha}.
\]
The local coordinate $z(W_0,\alpha)$ on $S(B_+,B_-)$ is
a local coordinate on a line of type $\alpha$ in the sense
of \cite{St74}. This observation will be useful when dealing with the
comparison of the subregular germ.

We have

(a) $z(W,\rho)=0$ if and only if $B(W)=B(W')$ and

(b) the collection of the $z(W,\rho)$, for all pairs $(W,\rho)$,
uniquely determines the star $(B(W))_{W\in{\mathcal W}}\in S(B_+,B_-)$,

(c) the coordinate functions $z(W,\rho)$ are regular and generate
the coordinate ring of $S(B_+,B_-)$.

The coordinates $z(W,\rho)$ are not algebraically independent.
The coordinate relations they satisfy can be most naturally
described using the notion of a gallery in the spherical building
associated to the Weyl group $\Omega$ of $T_G$. We may identify the
spherical building of $T_G$ with that of $T_+$.
Recall that a \textbf{gallery} is a sequence $W_0,W_1,\cdots, W_r$
of Weyl chambers such that $W_i$, $W_{i+1}$ are distinct but share a wall
$ker(\rho_{i+1})$. We always take $W_0$ to be the positive Weyl chamber of
$T_+$ with respect to $B_+$ and define
$\omega_i\in\Omega$ to be that element for which
$W_i=\omega_i^{-1}W_0$. Let $\alpha_{i+1}:=\omega_i\rho_{i+1}$,
so $\alpha_{i+1}=\alpha(W_i,\rho_{i+1})$. If $z_i:=z(W_{i-1},\rho_i)$
for $1\leq i\leq r$,
then
\[
B(W_{i-1})=h^{-1}B_-h,\ \ \
B(W_i)=h^{-1}\theta_{-\alpha_i}(-z_i)
B_-\theta_{-\alpha_i}(z_i)h,
\tag{1.9}
\]
for some $h\in N_+(\overline{F})$ possibly depending on $i$. We
define $n,n_{W_j}$, for $1\leq j\leq r$, by
$B(W_j)=B_-^{n_{W_j}n}$, where $n,n_{W_j}\in N_+(\overline{F})$
and
\[
n_{W_j}=\exp (z_jX_{-\alpha_j})
\exp (z_{j-1}X_{-\alpha_{j-1}})\cdots \exp (z_1X_{-\alpha_1}).
\tag{1.10}
\]
In particular, if $W_r=W_0$ (so that the gallery is closed),
then
\[
\exp (z_rX_{-\alpha_r})
\exp (z_{r-1}X_{-\alpha_{r-1}})\cdots \exp (z_1X_{-\alpha_1})=1.
\tag{1.11}
\]
In particular, if $W$, $W'$ are adjacent and share a wall
$ker(\rho)$ then $z(W,\rho)=-z(W',\rho)$.

We now state an observation of Langlands \cite[pp.483-484]{La83}

\begin{proclaim}{Lemma 1.12} TFAE:

\noindent
(a) A star $(B(W))_{W\in{\mathcal W}}\in S(B_+,B_-)$ is regular.

\noindent
(b) $B(W)\not= B(W')$ for all adjacent $W$, $W'$.

\noindent
(c) $z(W,\rho)\not= 0$ for all pairs $(W,\rho)$.
\end{proclaim}

We define coordinate patches on a resolution $S_1$ of $S$ (\cite{Ha92a}).

Recall that all the short (resp. long) roots lie in the same orbit
under the Weyl group. With this fact and the above lemma in mind, define
\[
S'(B_+,B_-):=\{(B(W))_{W\in{\mathcal W}}\in S(B_+,B_-)
\ |\ \forall \alpha\in\Delta ,\exists \omega\in\Omega
\text{ s.t. }z(\omega^{-1}W_0,\omega\alpha)\not= 0\}.
\]
Suppose $(B(W))_{W\in{\mathcal W}}\in S(B_+,B_-)$ is not in
$S'(B_+,B_-)$. Then $z(\omega^{-1}W_0,\omega\alpha)=0$
for all $\alpha\in\Delta$ and $\omega\in\Omega$. By the lemma,
this implies $B(W)=B(W')$ for all adjacent Weyl chambers. This is
the diagonal in ${\mathcal B}^{|\Omega|}$ and the diagonal contains no
regular stars.

Let $\aaa^\Delta$ denote the $|\Delta|$-fold Cartesian product
of affine spaces and let $T_+$ act on $\aaa^\Delta$ by
\[
t:(z(\alpha))_{\alpha\in\Delta}\longmapsto
(\alpha(t)^{-1}z(\alpha))_{\alpha\in\Delta},
\]
and on $S'(B_+,B_-)$ by
\[
t:(z(W,\rho))_{(W,\rho)\text{ a pair}}\longmapsto
(\alpha(t)z(W,\rho))_{(W,\rho)\text{ a pair}},\ \ \ \alpha=\alpha(W,\rho).
\]
Let
\[
S_1(B_+,B_-):=S'(B_+,B_-)\times_{T_+}\aaa^\Delta:=
(S'(B_+,B_-)\times\aaa^\Delta)/T_+.
\tag{1.13}
\]
Observe that there is a ``blow-down'' map
\[
S_1(B_+,B_-)\rightarrow S(B_+,B_-),
\tag{1.14}
\]
given by $(s',z)T_+\longmapsto s$, where
$s'=(z_1(W,\rho))_{(W,\rho)\text{ a pair}}$,
$z=(z(\alpha))_{\alpha\in\Delta}$ (note $z(\alpha)$ is only defined up to
multiplication by $\alpha(t)^{-1}$), and
\[
s=(z(W,\rho))_{(W,\rho)\text{ a pair}}
=(z(\alpha(W,\rho))z_1(W,\rho))_{(W,\rho)\text{ a pair}}.
\]
This is clearly well-defined (i.e., independent of the chosen
representative
$z=(z(\alpha))_{\alpha\in\Delta}$ of the $T_+$-orbit in $\aaa^\Delta$).
The subscript 1 on $z_1(W,\rho)$ has been added to distinguish it
from its image
$z(W,\rho):=z(\alpha(W,\rho))z_1(W,\rho)$ in $S(B_+,B_-)$.
Later we shall fix a canonical choice of a generic representative
$(z(\alpha))_{\alpha\in\Delta}$ of the $T_+$-orbit in $\aaa^\Delta$:
$z(\alpha)$ is chosen to be that multiple of $\alpha(t)^{-1}$
such that $z(\alpha)=\lambda /x(\alpha)$, where $x(\alpha)$ is as in
(1.17) below (equivalently: $z(\alpha)$ is chosen to be that multiple
which makes $w(\alpha)=1$ for $\alpha >0$ simple \cite[ch II, \S 2]{Ha92a}).

Let $S'(B_+)\subset S(B_+)$ be the image of
$S'(B_+,B_-)\times N_+$ in $S(B_+)$ under the map
\[
((B(W))_{W\in{\mathcal W}},n)\longmapsto (B(W)^n)_{W\in{\mathcal W}},
\]
and let $S_1(B_+)\subset S'(B_+)\times_{T_+}\aaa^\Delta$
be the image of $S_1(B_+,B_-)\times N_+$ under
\[
((B(W))_{W\in{\mathcal W}},(z(\alpha))_{\alpha\in\Delta},n)
\longmapsto ((B(W)^n)_{W\in{\mathcal W}},(z(\alpha))_{\alpha\in\Delta}).
\]
Therefore,
\[
S'(B_+)\cong S'(B_+,B_-)\times N_+,
\tag{1.15a}
\]
and
\[
S_1(B_+)\cong S_1(B_+,B_-)\times N_+.
\tag{1.15b}
\]
The apparent dependence of $S_1(B_+)$ on $B_-$ is not serious:
by \cite[Lemma 3.6]{La83}, if we replace $B_-$ by another
Borel $B'_-$ opposite
to $B_+$ then the corresponding variety is canonically isomorphic
to $S_1(B_+)$.
Let $B'_+$ be another Borel and let $B_-$ be a Borel opposite to both
$B_+$ and $B'_+$. There is a birational morphism from
$S_1(B_+)$ to $S_1(B'_+)$ as defined in \cite{La83}. Using this birational map
to paste $S_1(B_+)$ and $S_1(B'_+)$ together, we obtain
a variety $S_1$.

Observe that the ``obvious'' map $S_0(B_+)\rightarrow S_1(B_+)$
extends to a morphism $S_0\rightarrow S_1$. This map yields a map
$X_0\rightarrow G\times S_1$.

\begin{proclaim}{Definition 1.16}
Let $X$ be the Zariski closure of $X_0$ in $G\times S$.
Let $X_1$ be the closure of $X_0$ in $G\times S_1$ and let $Y_1$
be the closure of $Y_0$ in $X_1$. Let $X_1(B_+)$ denote the
intersection of $X_1$ with $G\times S_1(B_+)$ and similarly for
$Y_1(B_+)$. Similarly, we define $X''$, $Y''$, $X''(B_+)$,
$Y''(B_+)$, replacing $G\times S_1(B_+)$
by $G\times S''$ or $G\times S''(B_+)$ in the definition of
$X_1$ and $Y_1$, where $S''$ is the image of $S_0$ in $Y_1$.
\end{proclaim}

In particular, the first coordinate of a point
$p\in Y''(B_+,B_-)\subset G\times S''$ is an element
$b$ of $\cap_W B_-^{n_W} \subset B_-$, which we may write as
an ordered product
\[
b=t_0\prod_{\rho\in\Sigma_G^+}\exp (x(\rho)X_\rho),
\tag{1.17}
\]
where $\Sigma_G^+$ denotes the set of roots of $T_+$
which are positive with respect to the Borel $B_-$, on which
we fix an ordering once and for all.

\subsection*{Coordinates on $Y''$}

Consider a point $x\in X_1\subset G\times S_1$.
On the patch of $X_1$ meeting
$G\times S_1(B_+,B_-)\times N_+\cong
G\times S_1(B_+)$, we write $x$ as
$(b,\allowbreak(B(W))_{W\in{\mathcal W}},\allowbreak(z(\alpha))_{\alpha\in\Delta},\allowbreak n)$,
where $b\in\cap_{W\in{\mathcal W}}B(W)\subset B_-$ and
$B(W_0)=B_-$ (the condition that $b\in\cap_{W\in{\mathcal W}}B(W)$
is inherited from the definition of $X_0$).

Other coordinates on $X''(B_+)$ may be defined in terms of a point
$(t^g,(B(W)^g)_{W\in{\mathcal W}})$ in
$G\times S(B_+,B_-)$, where $t\in T_G(\ovF )$ is regular. We may write
\[
g=g_0t_+n_-n_+,
\]
where $T_G^{g_0}=T_+$, $t_+\in T_+(\ovF)$, $n_-\in N_-(\ovF)$,
and $n_+\in N_+(\ovF)$. Let $t^{g_0}=t'_+$, for some
$t'_+\in T_+(\ovF)$, and write $t^g=(t'_+n'_-)^{n_+}$,
for some $n'_-\in N_-(\ovF)$. Therefore,
\[
(t^g,(B(W)^g)_{W\in{\mathcal W}})=(t'_+n'_-,B_-^{n_W})^{n_+},
\]
where $n_W\in N_+(\ovF)$. We write
\[
n'_-=\prod_{\rho\in\Sigma_G^+}\exp (x(\rho)X_\rho),
\]
\[
n_+=\prod_{\rho\in\Sigma_G^+}\exp (v(\rho)X_{-\rho}),
\]
where, by (1.10),
\[
n_W=\exp (z(W_r,\rho_r)X_{-\alpha_r})
\exp (z(W_{r-1},\rho_{r-1})X_{-\alpha_{r-1}})\cdots
\exp (z(W_1,\rho_1)X_{-\alpha_1}),
\tag{1.18}
\]
where $W=\omega^{-1}W_0$, $\omega=s_{\alpha_r}\cdots s_{\alpha_1}$,
and $W_0$, ..., $W_r$ being chambers in a path from the positive Weyl chamber
$W_0$
to $W_r=W$. Recall we defined $z_1(W,\rho)$, for any pair $(W,\rho)$, by
\[
z(W,\rho)=z(\alpha(W,\rho))z_1(W,\rho).
\tag{1.19}
\]

Define the coordinate function $w(\rho)$, $\rho >0$, by
$w(\alpha)=1$, if $\alpha >0$ is simple, and otherwise by
\[
(t'_+)^{-1}(t'_+)^{n_-}=
\prod_{\rho\in\Sigma_G^+}
\exp (\frac{w(\rho)\lambda}{\prod_{\alpha\in\Delta}z(\alpha)^{m(\alpha)}}X_\rho),
\tag{1.20}
\]
where $\rho=\sum_{\alpha\in\Delta}m(\alpha)\alpha$. This describes a point
in the image
\[
C'\times G/T_G\hookrightarrow Y_1(B_+,B_-)
\subset X_1(B_+,B_-)\subset G\times S_1(B_+,B_-),
\]
where $C'=C\cap T_G'=C-\{1\}$.

\subsection*{Coordinates on $Y_1$}

We begin with the following

\begin{proclaim}{Lemma 1.21}(\cite[ch II, Lemma 3.3]{Ha92a}) If $t\in T_+$ and
\[
n=\prod_{\rho >0}\theta_\rho(n_\rho)\in N_+
\ \ \ \ \ (\text{ordered product}),
\]
then
\[
t^n=t\prod_{\rho >0}\theta_\rho(y(\rho)),
\]
where
\[
y(\rho)=\sum_{\rho_1,...,\rho_m} c_{\rho_1,...,\rho_m}(t)
n_{\rho_1}...n_{\rho_m},
\]
the sum ranging over all roots $\rho_i>0$ such that
$\rho=\rho_1+...+\rho_m$ and
\[
c_{\rho_1,...,\rho_m}(t)=
\sum_{\sigma_1,\sigma_2}(1-\sigma_1(t)^{-1})\sigma_2(t),
\]
where the sum over $\sigma_1,\sigma_2$ such that $\sigma_1$ is a root
and $\sigma_2$ is a sum of roots. Moreover, $c_\rho(t)=(1-\rho(t)^{-1})$.
\end{proclaim}

We want to apply this lemma to obtain coordinates of a point
\[
(t^g,(B(W))_{W\in{\mathcal W}},(z(\alpha))_{\alpha\in\Delta})^n\in
\chi^{-1}(C)\times S_1(B_+,B_-)\times N_+.
\]
For the case $G=G_2$, we
give coordinate relations between the coordinates in (1.20) and
the $n_\rho$'s and the $z(W,\alpha)$'s (see \cite[chap II]{Ha92a}):
\[
x(\alpha)=(1-\alpha^{-1})n_\alpha = (1-\alpha^{-1})/z(W_0,\alpha),
\tag{1.22}
\]
\[
x(\beta)=(1-\beta^{-1})n_\beta = (1-\beta^{-1})/z(W_0,\beta),
\]
\[
x(\gamma)=(1-\gamma^{-1})n_\gamma +
c_1\cdot \frac{(1-\alpha^{-1})}{{z(W_0,\alpha)z(W_0,\beta)z(\sigma_\beta^{-1}W_0,\alpha)}}
+\text{ higher order terms in }\lambda,
\]
\[
x(\delta)=(1-\delta^{-1})n_\delta +
c_2\cdot \frac{(1-\alpha^{-1})}{{z(W_0,\alpha)z(W_0,\beta)z(\sigma_\beta^{-1}W_0,\alpha)}}
+\text{ higher order terms in }\lambda,
\]
\[
x(\epsilon)=(1-\epsilon^{-1})n_\epsilon +
c_3\cdot \frac{(1-\alpha^{-1})n_\delta}{{z(W_0,\alpha)}}
+\text{ higher order terms in }\lambda,
\]
\[
\begin{split}
x(\zeta)
&=(1-\zeta^{-1})n_\zeta +
c_4\cdot \frac{(1-\alpha^{-1})n_\gamma^2}{{z(W_0,\alpha)}}+
c_5\cdot \frac{(1-\alpha^{-1})n_\delta}{{z(W_0,\alpha)z(W_0,\beta)}}+\\
&+c_6\cdot \frac{(1-\beta^{-1})n_\epsilon}{{z(W_0,\beta)}}+
c_7\cdot \frac{(1-\gamma^{-1})n_\delta}{{z(W_0,\beta)z(\sigma_\beta^{-1}W_0,\alpha)}}
+\text{ higher order terms in }\lambda,
\end{split}
\]
\[
\begin{split}
w(\gamma)
&=\frac{x(\gamma)z(\beta)z(\alpha)}{\lambda}\\
&=(\gamma(X)n_\gamma + c_1\frac{\alpha(X)}{{z(W_0,\alpha)z(W_0,\beta)}})
\alpha(X)^{-1}\beta(X)^{-1}z(W_0,\alpha)z(W_0,\beta)+\\
&\quad\quad +\text{ higher order terms in }\lambda,
\end{split}
\]
\[
\begin{split}
w(\delta)
&=\frac{x(\delta)z(\beta)z(\alpha)^2}{\lambda}\\
&=(\delta(X)n_\delta +
c_2\frac{\alpha(X)}{{z(W_0,\alpha)z(W_0,\beta)z(\sigma_\beta^{-1}W_0,\alpha)}})
\alpha(X)^{-2}\beta(X)^{-1}z(W_0,\alpha)^2z(W_0,\beta)+\\
&\quad\quad +\text{ higher order terms in }\lambda,
\end{split}
\]
\[
\begin{split}
w(\epsilon)
&=\frac{x(\epsilon)z(\beta)z(\alpha)^3}{\lambda}\\
&=(\epsilon(X)n_\epsilon + c_3\frac{\alpha(X)n_\delta}{{z(W_0,\alpha)}})
\alpha(X)^{-3}\beta(X)^{-1}z(W_0,\alpha)^3z(W_0,\beta)+\\
&\quad\quad +\text{ higher order terms in }\lambda,
\end{split}
\]
\[
\begin{split}
w(\zeta)
&=\frac{x(\zeta)z(\beta)^2z(\alpha)^3}{\lambda}\\
&=(\zeta(X)n_\zeta + c_4\frac{\alpha(X)n_\gamma^2}{{z(W_0,\alpha)}}
+c_5\frac{\alpha(X)n_\delta}{{z(W_0,\alpha)z(W_0,\beta)}}
+c_6\frac{\beta(X)n_\epsilon}{{z(W_0,\beta)}}+\\
&\quad +\frac{\gamma(X)n_\zeta}{{z(\sigma_\beta^{-1}W_0,\alpha)z(W_0,\beta)}})
\alpha(X)^{-3}\beta(X)^{-2}z(W_0,\alpha)^3z(W_0,\beta)^2+\\
&\quad\quad +\text{ higher order terms in }\lambda,
\end{split}
\]
Because of (1.22), on $Y_1(B_+,B_-)$, we may
use the algebraically independent coordinates
\[
z(W_0,\alpha),\ z(W_0,\beta),\ z(\sigma_\beta^{-1}W_0,\alpha),\
n_\gamma,\ n_\delta,\ n_\epsilon,\ n_\zeta,
\]
in place of the coordinates
\[
z(\alpha),\ z(\beta),\ x(\alpha),\ x(\beta),\ x(\gamma),\
x(\delta),\ x(\epsilon),\ x(\zeta),\
w(\gamma),\ w(\delta),\ w(\epsilon),\ w(\zeta),
\]
on $Y''(B_+,B_-)$ appearing in (1.17) and (1.20).
% ======================================================================
% Claude: Section 2, source file g2new2.tex (\pageno 19). Here \off = \overline{F}.
% ======================================================================
\section*{\S\,2 Galois action on the coordinates:
explicit computations for $G_2$}

Let $G$ denote a connected semisimple $F$-quasi-split
group with no outer automorphisms.

If $T$ is a Cartan of $G$ defined over $F$ then
$\Gamma=Gal(\overline{F}/F)$ preserves
$T$ and the set $\Sigma$ of roots of $G$ with respect to $T$,
$\sigma\in \Gamma$ sending $U_\rho$ to $U_{\sigma\rho}$, where
$U_\rho$ is a root group as in \S1. Generally
speaking, $\Gamma$ will not preserve the set $\Delta$ of positive
simple roots with respect to a Borel $B$ containing $T$, unless
$B$ is defined over $F$. The character lattice
$X:=X^*(T)=X^*(T/\overline{F})$ is a $\zzz[\Gamma]$-module and the action
of $\Gamma$ on $X_{\rrr}:=X^*(T)\otimes \rrr$ preserves the
set of root hyperplanes. Consequently, $\Gamma$ acts on the
set of Weyl chambers in $X_{\rrr}$ and preserves the positive
Weyl chamber if $B$ is defined over $F$.
The action of $\sigma\in\Gamma$ on the set
${\mathcal W}$ of Weyl chambers in $X^*(T_G)\otimes \rrr$
will be denoted $\sigma_T$.

As in section 1, we fix a Borel $\underline{B}$ of $G$
defined over ${F}$, a Cartan
$\underline{T}$ defined over $F$ contained in $\underline{B}$,
and a positive Weyl chamber $W_0$ with respect to the linear order
on $\Sigma_G^+$ which makes all the roots occuring in $\underline{B}$
positive. Let $T_G$ denote our fixed Cartan of $G$ defined over $F$
and let $B_G$ denote a Borel containing
$T_G$. Note that, for any $\sigma\in\Gamma$ both $B_G$ and
$B_G^\sigma$ are Borels containing $T_G$, since $T_G$ is
defined over $F$.
The Galois action on $\underline{T}$ induces a homomorphism
\[
\rho:\Gamma\rightarrow Aut(X,\Delta,X^{\vee},\Delta^{\vee}),
\tag{2.1a}
\]
where $X=X^*(\underline{T})$ denotes the root lattice,
$X^{\vee}=X_*(\underline{T})$ denotes the cocharacter lattice,
$\Delta=\Delta(\underline{T},\underline{B})$ the simple roots of
$\underline{T}$ with respect to $\underline{B}$,
and $\Delta^{\vee}=\Delta^\vee(\underline{T},\underline{B})$
denotes the simple coroots.
This induces (\cite{Sp79}) a homomorphism
\[
\tilde{\rho}:\Gamma\rightarrow
Aut(G,B,T,\{X_{\alpha}\}_{\alpha\in\Delta}\}),
\tag{2.1b}
\]
where $X_\alpha$ denotes a root vector associated to $\alpha$
in the Lie algebra of $G$.

The Borels
$\{{B}_G(W))\ |\ W\in {\mathcal W}\}$ are precisely the Borel
subgroups of $G$ which contain ${T}_G$ and, hence, are preserved by the
Galois action since ${T}_G$ is defined over $F$. Define
a map $\sigma_T:{\mathcal W}\rightarrow {\mathcal W}$ by
\[
\sigma {B}_G(W)={B}_G(\sigma_T W),
\]
for $\sigma\in\Gamma$. This defines a continuous homomorphism
$\Gamma\rightarrow \Omega$,
$\sigma\longmapsto \sigma_T$, which factors through $Gal(\overline{F}/E)$
where $E/F$ denotes the splitting field of $T_G$.
We may write $B_G=\underline{B}^g$,
for some $g\in G(\overline{F})$ such that
$g^{-1}\cdot \sigma g\in T_G(\overline{F})$ for all $\sigma\in \Gamma$.
This action $\sigma_T$ on the chambers $W\in {\mathcal W}$ induces an action,
also denoted $\sigma_T$, on the positive roots $\Sigma_G^+$ and hence on
the walls $ker(\rho)$ of $W$ to the walls
$ker(\sigma_T(\rho))$ of
$\sigma_TW$.

We define an action of
$\Gamma$ on the set of regular stars, denoted $S_0$, by
\[
\sigma (B(W)):={B}_G(W)^{\sigma h}
=\sigma{B}(\sigma_T^{-1}W),
\tag{2.2}
\]
where $B(W)=B_G(W)^h$.
The set of $F$-rational points of $S_0(\overline{F})$
under this $\Gamma$-action is the set
of regular stars $(B(W))_{W\in {\mathcal W}}$ such that
$B(W)=B_G(W)^h$, with $\sigma(h)h^{-1}\in T_G(\overline{F})$
for all $\sigma\in\Gamma$.

For $\gamma\in T_G(F)$, let
\[
\phi_\gamma^{-1}(F):=\phi^{-1}(\gamma)^\Gamma=
\{x\in X^0(F)\ |\ \phi(x)=\gamma\},
\]
where $\phi$ is defined following (1.2) above.
In fact, the projection map $\pi$ of \S 1.1
yields a set-theoretic bijection
\[
\phi_\gamma^{-1}(F)
\rightarrow \{g^{-1}\gamma g\ |\ g\in (T_G\setminus G)(F)\}
\cong (T_G\setminus G)(F),
\tag{2.3}
\]
onto the stable conjugacy class of $\gamma$.

The action of $\sigma\in\Gamma$ on a point
$((B_-^{n_Wn})_{W\in{\mathcal W}})\in
S(B_+)$, for some
$n,n_W\in N_+$, which we identify with the point
$((B(W))_{W\in{\mathcal W}},n)\in S(B_+,B_-)\times N_+$,
is given by
\[
\sigma: S(B_+,B_-)\times N_+ \rightarrow S(B_+,B_-)\times N_+,
\]
\[
((B(W))_{W\in{\mathcal W}},n)\longmapsto
((B'(W))_{W\in{\mathcal W}},n'),
\tag{2.4a}
\]
where $\sigma B_-=B_-^{c_\sigma}$,
$c_\sigma\in N_{Aut\ G}(T_+)(\ovF)=N_G(T_+)(\ovF)$, and
\[
B'(W)^{n'}=\sigma(B(\sigma_T^{-1}(W))^n)=
(\sigma B_-)^{\sigma(n_{\sigma_T^{-1}(W)}n)},
\tag{2.4b}
\]
and
\[
n'=c_\sigma\sigma(n_{\sigma_T^{-1}(W_0)}n),\ \ B'(W_0)= B_- \ .
\tag{2.4c}
\]

How does the action of $\Gamma$ on $Y_1(B_+)$ affect the
point determined by $(b,(B_-^{n_W})_{W\in{\mathcal W}})^{n}$?
Here $b\in B_-$, $n_W,n\in N_+(\ovF)$.
Let $\alpha(W,\rho)\in\Delta$ denote the simple root associated to
a pair $(W,\rho)$ as above.
To be consistent with (2.4a), define an action of $\sigma\in\Gamma$
on $\Delta$, denoted by $\sigma_T^*$, as follows:
\[
(\sigma_T^*)^{-1}(\alpha(W,\rho))
=\alpha(\sigma_T^{-1}(W),\sigma_T^{-1}(\rho)).
\tag{2.5}
\]
The action of $\Gamma$ on
$S_1(B_+)\cong (S'(B_+,B_-)\times_{T_+}\aaa^\Delta)\times N_+$,
is given by
\[
\sigma:((B(W))_{W\in{\mathcal W}},(z(\alpha))_{\alpha\in \Delta},n)\longmapsto
((B'(W))_{W\in{\mathcal W}},(z'(\alpha))_{\alpha\in \Delta},n'),
\]
where $B'(W)$ and $n'$ are as above, and
\[
z'(\alpha)=\sigma(z((\sigma_T^*)^{-1}(\alpha))).
\]
By the formulas
for the Galois action on $S(B_+)$, we have
\[
\sigma:(b,(B_-^{n_W})_{W\in{\mathcal W}})^{n}\longmapsto
(\sigma b,(\sigma B_-^{\sigma (n_{\sigma^{-1}_T (W)})}
)_{W\in{\mathcal W}})^{\sigma(n_{\sigma^{-1}_T (W_0)}n)}.
\tag{2.6}
\]
Here $b$ is as in (1.17) and
\[
n_W=\exp (z(W_r,\rho_r)X_{-\alpha_r})
\exp (z(W_{r-1},\rho_{r-1})X_{-\alpha_{r-1}})\cdots
\exp (z(W_1,\rho_1)X_{-\alpha_1}),
\]
where $W_0,...,W_r=W$ is a gallery as in (1.9-10).
How does the action of $\Gamma$ on $S(B_+)$ affect the coordinates
$z(W,\rho)$? A point of $S(B_+,B_-)$ is determined by
$(z(W,\rho))_{(W,\rho)\ a \ pair}$, so by (2.2) and (2.6),
\[
\sigma:(z(W,\rho))_{(W,\rho)\ a\ pair}\longmapsto
(z'(W,\rho))_{(W,\rho)\ a\ pair},
\]
where
\[
z'(W,\rho)=\sigma(z(\sigma_T^{-1}(W),\sigma_T^{-1}(\rho)).
\tag{2.7}
\]
Suppose now that we are on a patch $S(B_+,B_-)$ with
$T_+=B_+\cap B_-$ defined over $F$.
There is an $c_\sigma\in N_{Aut\ G}(T_+)(\ovF)=
N_G(T_+)(\ovF)$ such that
$\sigma B_-=c_\sigma B_-c_\sigma^{-1}$. We may then write
\[
\sigma:(b,(B_-^{n_W})_{W\in{\mathcal W}})^{n}\longmapsto
(b',(B_-^{n'_W})_{W\in{\mathcal W}})^{n'},
\tag{2.8}
\]
where
\[
b'= \sigma({n_{\sigma^{-1}_T (W_0)}})c_\sigma b
c_\sigma\sigma({n^{-1}_{\sigma^{-1}_T (W_0)}})^{-1},
\]
\[
n'_W=c_\sigma
\sigma(n_{\sigma^{-1}_T (W)}n^{-1}_{\sigma^{-1}_T (W_0)})c_\sigma^{-1},
\]
\[
n'= c_\sigma \sigma(n_{\sigma^{-1}_T (W_0)}n).
\]
The action of $c_\sigma$ by conjugation commutes with the action of
$\sigma_T$ by permutation of the Borels in a star.
If $B_-$ is defined over $F$ then the
$W_0$-component of $(B_-^{n'_W})_{W\in{\mathcal W}}$ is
$\sigma B_-=B_-$. In particular, if the point in $Y_1(B_+)$
determined by $(b,(B_-^{n_W})_{W\in{\mathcal W}},
(z(\alpha))_{\alpha\in\Delta},{n_+})$ has the
coordinates $\{x(\rho)\ |\ \rho >0\}$ then the image of this point
under $\sigma$ has the coordinates $\{x'(\rho)\ |\ \rho >0\}$,
where $x'(\rho)$ is determined by
\[
\sigma({n_{\sigma^{-1}_T (W_0)}})c_\sigma b
c_\sigma\sigma({n^{-1}_{\sigma^{-1}_T (W_0)}})^{-1}
=b'=n_0t_0\prod_{\rho\in\Sigma^+}\exp (x'(\rho)X_\rho),
\tag{2.9}
\]
for some $n_0\in N_+=<U_\rho\ |\ \rho>0>$,
where $\Sigma^+$ denotes the set of roots of $T_+$
which are positive with respect to $B_-$.
Using Chevalley's commutator formula, this formula for
$x'(\rho)$ will be used for explicit calculations below.

\begin{proclaim}{Lemma 2.10 \cite[pf of Lemma 5.3]{Ha89}}
(a) Let $\alpha$ denote an
arbitrary simple positive root of $G$, an arbitrary quasi-split group. If
$(u,(B(W))_{W\in {\mathcal W}})\in E_{\Delta-\alpha}$
then every such Borel subgroup $B(W)$ lies in the same line $\ell_\beta$
in the Dynkin curve for $G$. Here
$E_{\Delta-\alpha}$ is a $1$-regular Richardson branch
of $\phi^{-1}(1)$ on $Y_C$ as in \S 3 below.

(b) The action of Galois on the points of $E_{\Delta-\alpha}$
is determined by the action on the lines in the Dynkin curve.
\end{proclaim}

Now we restrict to the case $G=G_2$.

Recall the Dynkin curve for $G_2$ \cite{St74}:

% Claude: The source has five empty display pairs ($$ $$) here, followed by
% \vskip .4in. The figure of the Dynkin curve is missing from the source.
% Only the vertical space is reproduced.
\vspace{0.4in}

\noindent
Possible Galois actions on this diagram include:

\noindent
(i) exchanging $\ell_\beta'$ with $\ell_\beta'''$, leaving
$\ell_\beta''$ fixed, and flipping $\ell_\alpha$;

\noindent
(ii) flipping $\ell_\beta'$, $\ell_\beta'''$, $\ell_\beta''$
and fixing $\ell_\alpha$;

\noindent
(iii) cyclically permuting the lines $\ell_\beta'$,
$\ell_\beta'''$, $\ell_\beta''$
and ``rotating'' $\ell_\alpha$.

Label and order the positive roots of the maximal torus
$T_+\subset B_-$ of $G_2$ as
\[
\alpha \text{  (short) },\beta \text{  (long) },
\gamma=\alpha +\beta, \delta =2\alpha +\beta,
\epsilon =3\alpha +\beta,  \zeta=3\alpha +2\beta .
\]
We have
\[
\theta_\delta(r)\theta_\gamma(s)=\theta_\gamma(s)\theta_\delta(r)\theta_{\zeta}(3rs),
\]
so
\[
\begin{split}
\theta_{\eps}(x)^{\theta_{-\alpha}(z)}
&=\theta_{\eps}(x)\theta_\delta(c_1xz)\theta_\gamma(c_2xz^2)
\theta_\beta(c_3xz^3)\\
&=\theta_{\eps}(x)\theta_\beta(c_3xz^3)\theta_\delta(c_1xz)\theta_\gamma(c_2xz^2)\\
&=\theta_\beta(c_3xz^3)\theta_\gamma(c_2xz^2)\theta_\delta(c_1xz)\theta_{\eps}(x)
\theta_\zeta(c_4x^2z^3),
\end{split}
\]
where the $c_i$'s are computable constants. Likewise,
\[
\begin{split}
\theta_{\delta}(x)^{\theta_{-\alpha}(z)}
&=\theta_{\delta}(x)\theta_\gamma(d_1xz)\theta_\beta(d_2xz^2)\\
&=\theta_\beta(d_2xz^2)\theta_\gamma(d_1xz)\theta_{\delta}(x)\theta_\zeta(d_3x^2z),
\end{split}
\]
\[
\theta_{\gamma}(x)^{\theta_{-\alpha}(z)}
=\theta_{\gamma}(x)\theta_\beta(exz)=\theta_\beta(exz)\theta_{\gamma}(x),
\]
\[
\theta_{\beta}(x)^{\theta_{-\alpha}(z)}=\theta_{\beta}(x),
\]
and $\theta_{\alpha}(x)^{\theta_{-\alpha}(z)}=
\theta_{-\alpha}(y_\alpha)t\theta_{\alpha}(x)$, for some $y_\alpha$ depending on $x$,
$z$, and some $t$ belonging to a
maximal torus.

In general, if $\sigma\in\Gamma$ acts on
$T_G(\overline{F})$ pointwise as as simple
reflection
$s_\rho$, some $\rho\in\Delta$, then we write $\sigma=\sigma_\rho$.
This amounts to taking $\sigma(n_{\sigma_T^{-1}(W_0)})=\theta_\rho (z)$,
where $z=\sigma(z(\rho)z_1(W_0,\rho))$ \cite[\S 3]{Ha89}. We will
compute the Galois action on the coordinate of a point $P$ in
$Y''(B_+,B_-)$, where we assume that $B_+\cap B_-=T_+$ is defined over
$F$ and $c_\sigma =1$ in (2.9). (We will allow $c_\sigma$ to
be a reflection below, see (2.16-17).)

Taking $\sigma=\sigma_\alpha$, we have
\[
\prod_{\rho >0}\theta_\rho(x(\rho))^{\theta_{-\alpha}(z)}
=\theta_{-\alpha}(y_\alpha)
t\prod_{\rho >0}\theta_\rho(\phi_\rho^\alpha(\unx,z)),
\tag{2.11}
\]
where $\unx =(x(\alpha),\cdots, x(\zeta))$,
\[
\begin{split}
\phi_\alpha^\alpha(\unx,z)
&=x(\alpha),\\
\phi_\beta^\alpha(\unx,z)
&=x(\beta)+ezx(\gamma)+d_2z^2x(\delta)+c_3z^3x(\eps),\\
\phi_\gamma^\alpha(\unx,z)
&=x(\gamma)+d_1zx(\delta)+c_2z^2x(\eps),\\
\phi_\delta^\alpha(\unx,z)
&=x(\delta)+c_2zx(\eps),\\
\phi_{\eps}^\alpha(\unx,z)
&=x(\eps),\\
\phi_\zeta^\alpha(\unx,z)
&=x(\zeta)+d_3zx(\delta)^2+c_4z^3x(\eps)^2+3c_2z^2x(\delta)x(\eps).
\end{split}
\]
Therefore, if $c_{\sigma_\alpha}=1$, we have
\[
\sigma_\alpha :x(\rho)\longmapsto \phi_\rho^\alpha(\unx,z),
\tag{2.12}
\]
where $z=\sigma_\alpha(z(\alpha)z_1(W_0,\alpha))$.

Take $\sigma =\sigma_\beta$ and note
\[
\theta_\gamma(r)\theta_\alpha (s)=\theta_\alpha (s)\theta_\gamma(r)
\theta_\delta(2rs)\theta_\epsilon(3rs^2)\theta_\zeta(3r^2s),
\]
\[
\theta_\beta(r)\theta_\alpha (s)=\theta_\alpha (s)\theta_\beta(r)
\theta_\gamma(rs)
\theta_\delta(r^2s)\theta_\epsilon(r^3s)\theta_\zeta(r^3s^2).
\]
We have, by Chevalley's commutator formula,
\[
\theta_\alpha(x)^{\theta_{-\beta}(z)}=\theta_\alpha(x),
\]
\[
\theta_\beta(x)^{\theta_{-\beta}(z)}=
\theta_{-\beta}(y_\beta)
t\theta_\beta(x),\ \ \ (t\in T_G ,\
\alpha(t)=1+xz,\  y_\beta:=\frac{-x^2z}{1+xz}),
\]
\[
\begin{split}
\theta_\gamma(x)^{\theta_{-\beta}(z)}
&=\theta_\gamma(x)\theta_\alpha(czx)\\
&=\theta_\alpha(czx)\theta_\gamma(x)\theta_\delta(2czx^2)
\theta_\epsilon(3c^2z^2x^3)\theta_\zeta(3czx^3),
\end{split}
\]
\[
\theta_\delta(x)^{\theta_{-\beta}(z)}=\theta_\delta(x),
\]
\[
\theta_\epsilon(x)^{\theta_{-\beta}(z)}=\theta_\epsilon(x),
\]
\[
\theta_\zeta(x)^{\theta_{-\beta}(z)}=\theta_\zeta(x)\theta_\epsilon(dzx)
=\theta_\epsilon(dzx)\theta_\zeta(x).
\]
Therefore,
\[
\prod_{\rho >0}\theta_\rho(x(\rho))^{\theta_{-\beta}(z)}
=t\prod_{\rho >0}\theta_\rho(\phi_\rho^\beta(\unx,z)),
\tag{2.13}
\]
where
\[
\begin{split}
\phi_\alpha^\beta(\unx,z)
&=x(\alpha)(1+x(\beta)z)+cx(\gamma)z,\\
\phi_\beta^\beta(\unx,z)
&=x(\beta),\\
\phi_\gamma^\beta(\unx,z)
&=x(\gamma)+czx(\beta)x(\gamma),\\
\phi_\delta^\beta(\unx,z)
&=x(\delta)+2czx(\gamma)^2+c^2x(\beta)x(\gamma)^2z^2,\\
\phi_{\eps}^\beta(\unx,z)
&=x(\eps)+dzx(\zeta)+3c^2x(\gamma)^3z^2+c^3x(\beta)x(\gamma)^3z^3,\\
\phi_\zeta^\beta(\unx,z)
&=x(\zeta)+3zx(\gamma)x(\delta)+9czx(\gamma)^3+c^3z^3x(\beta)^2x(\gamma)^3.
\end{split}
\]
Therefore, if $c_{\sigma_\beta}=1$, we have
\[
\sigma_\beta :x(\rho)\longmapsto \phi_\rho^\beta(\unx,z),
\tag{2.14}
\]
where $z=\sigma_\beta(z(\beta)z_1(W_0,\beta))$.

Formulas (2.12) and (2.14) help to determine the set of $F$-rational
points of $Y_C$. We also need to compute the effect of conjugating $b$ by
$c_\sigma$. For this we recall a basic lemma \cite{Ca89}:

\begin{proclaim}{Lemma 2.15}
Let $\phi_\rho:SL(2)\rightarrow G$ be the
homomorphism associated to the root
$\rho$ and let $\theta_\rho:\gga_a\rightarrow U_\rho$ denote the root
group embedding, so
$\theta_\rho(x)=\phi_\rho(\begin{pmatrix} 1&x\\ 0&1\end{pmatrix})$. Let
$s_\rho:=s_\rho(1)$ where
\[
s_\rho(x)=\theta_\rho(x)\theta_{-\rho}(-x^{-1})\theta_\rho(x)
=\phi_\rho(\begin{pmatrix} 0&x\\ -1/x&0\end{pmatrix}).
\]
Then
\[
s_\rho \theta_\sigma (x)s_\rho^{-1}
=\theta_{w_\rho(\sigma)}(\eta_{\rho,\sigma}x),
\]
where $\eta_{\rho,\sigma}$ is as in \cite[pp 94-96]{Ca89} and $w_\rho$ is
the reflection of the Weyl group associated to $\rho$.
\end{proclaim}

As a consequence of this, we may write
\[
\begin{split}
&s_\alpha\cdot \prod_{\rho>0}\theta_{\rho}(x(\rho))\cdot s_\alpha^{-1}\\
&=\theta_{-\alpha}(x'(\alpha))
\prod_{\rho>0, \rho\not=\alpha}\theta_{\rho}(x'(\rho)),
\end{split}
\tag{2.16}
\]
where $x'(\alpha)=\eta_{\alpha,\alpha}x(\alpha)$,
$x'(\beta)=\eta_{\alpha,\epsilon}x(\epsilon)$,
$x'(\gamma)=\eta_{\alpha,\delta}x(\delta)$,
$x'(\delta)=\eta_{\alpha,\gamma}x(\gamma)$,
$x'(\epsilon)=\eta_{\alpha,\beta}x(\beta)$,
and
\[
x'(\zeta)=\eta_{\alpha,\zeta}x(\zeta)+
f_1x(\delta)x(\gamma)+f_2x(\beta)x(\epsilon),
\]
where $f_1,f_2$ are constants. Likewise, we may write
\[
\begin{split}
&s_\beta\cdot \prod_{\rho>0}\theta_{\rho}(x(\rho))\cdot s_\beta^{-1}\\
&=\theta_{-\beta}(x'(\beta))
\prod_{\rho>0, \rho\not=\beta}\theta_{\rho}(x'(\rho)),
\end{split}
\tag{2.17}
\]
where
\[
x'(\alpha)=\eta_{\alpha,\gamma}x(\gamma)+f_3x(\alpha)x(\beta)
\]
$x'(\beta)=\eta_{\alpha,\beta}x(\beta)$,
$x'(\gamma)=\eta_{\alpha,\alpha}x(\alpha)$,
\[
x'(\delta)=\eta_{\alpha,\delta}x(\delta)+f_4x'(\alpha)x'(\gamma),
\]
$x'(\epsilon)=\eta_{\alpha,\zeta}x(\zeta)$,
and
\[
x'(\zeta)=\eta_{\alpha,\epsilon}x(\epsilon).
\]

% Claude: UNTYPESET. The following text follows \enddocument in g2new2.tex
% and was never typeset. It reappears, essentially verbatim, as the
% discussion before Lemma 6.4 in Section 6.
\iffalse
How do these coordinates relate to the coordinates on a unipotent orbit
of $G$? In general this does not appear to be an easy question, but if the
orbit is a Richardson orbit then the answer is quite simple. The discussion
below is motivated by \cite[Remark V.2.2]{Ha92a}.

Let ${\mathcal O}$ be the Richardson orbit associated to a parabolic $P$
as in \cite[Theorem 2, \S 3.9]{St74}. Write $P=MN_P$, so that $N_P^G$
contains ${\mathcal O}$ as a dense subset, ${\mathcal O}\cap N_P$ is dense in
$N_P$ and ${\mathcal O}\cap N_P$ forms a single class under $P$. We have
dim${\mathcal O}$ = dim$G$ - dim$M$. Let $u\in {\mathcal O}\cap N_P$ and recall
$N_P^- P$ is dense in $G$ by the Bruhat decomposition. The $P$-orbit
$\{u^p\ |\ p\in P\}$ is equal to ${\mathcal O}\cap N_P$ and is dense in $N_P$.
Therefore, the coordinates of $N_P$ yield coordinates for the $P$-orbit
of $u\in {\mathcal O}\cap N_P$. The $G$-orbit of this $u$, that is
${\mathcal O}$ itself, contains the $N_P^- P$-orbit of $u$ as an open
subvariety. We may regard the $N_P^- P$-orbit of $u$ as the
$N_P^-$-orbit of ${\mathcal O}\cap N_P$. By Corollary 1.6, for each
$\overline{n}\in N_P^-$, $B^{\overline{n}}$ is opposite to
$B^-$ and $B^{\overline{n}}\not= B$ unless $\overline{n}=1$.
Therefore, for our $u$ as above, $u^{\overline{n}}\notin {\mathcal O}\cap N_P$.
From this it follows that the $N_P^-$-orbit of ${\mathcal O}\cap N_P$
is isomorphic to $N_P^-$.

\begin{proclaim}{Lemma 1.21} If ${\mathcal O}$ is the Richardson orbit associated to a
parabolic $P$ then the coordinates on $N_P^- \times N_P$ serve as local
coordinates on ${\mathcal O}$.
\end{proclaim}
\fi
% ======================================================================
% Claude: Section 3, source file g2new3.tex (\pageno 26). Here \off = {\Cal O}_F.
% ======================================================================
\section*{\S\,3 Construction of regular patches}

This section is a preliminary step to constructing the regular
variety $Y_C$ for $G_2$.

% Claude: In g2new3.tex \off is {\Cal O}_F, so the next sentence prints
% "over O_F". The phrase "over an algebraically closed field again" later in
% this section suggests that \overline{F} was intended.
Let $G$ denote an arbitrary connected semisimple group over $\OF$.
Let $\chi:G\rightarrow T_G/\Omega$ denote the adjoint quotient \cite{Sl80}.
There is an open patch
\[
Y''(B_+,B_-)\times N_+\subset
\chi^{-1}(C)\times S''(B_+,B_-)\times N_+
\subset G\times S''(B_+,B_-)\times N_+
\]
of $Y''$ which is isomorphic to an open set of the product
$\aaa^p\times Z$, where $Z\subset \aaa^{2p+1}$ is a
``monomial'' $p+1$ dimensional
affine variety having coordinates
\[
w(\rho),\quad (\rho\in \Sigma_G^+ -\Delta),\quad
z(\alpha)\quad (\alpha\in\Delta),\quad x(\rho)
\quad (\rho\in\Sigma_G^+),\quad \lambda,
\]
(we set $w(\alpha)=1,\ \alpha\in\Delta$), defined by the relations:
for $\lambda\not= 0$, $\rho\in\Sigma_G^+$, we have
\[
\lambda w(\rho )=x(\rho)\prod_{\alpha\in\Delta }z(\alpha )^{m(\alpha)},
\tag{3.1}
\]
where $m(\alpha)=m_\rho(\alpha)\geq 0$ is the integral coefficient in
$\rho = \sum_{\alpha\in\Delta }m(\alpha)\alpha $.
Here $\Delta\subset \Sigma_G$ denotes a basis of simple roots
of $T_+$ in $B_-$.
There is a surjection $Y''(B_+)\rightarrow Z$
which allows us to pull-back open subvarieties $Z_U\subset Z$ to open
subvarieties $Y''_U\subset Y''$.
Moreover,
if $\rho =\alpha$ is simple then $\lambda =x(\alpha)z(\alpha)$. This variety
$Z$ is regular except possibly for a singularity contained
in the hypersurface $\lambda =0$.

We first claim that, for the case $G=G_2$,
the singularities of $Y''(B_+,B_-)$, if any, lie in
\[
K:\ \ \ x(\alpha)=x(\beta)=x(\gamma)=z(\alpha)=w(\zeta)=0
\]
or in
\[
K_0:\ \ \ x(\alpha)=x(\beta)=x(\gamma)=x(\delta)=x(\epsilon)=x(\zeta)=w(\zeta)=0.
\]
The points in $Y''(B_+,B_-)$ with coordinates in the latter set all map to the
identity in $G$ under $\pi$, hence correspond to the trivial unipotent orbit.
Therefore, blowing up $Z$ along $K_0$ creates exceptional divisors
associated with the trivial orbit, which do not concern us here.

To prove this claim, we exhibit 5 regular patches, give generators of the
local coordinate ring, and express the remaining coordinates as regular
functions of the generators on that patch.

\medskip
\noindent
$x(\alpha)\not= 0$: Generators: $x(\alpha)$, $x(\beta)$, $z(\beta)$, $x(\gamma)$, $x(\delta)$, $x(\epsilon)$, $x(\zeta)$. Remaining coordinates:
$\lambda=x(\beta)z(\beta)$, $z(\alpha)=x(\beta)z(\beta)/x(\alpha)$,
$w(\gamma)=x(\gamma)z(\beta)/x(\alpha)$,
$w(\delta)=x(\delta)z(\alpha)z(\beta)/x(\alpha)$,
$w(\epsilon)=x(\epsilon)z(\alpha)^2 z(\beta)/x(\alpha)$,
$w(\zeta)=x(\zeta)z(\alpha)^2 z(\beta)^2 /x(\alpha)$.

\medskip
\noindent
$x(\beta)\not= 0$: Generators: $x(\alpha)$, $x(\beta)$, $z(\alpha)$, $x(\gamma)$, $x(\delta)$, $x(\epsilon)$, $x(\zeta)$. Remaining coordinates:
$\lambda=x(\alpha)z(\alpha)$, $z(\beta)=x(\alpha)z(\alpha)/x(\beta)$,
$w(\gamma)=x(\gamma)z(\alpha)/x(\beta)$,
$w(\delta)=x(\delta)z(\alpha)^2 /x(\beta)$,
$w(\epsilon)=x(\epsilon)z(\alpha)^3 /x(\beta)$,
$w(\zeta)=x(\zeta)z(\alpha)^5 x(\alpha)^2 /x(\beta)^3$.

\medskip
\noindent
$x(\gamma)\not= 0$: Generators: $x(\alpha)$, $x(\beta)$, $x(\gamma)$, $w(\gamma)$, $x(\delta)$, $x(\epsilon)$, $x(\zeta)$. Remaining coordinates:
$\lambda=x(\alpha)x(\beta)w(\gamma)/x(\gamma)$,
$z(\alpha)=x(\beta)w(\gamma)/x(\gamma)$, $z(\beta)=x(\alpha)w(\gamma)/x(\gamma)$,
$w(\delta)=x(\delta)w(\gamma)^2 x(\beta)/x(\gamma)^2$,
$w(\epsilon)=x(\epsilon)w(\gamma)^3 x(\beta)^2 /x(\gamma)^3$,
$w(\zeta)=x(\zeta)x(\alpha) x(\beta)^2 w(\gamma)^4 /x(\gamma)$.

\medskip
\noindent
$w(\zeta)\not= 0$: Generators: $z(\alpha)$, $z(\beta)$, $w(\gamma)$, $w(\delta)$, $w(\epsilon)$, $w(\zeta)$, $x(\zeta)$. Remaining coordinates:
$\lambda=x(\zeta)z(\alpha)^3 z(\beta)^2 /w(\zeta)$,
$x(\alpha)=x(\zeta)z(\alpha)^2 z(\beta)^2 /w(\zeta)$,
$x(\beta)=x(\zeta)z(\alpha)^3 z(\beta) /w(\zeta)$,
$x(\gamma)=w(\gamma)x(\zeta) z(\alpha)^2 z(\beta) /w(\zeta)$,
$x(\delta)=w(\delta)x(\zeta) z(\alpha) z(\beta) /w(\zeta)$,
$x(\epsilon)=w(\epsilon)x(\zeta)z(\beta) /w(\zeta)$.

\medskip
\noindent
$z(\alpha)\not= 0$ (assuming $x(\epsilon)\not=0$):
Generators: $z(\alpha)$, $z(\beta)$, $x(\beta)$, $w(\gamma)$, $w(\delta)$, $w(\epsilon)$, $x(\zeta)$. Remaining coordinates:
$\lambda=x(\beta)z(\beta)$, $x(\alpha)=x(\beta)z(\beta)/z(\alpha)$,
$x(\gamma)=w(\gamma)x(\beta)/z(\alpha)$,
$x(\delta)=w(\delta)x(\beta)/z(\alpha)^2$,
$x(\epsilon)=w(\epsilon)x(\beta)/z(\alpha)^3$,
$w(\zeta)=\frac{x(\zeta)}{x(\epsilon)} w(\epsilon)z(\beta) /z(\alpha)^3$.
We remark that the assumption $x(\epsilon)\not=0$ may be
replaced by $x(\delta)\not= 0$ or $x(\zeta)\not= 0$
without affecting regularity (though
of course the generators, etc., will change).

The above claim follows.

As remarked above, the ``interesting'' singularities of $Z$
lie in the subvariety $K$ above, not in $K_0$ since the exceptional
divisor associated to the blow-up along $K_0$ corresponds to the trivial orbit,
whose matching has be deduced from the
work of Rogawski, Howe, and Harish-Chandra \cite{Ro80}.

By Hironaka's embedded resolution of
singularities theorem (see also Mumford's semi-stable reduction theorem
\cite{Ke72}), we may obtain a smooth variety
$\tilde{Z}\rightarrow Z$, proper over $Z$, by blowing
up a finite number of times, beginning with $K$
as the center of the first blow up. Moreover, the exceptional
divisor $E_Z$ describing $\lambda =0$ in
$\tilde{Z}$ will have normal crossings.
Let $\tilde{Y''}\rightarrow Y''$ be the
corresponding resolution of $Y''$ with exceptional divisor
$E_{Y''}\subset \tilde{Y''}$; the branches of $E_{Y''}$ are reduced and
cross normally. This can (and has) been done explicitly.
However, as the details need not concern us here they will be omitted.

Let $G$ be an arbitrary connected semisimple
group over an algebraically closed field again.
Let $\Delta=\{\alpha_1,...,\alpha_\ell\}$, which we sometimes identify
with $\{1,...,\ell\}$, $\ell :=rk\, G$. Let $S_1,S_2\subset \{1,...,\ell\}$
be disjoint sets, and if they are
non-empty then let $i_{min}:=min_{i\in S_1}i$,
$j_{min}:=min_{j\in S_2}j$, and similarly for $i_{max},\ j_{max}$. If
$S_1,S_2$ are non-empty, not only do we assume
they are disjoint but also that
(i) $i_{max}<j_{min}$ or $j_{max}<i_{min}$
and (ii) if $S_2\not=\emptyset$ then
$S_1\cup S_2\not=\Delta$.

For each map $\sigma :S_2 \rightarrow \Sigma_G^+ -\Delta$,
define the subset $Z_{S_1,S_2,\sigma}\subset Z$ by

\noindent
(3.2a) $z(\alpha)\not= 0$, for all
$\alpha\notin S_1\cup S_2$,

\noindent
(3.2b) $x(\alpha)\not= 0$, for all $\alpha\in S_1$,

\noindent
(3.2c) $x(\sigma(j))\not= 0$, for all $j\in S_2$,

\noindent
(3.2d) $w(\sigma(j))\not= 0$, for all $j\in S_2$.

\noindent
These are \textbf{fundamental patches} of $Z$ in the sense of \cite{Ha92a}.
This collection of patches depends, a priori, on the
identification of $\Delta$ with $\{1,...,\ell\}$ above.
For $G=G_2$, they will turn out to have the
nice property that the divisor $\lambda =0$ on $Z$ will only meet
$Z_{S_1,S_2,\sigma}$ along a single irreducible branch which we will denote
by $E_{S_1,S_2,\sigma}$, possibly
occurring with some multiplicity $a(E_{S_1,S_2,\sigma})$.
The $E_{S_1,S_2,\sigma}$ are called the \textbf{fundamental branches} of
$\lambda =0$ on $Z$. If $S_2=\{\rho\}$ is a singleton and
$\tau=\sigma(\rho)$ then for brevity
we write $E_{S_1,\rho,\tau}$ in place of $E_{S_1,S_2,\sigma}$.
If $S_2 =\emptyset$ the we abbreviate this patch by
$Z_{S_1}$ and call it a \textbf{Richardson patch}.
Similarly for $E_{S_1}$, which we call a \textbf{Richardson branch}.
(I call the $Z_{S_1}$
``Richardson patches'', since by Proposition 9.1 of Ch II of \cite{Ha92a}
they will lead to divisors $E$ of $Y$ for which
$\pi(E)$ is the Zariski closure of a Richardson orbit.) Notice that the
\textbf{regular branch} $E_\emptyset$ ($S_1=S_2=\emptyset$) satisfies
$\pi(E_\emptyset)=\overline{\mathcal O}_{reg}=V$ where $V$ denotes the unipotent
variety of $G$ and ${\mathcal O}_{reg}$ denotes the regular orbit of
$G$. Since $E_\emptyset\subset Y_C$ is regular,
the map $\pi:E_\emptyset \rightarrow V$ yields a resolution of singularities
% Claude: [Sp] is not a label in the bibliography of the source (which has
% [Sp1] and [Sp2]); the citation is typeset literally.
similar to the Springer resolution \cite{Sl80}, [Sp], \cite{St74}.

Fix a map $\sigma :S_2 \rightarrow \Sigma_G^+ -\Delta$ satisfying

\noindent
(3.3a) for $j,j'\in S_2$, the simple root expansion of $\sigma(j)$ and
$\sigma(j')$ contain no common simple roots unless $j=j'$ or
$j,j'$ are consecutive.

\noindent
(3.3b) If $j,j'$ are consecutive in $S_2$ then
\[
\sigma(j)=\sum_{k=j}^{j'}m_{jk}\alpha_k,
\]
for some integers $m_{jk}\geq 0$.

\noindent
(3.3c) There is an $a\geq 1$ for which $\sigma(j_{max})
=a\alpha_{j_{max}}+\sum_{\alpha\notin S_1\cup S_2}m(\alpha)\alpha$,
for some integers $m(\alpha)\geq 0$.

Such a map will be called \textbf{essential} (in the sense that there
are lots of other possible patches $Z_{S_1,S_2,\sigma}$ but the
patches which are not essential are expected to be
redundant in the covering of $Z$ by patches up to a set of codimension $2$).
The corresponding fundamental patches are called \textbf{essential}.
This terminology is admittedly a bit optimistic: there are, in general,
many extraneous essential patches. Therefore this definition is
only provisional. In any case, these patches generalize the
sequences $\gamma(i,j),...$ of roots used in the paper \cite{Ha89}.

\begin{proclaim}{Notation} The notation for these divisors is not
exactly the same as that of \cite{Ha92a}, \cite{Ha89}, ... , though the notation
for the patches seems to agree. Hales' $E_S$ is usually my $E_{\Delta -S}$.
\end{proclaim}

Rather than discuss these patches in greater generality
than is necessary for
our purposes here, we merely state the following

\begin{proclaim}{Conjecture 3.4} If $G$ is not of type $D_n$, $E_n$, $F_4$ then
$Z$ is covered by the $Z_{S_1,S_2,\sigma}$, except for
a subset of codimension $\geq 2$. Moreover, the $Z_{S_1,S_2,\sigma}$
are all regular and $Z$ is regular in codimension one.
\end{proclaim}

In the case of $G_2$ this follows from the next two lemmas.
The conjecture is known to be true in the case of $A_n$ and
$C_n$ \cite{Ha92a}, \cite{Ha89}, and false in the case of $D_n$.

\begin{proclaim}{Lemma 3.5 (Hales \cite{Ha92a})}
For $S_1\subset \Delta$ non-empty, $Z_{S_1}$ is regular, with
generators $\lambda,\  z(\alpha)\ (\alpha\in\Delta -S_1),
\  x(\alpha)\ (\alpha\in S_1),\ x(\rho)\
(\rho\in\Sigma_G^+-\Sigma_{S_1}^+),
\ w(\rho)\ (\rho\in \Sigma_{S_1}^+)$, where $\Sigma_{S_1}^+$ denotes the
subroot system generated by $S_1\subset \Delta$.
It is therefore $p+1$ dimensional and the divisor $\lambda =0$ occurs with
multiplicity one.
\end{proclaim}

It is also easy to prove the conjecture when $S_1$ is empty and $S_2$ is
a singleton.
For brevity, let $m_\alpha$ denote the multiplicity with
which $z(\alpha)=0$ occurs in $\lambda =0$ and let
$\Delta_{\lambda =0}:=\{\beta\in\Delta\ |\ m_{\beta}\not= 0\}$.

\begin{proclaim}{Lemma 3.6} For any connected semi-simple group $G$, any
simple root $\alpha\in\Delta$, and any non-simple positive root
$\sigma\in\Sigma_G^+ -\Delta$, if
$Z_{\emptyset,\alpha,\sigma}$ is an essential patch then
$Z_{\emptyset,\alpha,\sigma}$ is regular and its coordinate ring has generators
\[
\begin{split}
&z(\beta )\ \ (\beta\in\Delta ),\\
&x(\rho) \ \ \ \ \text{ if }
\rho >\sigma, \\
&w(\rho) \ \ \ \
\text{ if }
\rho \leq \sigma
\text{ and }\rho\notin\Delta.
\end{split}
\]
Moreover, the divisor $\lambda =0$ occurs as $z(\alpha)=0$
with multiplicity $m$, where $m\geq 1$ is the integral coefficient of
$\alpha$ in the simple root expansion of $\sigma$.
\end{proclaim}

\begin{proof} Let $\sigma=m\alpha +\sum_{\beta\not= \alpha}m_\beta\beta$
and notice that we may assume that
$\sigma\notin \Delta$ for otherwise the patch is Richardson.

Recall the definition of $Z_{\emptyset,S_2,\sigma}$: $z(\beta)\not= 0$
for $\beta\not=\alpha$ simple, $x(\sigma)\not= 0$, and
$w(\sigma)\not= 0$. By (3.3c) and our hypothesis,
$\lambda w(\sigma)
=x(\sigma)z(\alpha)^m\prod_{\beta\not= \alpha}z(\beta)^{m_\beta}$,
it is clear that
$\lambda=0$ occurs as $z(\alpha)=0$
with multiplicity $m_\alpha$. Along $\lambda=0$ the other $z(\beta)'s$
are non-zero,
and the $\{z(\beta)\ |\ \beta\in\Delta\}$
give us $\ell$ independent coordinates. We need
$p+1-\ell$ more having the property that the remaining coordinates
are regular functions of these along $\lambda =0$. Suppose
$\rho \leq\sigma$, so $\sigma=\rho +\sum_\beta m(\beta)\beta$.
We have
\[
x(\rho)=\frac{x(\sigma)w(\rho)}{w(\sigma)}
\prod_\beta z(\beta)^{m(\beta)},
\]
and therefore we want to choose the $w(\rho)$ with $\rho \leq\sigma$,
$\rho\notin \Delta$,
as generators. Likewise, we want to choose the $x(\rho)$ with
$\rho \geq\sigma$ as generators since a similar line of reasoning applied
to the quotients
\[
\frac{w(\rho)}{w(\sigma)}=\frac{x(\rho)}{x(\sigma)}\prod_\beta
z(\beta)^{m_\rho(\beta)-m_{\sigma}(\beta)}.
\]
This yields the claim.
\end{proof}

\begin{proclaim}{Lemma 3.7} For any connected semisimple group,
the variety $Z$ is normal and Cohen-Macauley.
Moreover, $Z$ is a set-theoretic complete intersection.
\end{proclaim}

\begin{proof} $Z$ is ``monomial'' and hence has a toroidal
embedding \cite[Ch I]{Ke72}. From the definition of saturated
\cite[Ch I, Definition 2]{Ke72} and the equations (3.1)
defining $Z$, we see that any toroidal embedding
$T\hookrightarrow Z$ corresponds to a saturated sub-semigroup of
$X^*(T)=Hom_{alg.-grps.}(T,{\mathbb G}_m)$. $Z$ is therefore
normal by \cite[Ch I, Theorem 1]{Ke72} and Cohen-Macauley
by \cite[Ch I, Theorem 14]{Ke72}.
The fact that $Z$ is a set-theoretic
complete intersection follows from the definition
and counting the number of relations defined by (3.1).
\end{proof}
% ======================================================================
% Claude: Section 4, source file g2new4.tex (\pageno 30). Here \off = {\Cal O}_F
% (not used in this section).
% ======================================================================
\section*{\S\,4 Regularity of patches for $G=G_2$}

It is not difficult to write all these patches of \S 3
down and show that they are
regular for $G_2$. Recall $\Sigma_G^+=\{\alpha,\beta,
\gamma:=\alpha+\beta,\delta:=2\alpha +\beta,
\epsilon:=3\alpha +\beta,\zeta:=3\alpha +2\beta\}$, with
$\alpha$ short and $\beta$ long, so $p=6$ and dim $Z=7$. We have variables
\[
\begin{split}
&x(\alpha),x(\beta),
x(\gamma),x(\delta),x(\epsilon),x(\zeta)\\
&z(\alpha),z(\beta),
w(\gamma),w(\delta),w(\epsilon),w(\zeta)\\
&\lambda,
\end{split}
\tag{4.1}
\]
subject to the usual relations (3.1).

With $B_+$ choosen as in \S 3, we define $\omega_{Y''(B_+)}$ to be the
$G$-invariant form
\[
\omega_{Y''(B_+)}=d\lambda \wedge
\wedge
dx(\alpha)\wedge dx(\beta)
\wedge dx(\gamma)\wedge dx(\delta)\wedge dx(\epsilon)\wedge dx(\zeta)
\wedge dv,
\]
where $dv$ is a form on $N_+$. We use this to compute the Igusa constants
$\beta$ in this section as in \cite{Ha89}.

As we know (from lemma 3.5), the Richardson patches $Z_{S_1}$ are regular.
Now let us consider the remaining patches. Since $S_1\cup S_2\not= \Delta$
if $S_2\not= \emptyset$, it is easy to list
them all: $S_1=\emptyset$, $S_2=\{\alpha\}$, where
$\sigma(\alpha)=\alpha,\gamma,\delta,\epsilon,\zeta$, or
$S_1=\emptyset$, $S_2=\{\beta\}$, where
$\sigma(\beta)=\beta,\gamma,\delta,\epsilon,\zeta$.
This may be deduced from the discussion above. We next list the generators of
the local coordinate ring in each case:

\noindent
$S_2=\{\alpha\}, \sigma(\alpha)=\alpha$: Same as $Z_{\{\alpha\}}$. The divisor
$\lambda =0$ occurs with multiplicity one.

\noindent
$S_2=\{\alpha\}, \sigma(\alpha)=\gamma$: Generators: $z(\alpha), z(\beta),
x(\gamma), w(\gamma), x(\delta), x(\epsilon),
x(\zeta)$. Regular at $\lambda =0$. The divisor $\lambda =0$
occurs with multiplicity one (represented by $z(\alpha)=0$).

\noindent
$S_2=\{\alpha\}, \sigma(\alpha)=\delta$: Generators: $z(\alpha), z(\beta),
w(\gamma), x(\delta),w(\delta), x(\epsilon),
x(\zeta)$. Regular at $\lambda =0$. The divisor $\lambda =0$
occurs (as $z(\alpha)=0$) with multiplicity two.

\noindent
$S_2=\{\alpha\}, \sigma(\alpha)=\epsilon$: Generators: $z(\alpha), z(\beta),
w(\gamma), w(\delta),w(\epsilon),x(\epsilon),x(\zeta)$.
Regular at $\lambda =0$. The divisor $\lambda =0$
occurs (as $z(\alpha)=0$) with multiplicity three.

\noindent
$S_2=\{\alpha\}, \sigma(\alpha)=\zeta$: Generators: $z(\alpha), z(\beta),
w(\gamma), w(\delta),w(\epsilon), w(\zeta),
x(\zeta)$. Regular at $\lambda =0$. The divisor $\lambda =0$
occurs (as $z(\alpha)=0$) with multiplicity three.

\noindent
$S_2=\{\beta\}, \sigma(\beta)=\beta$: Same as $Z_{\{\beta\}}$.
The divisor $\lambda =0$
occurs with multiplicity one.

\noindent
$S_2=\{\beta\}, \sigma(\beta)=\gamma$: Generators: $z(\alpha), z(\beta),
x(\gamma), w(\gamma), x(\delta), x(\epsilon),
x(\zeta)$. Regular at $\lambda =0$. The divisor $\lambda =0$
occurs with multiplicity one (represented by $z(\beta)=0$).

\noindent
$S_2=\{\beta\}, \sigma(\beta)=\delta$: Generators: $z(\alpha), z(\beta),
w(\gamma), x(\delta),w(\delta), x(\epsilon),
x(\zeta)$. Regular at $\lambda =0$. The divisor $\lambda =0$
occurs (as $z(\beta)=0$) with multiplicity one.

\noindent
$S_2=\{\beta\}, \sigma(\beta)=\epsilon$: Generators: $z(\alpha), z(\beta),
w(\gamma), w(\delta),w(\epsilon),x(\epsilon),x(\zeta)$.
Regular at $\lambda =0$. The divisor $\lambda =0$
occurs (as $z(\beta)=0$) with multiplicity one.

\noindent
$S_2=\{\beta\}, \sigma(\beta)=\zeta$: Generators: $z(\alpha), z(\beta),
w(\gamma), w(\delta),w(\epsilon), w(\zeta),
x(\zeta)$. Regular at $\lambda =0$. The divisor $\lambda =0$
occurs (as $z(\beta)=0$) with multiplicity two.

There is some redundancy here since $E_{\emptyset,\{\beta\},\sigma}=
E_{\{\beta\}}$ if $\sigma(\beta)=\beta,\gamma,\delta,\epsilon$.

We denote the (irreducible) branch of $\{\lambda =0\}\subset Z$
occuring in $Z_{S_1,S_2,\sigma}$ by $E_{S_1,S_2,\sigma}$. The same notation
will be used for the corresponding branch in $Y_U$ and for its Zariski closure
in $Y_1$. To determine which of the $\kappa$-stable germs, if any,
the divisor $E_{S_1,S_2,\sigma}$
contributes to by formula 0.6,
we must compute the Igusa constant
\[
\beta(\tilde{E}_{S_1,S_2,\sigma})=
\frac{b(\tilde{E}_{S_1,S_2,\sigma})}{a(\tilde{E}_{S_1,S_2,\sigma})}
\tag{4.2}
\]
where $\tilde{E}$ denotes the preimage of $E$ in the
desingularized variety $Y_C$.

We record two observations regarding the Igusa constants $\beta(E)$.

\noindent
(I)
First, if $\phi:V\rightarrow W$ is an isomorphism of smooth varieties
which sends a divisor $E\subset V$ to a divisor $F=\phi(E)\subset W$
then $\beta(E)=\beta(F)$.

\noindent
(II)
Secondly, let $V'$ denote a blow up of a variety
$V$ with exceptional divisor $E_{V'}$.
Recall that for each irreducible branch $E$ on $V$,
the pull-back ${\xi'}^*$ of the blow down map $\xi':V'\rightarrow V$
satisfies ${\xi'}^*(E)=\tilde{E}+\delta E_{V'}$, where $\tilde{E}$
is the strict transform
of $E$, and
\[
\delta:=
\begin{cases} 1, &\text{ if $E$ contains the center of the blow up},\\
0, &\text{ otherwise}.
\end{cases}
\]
A repeated use of this global identity can be used to obtain
local information about the
Igusa constants on $Z$.

For example, if $E$ denotes a branch of $\{\lambda =0\}\subset Z$ and
$\tilde{E}$ its strict transform in a desingularization $\tilde{Z}$,
we have
\[
\beta(\tilde{E})=\beta(E),
\]
provided $E$ does not contain any of the images of the centers of
the blow ups used to create $\tilde{Z}$. The only other divisors
in $\{\lambda =0\}\subset \tilde{Z}$ other than the strict
transforms $\tilde{E}$ above are the exceptional divisors.
To compute all the $\beta$-values of the
branches of $\phi^{-1}(1)$ in $Y_C$ we need only compute the $\beta$-values
of the divisors $E_{S_1,S_2,\sigma}$ on $Z$ and the branches of
the exceptional divisors of the blow-up.

We digress to
compute $\beta({E}_{S_1,S_2,\sigma})$ on $Z_{S_1,S_2,\sigma}$. We
need not compute the $\beta$-values for the exceptional divisors
because of lemma 0.8 and the fact that all the needed $F$-rational
${\mathcal O}$-divisors can be accounted for among the $E_{S_1,S_2,\sigma}$.

If $S_2=\emptyset$ then $\beta({E}_{S_1,S_2,\sigma})$
has been computed in \cite[ch II, Proposition 9.1]{Ha92a}
so some of our calculations are covered by Hales' general
calculations. For completeness, we list everything anyway.

\subsection*{Igusa forms on the fundamental patches of $Y''$}

We pull-back the patch $Z_{S_1,S_2,\sigma}$
to $Y''(B_+,B_-)$ via the open immersion
$Y''(B_+,B_-)\rightarrow Z$. Call this inverse image $Y''_{S_1,S_2,\sigma}$.
The Igusa form $\omega_U$ on $U=Y''_{S_1,S_2,\sigma}$ is related to the
Igusa form $\omega_{Z,U}$ on $Z_{S_1,S_2,\sigma}$ by
$\omega_U=\omega_{Z,U}\wedge dv$,
where $dv$ is regular. We compute $\omega_{Z,U}$below.

$Z_{\{\alpha,\beta\}}$: Here the multiplicity $a(E_{\{\alpha,\beta\}})=1$
and $b(E_{\{\alpha,\beta\}})=1$ since
\[
\omega_{Z,U}
=d\lambda\wedge dx(\alpha)\wedge dx(\beta)\wedge dx(\gamma)
\wedge dx(\delta)\wedge dx(\epsilon)\wedge dx(\zeta) .
\]
Therefore, $\beta(E_{\{\alpha,\beta\}})=1$.

\medskip
$Z_{\alpha}$: Here the multiplicity $a(E_{\alpha})=1$
and $b(E_{\alpha})=2$ since
\[
\begin{split}
\omega_{Z,U}
&=d\lambda\wedge dx(\alpha)\wedge dx(\beta)\wedge dx(\gamma)
\wedge dx(\delta)\wedge dx(\epsilon)\wedge dx(\zeta)\\
&=-\frac{\lambda}{z(\beta)^2}
d\lambda\wedge dx(\alpha)\wedge dz(\beta)\wedge dx(\gamma)
\wedge dx(\delta)\wedge dx(\epsilon)\wedge dx(\zeta).
\end{split}
\]
Therefore, $\beta(E_{\alpha})=2$.

\medskip
$Z_{\beta}$: Here the multiplicity $a(E_{\beta})=1$
and $b(E_{\beta})=2$ since
\[
\begin{split}
\omega_{Z,U}
&=d\lambda\wedge dx(\alpha)\wedge dx(\beta)\wedge dx(\gamma)
\wedge dx(\delta)\wedge dx(\epsilon)\wedge dx(\zeta)\\
&=-\frac{\lambda}{z(\alpha)^2}
d\lambda\wedge dz(\alpha)\wedge dx(\beta)\wedge dx(\gamma)
\wedge dx(\delta)\wedge dx(\epsilon)\wedge dx(\zeta).
\end{split}
\]
Therefore, $\beta(E_{\beta})=2$.

\medskip
$Z_{\emptyset,\alpha,\sigma}, \ \sigma(\alpha)=\gamma$:
Here the multiplicity $a(E_{\emptyset,\alpha,\sigma})=1$, by the
above list. Also, $b(E_{\emptyset,\alpha,\sigma})=2$ since
\[
\begin{split}
\omega_{Z,U}
&=d\lambda\wedge dx(\alpha)\wedge dx(\beta)\wedge dx(\gamma)
\wedge dx(\delta)\wedge dx(\epsilon)\wedge dx(\zeta)\\
&=c\frac{x(\gamma)^3z(\beta)z(\alpha)}{w(\gamma)^4}
dz(\alpha)\wedge dz(\beta)\wedge dw(\gamma)\wedge dx(\gamma)
\wedge dx(\delta)\wedge dx(\epsilon)\wedge dx(\zeta).
\end{split}
\]
Therefore, $\beta(E_{\emptyset,\alpha,\sigma})=2$.

\medskip
$Z_{\emptyset,\alpha,\sigma}, \ \sigma(\alpha)=\delta$:
Here the multiplicity $a(E_{\emptyset,\alpha,\sigma})=2$, by the
above list. Also, $b(E_{\emptyset,\alpha,\sigma})=6$ since
\[
\begin{split}
\omega_{Z,U}
&=d\lambda\wedge dx(\alpha)\wedge dx(\beta)\wedge dx(\gamma)
\wedge dx(\delta)\wedge dx(\epsilon)\wedge dx(\zeta)\\
&=c\frac{x(\delta)^4z(\beta)z(\alpha)^5}{w(\delta)^5}
dz(\alpha)\wedge dz(\beta)\wedge dw(\gamma)\wedge dw(\delta)
\wedge dx(\delta)\wedge dx(\epsilon)\wedge dx(\zeta).
\end{split}
\]
Therefore, $\beta(E_{\emptyset,\alpha,\sigma})=3$.

\medskip
$Z_{\emptyset,\alpha,\sigma}, \ \sigma(\alpha)=\epsilon$:
Here the multiplicity $a(E_{\emptyset,\alpha,\sigma})=3$, by the
above list. Also, $b(E_{\emptyset,\alpha,\sigma})=10$ since
\[
\begin{split}
\omega_{Z,U}
&=d\lambda\wedge dx(\alpha)\wedge dx(\beta)\wedge dx(\gamma)
\wedge dx(\delta)\wedge dx(\epsilon)\wedge dx(\zeta)\\
&=c\frac{x(\epsilon)^5z(\beta)z(\alpha)^9}{w(\delta)^6}
dz(\alpha)\wedge dz(\beta)\wedge dw(\gamma)\wedge dw(\delta)
\wedge dw(\epsilon)\wedge dx(\epsilon)\wedge dx(\zeta).
\end{split}
\]
Therefore, $\beta(E_{\emptyset,\alpha,\sigma})=10/3$.

\medskip
$Z_{\emptyset,\alpha,\sigma}, \ \sigma(\alpha)=\zeta$:
Here the multiplicity $a(E_{\emptyset,\alpha,\sigma})=3$, by the
above list. Also, $b(E_{\emptyset,\alpha,\sigma})=11$ since
\[
\begin{split}
\omega_{Z,U}
&=d\lambda\wedge dx(\alpha)\wedge dx(\beta)\wedge dx(\gamma)
\wedge dx(\delta)\wedge dx(\epsilon)\wedge dx(\zeta)\\
&=c\frac{x(\zeta)^6z(\beta)^7z(\alpha)^{10}}{w(\delta)^7}
dz(\alpha)\wedge dz(\beta)\wedge dw(\gamma)\wedge dw(\delta)
\wedge dw(\epsilon)\wedge dw(\zeta)\wedge dx(\zeta).
\end{split}
\]
Therefore, $\beta(E_{\emptyset,\alpha,\sigma})=11/3$.

\medskip
$Z_{\emptyset,\beta,\sigma}, \ \sigma(\beta)=\zeta$:
Here the multiplicity $a(E_{\emptyset,\beta,\sigma})=2$, by the
above list. Also, $b(E_{\emptyset,\beta,\sigma})=8$ since
\[
\begin{split}
\omega_{Z,U}
&=d\lambda\wedge dx(\alpha)\wedge dx(\beta)\wedge dx(\gamma)
\wedge dx(\delta)\wedge dx(\epsilon)\wedge dx(\zeta)\\
&=c\frac{x(\zeta)^6z(\beta)^7z(\alpha)^{10}}{w(\delta)^7}
dz(\alpha)\wedge dz(\beta)\wedge dw(\gamma)\wedge dw(\delta)
\wedge dw(\epsilon)\wedge dw(\zeta)\wedge dx(\zeta).
\end{split}
\]
Therefore, $\beta(E_{\emptyset,\beta,\sigma})=4$.

\medskip
We summarize this in the following table:
\[
\begin{array}{cccc}
S_1&S_2&\sigma&\beta(E_{S_1,S_2,\sigma})-1\\
 & & &\\
\alpha,\beta&\emptyset&------&0\\
\alpha&\emptyset&------&1\\
\beta&\emptyset&------&1\\
\emptyset&\emptyset&------&\text{trivial}\\
\emptyset&\alpha&\alpha&1\\
\emptyset&\alpha&\gamma&1\\
\emptyset&\alpha&\delta&2\\
\emptyset&\alpha&\epsilon&7/3\\
\emptyset&\alpha&\zeta&8/3\\
\emptyset&\beta&\beta&1\\
\emptyset&\beta&\gamma&1\\
\emptyset&\beta&\delta&1\\
\emptyset&\beta&\epsilon&1\\
\emptyset&\beta&\zeta&3.
\end{array}
\tag{4.3}
\]
As was remarked above, there is some redundancy here... .

Let $H_1=PGL(3)$
and $H_2=(SL(2)\times SL(2))/(\pm 1)$.

\begin{proclaim}{Lemma 4.4} In the notation of Lemmas 3.5 and 3.6, we have
\[
\beta(E_S)=\frac{b(E_S)}{a(E_S)}=1+|\Sigma_S^+|=1+\frac{dim(Cent_G(u))-2}{2},
\]
where $u$ is an element of the Richardson orbit $\ooo_S$ associated to the
parabolic $P_S$ and the divisor $E_S$ is the
Richardson divisor contained in the patch $Z_S$.
The Igusa constants for
\[
E_{\ooo}=
\begin{cases}
E_{\emptyset,\alpha,\delta}, & \ooo \text{ is $2$-regular},\\
E_{\emptyset,\beta,\zeta}, &  \ooo \text{ is $3$-regular},
\end{cases}
\]
are given by
\[
\beta(E_{\ooo})=
\frac{b(E_{\ooo})}{a(E_{\ooo})}
=1+|\Sigma_{H_i}^+|=1+\frac{dim(Cent_G(u))-2}{2},
\]
where $i=1,2$ is determined by the orbit as indicated below and
$u$ is an element of the induced orbit
\[
\ooo=
\begin{cases}
\Ind_{H_2}^G(1), &\ooo \text{ is $2$-regular},\\
\Ind_{H_1}^G(1), &\ooo \text{ is $2$-regular},
\end{cases}
\]
associated to the endoscopic $H_i$ and the divisor $E_{\ooo}$.
\end{proclaim}

\begin{proof} The first part of the lemma is a special case of
\cite[ch II, Proposition 9.1]{Ha92a} (it also follows
from our case-by-case calculations).
The second part is verified
by the calculations above. (For various
facts about unipotent orbits, including the definition of
induction used here, see \cite{Spa82}.)
\end{proof}

% Claude: UNTYPESET. The following text follows \enddocument in g2new4.tex
% and was never typeset.
\iffalse
\subsection*{Igusa forms on $Y_1$}

Each of the Igusa forms $\omega_U$ on $Y''(B_+,B_-)$ can be pulled-back to
$Y_1(B_+,B_-)$ via the coordinate relations (1.22).
This allows us to compute the dependence of $\omega_U$ on $X$.
The result of this computation is essentially
independent of the patch as we will see.

On the preimage $U=Y_{1,S_1,S_2,\sigma}$ of $Z_{S_1,S_2,\sigma}$ under
the open immersion $Y_{1,S_1,S_2,\sigma}\rightarrow Z$, we have
$\omega_{U}=\omega_{Z,U}\wedge dv$, where $dv$ is regular and
\[
\omega_{Z,U}=h(t_0,\lambda)d\lambda\wedge dn_\alpha\wedge dn_\beta\wedge
...\wedge dn_\zeta ,
\]
where $t_0$ is as in (1.17), $t_0=exp(\lambda X)$, and
\[
\begin{split}
h(t_0,\lambda)
&=(1-\alpha(t_0)^{-1})...(1-\zeta(t_0)^{-1})+
\text{ higher order terms in }\lambda \\
&=\lambda^6 \alpha(X)...\zeta(X)+ \text{ higher order terms in }\lambda.
\end{split}
\]
\fi
% ======================================================================
% Claude: Section 5, source file g2new5.tex (\pageno 35). Here \off = {\Cal O}_F
% (not used in this section). VARIANT file g2new5d.tex differs from g2new5.tex
% in exactly two places, both marked below.
% ======================================================================
\section*{\S\,5 The Igusa expansion for $G_2$}

\subsection*{Notes on the derivation of the Igusa expansion}

Consider an $F$-morphism
\[
\phi:Y\rightarrow {\mathbb A},
\]
smooth outside the ``special fiber'' $\phi^{-1}(0)$,
% Claude: VARIANT. g2new5d.tex reads "from a smooth $F$-manifold $Y$" here.
from a smooth $F$-variety $Y$ to the affine line over $F$.
When we regard $Y(F)$ and ${\mathbb A}(F)=F$ as $F$-manifolds
in the p-adic topology it will be explicitly stated.
The ``default'' topology (to use the computer terminology)
on $Y$ is the Zariski topology. Let $f\in C_c^\infty(Y(F))$,
considering $Y(F)$ as a p-adic manifold, and consider
\[
F(\lambda):=
\int_{\phi^{-1}(\lambda)(F)}f(y)|\omega_\lambda|,
\]
for $\lambda\in F-\{0\}$. If $f$ is locally a product of
quasi-characters on a suitable local coordinate system of
a neighborhoood of $\phi^{-1}(0)$ in the manifold $Y(F)$ then Igusa
has given an ``asymptotic expansion'' for $F(\lambda)$
as $\lambda\rightarrow 0$. Actually, since the topology on
$Y(F)$ is totally disconnected, this ``asymptotic expansion''
is actually an equality valid for all $\lambda$ in a neighborhood
of $0$. We explain this next.

Assume that $Y(F)$ may be covered by patched
$U$ on which there exist coordinates $\mu_1,...,\mu_n$
such that

\noindent
(5.1) $\phi^{-1}(0)$ has normal crossings; more specifically,
the irreducible components $E_i$ of $\phi^{-1}(0)\cap U$
are of the form $\{\mu_i =0\}$ on $U$ and
\[
\lambda=\alpha\prod_i\mu_i^{a_i},
\]
for some integral $a_i>0$, where $\alpha$ is a non-vanishing
function in a p-adic neighborhood of $0$ and
$|\alpha|=q^{-m(\alpha)}$ is constant for all sufficiently
small $\lambda$.

\noindent
(5.2) On $U$, $f(y)$ is of the form
\[
f(\mu_1,...,\mu_n)=\gamma \prod_i\kappa_i(\mu_i),
\]
where $\gamma$ is constant for sufficiently small $\lambda$.

\noindent
(5.3) the form $\omega$ is given on $U$ by
\[
\omega=\prod_i\mu_i^{b_i-1}d\mu_1...d\mu_n.
\]
(5.4) We \textbf{assume} that $f$ vanishes off $U$ and that if $b_i=0$ then
$\kappa_i\not= 1$, $1\leq i\leq n$.

Let $F(\lambda)$ be given by (0.2) above.
We will sketch the derivation of an equality of the form
\[
F(\lambda) =\sum_{\theta,\beta}\theta(\lambda)|\lambda|^{\beta -1}
\sum_{r=1}^{e(\theta,\beta)}(-\log_q|\lambda|)^{r-
1}F_r(\theta,\beta,f).
\tag{5.5}
\]
The following lemma helps us to compute $F_1$:

\begin{proclaim}{Lemma 5.6} Let
\[
G(\lambda) =\sum_{\theta,\beta}\sum_{r=1}^{e(\theta,\beta)}
a(\theta,\beta,r)(-\log_q|\lambda|)^{r-1}\theta(\lambda)|\lambda|^{\beta -1},
\]
where the $a(\theta,\beta,r)$ are constants. If
$a(\theta,\beta,r)=0$ for $r>1$ then
the coefficients $b_n$ in the power series of the
truncated Mellin transform
\[
M_\epsilon G(s):=\int_{|\lambda|\leq \epsilon}
G(\lambda)|\lambda|^sd\lambda =\sum_{n\geq 0}b_nq^{-ns},
\ \ \ \ \ \epsilon=q^{-M},
\]
are given by
\[
b_n=
\begin{cases}
vol({\mathcal O}_F^\times,d\lambda)
\sum_{\substack{ \theta,\beta\\ \theta\text{ unramified}}}
a(\theta,\beta,1)(\theta(\pi)q^{-\beta})^n,&n\geq M,\\
0,& n<M.
\end{cases}
\]
\end{proclaim}

The proof of this is left to the reader.

Now, to verify the formula (5.5) for $F(\lambda)$ we may
verify it for its truncated Mellin transform
and use the lemma above to check the coefficients of
both sides of the resulting power series in
$q^{-s}$. The formula we have in mind for $F_1$ is stated below.
First, we need some notation: in the notation of the assumptions
(5.1-4) above, let $c_j=c_j(a_1,...,a_n)$ denote
the coefficients for which the principal part of
$\prod_i\frac{1}{1-t^{a_i}}$ at $t=1$ is
\[
\sum_j\frac{c_j}{(1-t)^j}.
\]
Let
\[
A(x)=\sum_j \frac{c_j}{(j-1)!}(x+1)...(x+j-1),
\]
and define $A_i(x)$ by $A(x-y)=\sum_j x^{j-1}A_j(y)$,
so $A_1(y)$ is the ``constant term'' of this polynomial in $x$.
(The function $A(x)$ arises via the undergraduate calculus equality
\[
\frac{(j-1)!}{(1-t)^j}
=\sum_{n\geq 0}(n+1)(n+2)...(n+j-1)t^n.)
\]

\begin{proclaim}{Lemma 5.7} Suppose that $Y(F)$, $f$, $\omega$, and
$\lambda$ satisfy the conditions of (5.1-4). Suppose furthermore
that the support of $f$ is contained in one such patch
$U$, where
\[
U=\{(\mu_1,...,\mu_n)\in F^n\ |\ |\mu_i|\leq q^{-M_i}\}.
\]
Finally, suppose that we reindexed the $\mu_i$, if
necessary, so that
\[
\theta^{a_i}=\kappa_i,\ \ \ \ \beta=b_i/a_i,\ \ \ \
1\leq i\leq p,
\]
and for each $i>p$ at least one of these equalities fails to hold.
(This $p$ is of course different from the fixed prime $p$.)
Then $F(\lambda)$ satisfies the formula (5.5)
where $F_1(\theta,\beta,f)$ is given by the value of
\[
\int_{D_p\cap U}A_1(M)h_p\prod_{i=p+1}^n|\mu_i|^{z_i}|\nu_p|,
\]
at $z_{p+1}=...=z_n=0$, where
\[
M=m(\alpha)+\sum_{i\leq p}a_iM_i +\sum_{i>p}a_i(-\log_q|\mu_i|),
\]
\[
h_p=\frac{\gamma}{\theta(\alpha)}
\prod_{i=p+1}^n\kappa_i(\mu_i)\theta(\mu_i)^{-a_i},
\]
and
\[
\nu_p=\frac{1}{|\alpha|^\beta}
\prod_{i=p+1}^n|\mu_i|^{b_i-\beta a_i-1}d\mu_{p+1}...d\mu_n.
\]
Moreover, if, in addition, any one of the characters $\kappa_i\theta^{-a_i}$
is ramified then $F_1(\theta,\beta,f)=0$.
\end{proclaim}

We now sketch the proof of this lemma following \cite{La83}.
By our assumptions,
\[
\begin{split}
M_\epsilon (F\theta^{-1})(s)
&=\frac{|\alpha|^s}{\theta(\alpha)}
\int_Uf(y)\prod_i\theta(\mu_i)^{-a_i}|\mu_i|^{a_is+b_i-1}
d\mu_1...d\mu_n\\
&=\frac{|\alpha|^s}{\theta(\alpha)}\gamma
\prod_i\int_{|\mu|\leq q^{-M_i}}
\kappa_i(\mu)\theta(\mu)^{-a_i}|\mu|^{a_is+b_i-1}d\mu,
\end{split}
\tag{5.8}
\]
since $|\omega_\lambda ||d\lambda |
=\prod_{i=1}^n|\mu_i|^{b_i-1}d\mu_1...d\mu_n$.
Recall that
\[
\int_{|\mu|\leq q^{-M_i}}
\kappa'(\mu)|\mu|^zd\mu =0
\]
if $\kappa'$ is a ramified character. Therefore we may assume
from now on that every character $\kappa_i(\mu)\theta(\mu)^{-a_i}$
is unramified and we let
\[
\epsilon_i:=\kappa_i(\pi)\theta(\pi)^{-a_i},\ \ \ \
1\leq i\leq n,
\]
where $\pi$ denotes a fixed generator of the maximal ideal
of the local ring of integers ${\mathcal O}_F$. We may write
\[
M_\epsilon (F\theta^{-1})(s)
=\gamma\frac{|\alpha|^s}{\theta(\alpha)}
(1-q^{-1})^n\prod_i
\frac{(\epsilon_iq^{-a_is-b_i})^{M_i}}{1-\epsilon_iq^{-a_is-b_i}}.
\]
Recall that we have reindexed the $\mu_i$ so that
\[
\theta^{a_i}=\kappa_i,\ \ \ \ \beta=b_i/a_i,\ \ \ \
1\leq i\leq p,
\]
and for each $i>p$ at least one of these equalities fails to hold.
Then
\[
M_\epsilon (F\theta^{-1})(s)
=M(s,z_{p+1},...,z_n),
\]
at $z_{p+1}=...=z_n=0$, where
\[
\begin{split}
&M(s,z_{p+1},...,z_n)\\
&\quad\quad =(1-q^{-1})^p
[\prod_{1\leq i\leq p}
\frac{(q^{-a_is-b_i})^{M_i}}{1-q^{-a_is-b_i}}]|\alpha|^{s+\beta}
\int_{\substack{ |\mu_i|\leq q^{-M_j}\\ p+1\leq i\leq n}}
\times\\
&\quad\quad\quad\quad \times
h_p|\alpha|^{-\beta}
\prod_{p+1\leq i\leq n}|\mu_i|^{z_i+a_is+b_i-1}
d\mu_{p+1}...d\mu_n\\
&\quad\quad
=(1-q^{-1})^p|\alpha|^{s+\beta}\frac{\gamma}{\theta(\alpha)}
[\prod_{1\leq i\leq p}
\frac{(q^{-a_is-b_i})^{M_i}}{1-q^{-a_is-b_i}}]\times \\
&\quad\quad \quad\quad \times
|\alpha|^{-\beta}\int_{\substack{ |\mu_i|\leq q^{-M_j}\\ p+1\leq j\leq n}}
\prod_{p+1\leq i\leq n}|\mu_i|^{a_i(s+\beta)}
\kappa_i(\mu_i)\theta(\mu_i)^{-a_i}|\mu_i|^{z_i-a_i\beta+b_i-1}
d\mu_{p+1}...d\mu_n\\
&\quad\quad
=(1-q^{-1})^p|\alpha|^{s+\beta}[\prod_{1\leq i\leq p}
\frac{(q^{-a_is-b_i})^{M_i}}{1-q^{-a_is-b_i}}]\times \\
&\quad\quad \quad\quad \times
\int_{\substack{ |\mu_i|\leq q^{-M_j}\\ p+1\leq i\leq n}}
\prod_{p+1\leq i\leq n}|\mu_i|^{z_i+a_i(s+\beta)}h_p|\nu_p|.
\end{split}
\tag{5.9}
\]
The last integral may be interpreted as a principal-value
integral of a function on $Y(F)$ over a variety of the form
$D_p=\cap_{1\leq i\leq p}E_i$, where $E_i$ is an irreducible
component of $\phi^{-1}(0)$ which satisfies
$E_i\cap U=\{(\mu_1,...,\mu_n)\ |\ \mu_i=0\}$. Moreover, this
multiple integral is a rational function of $q^{-s}$.

To see what this rational function is, let for the moment
$t=q^{-s-\beta}$ and recall that we have defined the $c_j$ in
such a way that the principal part of
\[
\prod_{1\leq i\leq p}
\frac{(q^{-a_i(s+\beta)})^{M_i}}{1-q^{-a_i(s+\beta)}}=
\prod_{1\leq i\leq p}
\frac{t^{a_iM_i}}{1-t^{a_i}}
\]
at $t=1$ is
\[
t^{\sum_{1\leq i\leq p}a_iM_i}
\sum_{1\leq i\leq p}\frac{c_j}{(j-1)!}(1-t)^{-j}
=\prod_{p+1\leq i\leq n}|\mu_i|^{a_i(s+\beta)}|\alpha|^{-s-\beta}
t^M\sum_{n\geq 0}A(n)t^{n},
\]
where $M$ is as in the statement of the lemma.
(Here we have multiplied and divided by
$\prod_{p+1\leq i\leq n}|\mu_i|^{a_i(s+\beta)}\allowbreak
|\alpha|^{-s-\beta}$ to account for all of $M$.)
The coefficient of $t^n$ in the latter expression is,
for $n\geq M$,
\[
A(n-M)\prod_{p+1\leq i\leq n}|\mu_i|^{a_i(s+\beta)}|\alpha|^{-s-\beta}
=\prod_{p+1\leq i\leq n}|\mu_i|^{a_i(s+\beta)}|\alpha|^{-s-\beta}
\sum_rn^{r-1}A_r(M),
\]
and, for $n<M$, the coefficient of $t^n$ is $0$. From this the
lemma follows.

\subsection*{Geometry of the Igusa expansion}

Let $E$ be an irreducible divisor in $Y''\subset G\times S''$,
which we know is regular in
codimension one, for which the image of $E$ under $\pi$ is contained in the
Zariski closure of a unipotent orbit ${\mathcal O}$ in $G$. Such a divisor is
called an \textbf{${\mathcal O}$-divisor} if ${\mathcal O}$ is minimal (with respect to
the partial ordering given by inclusion of Zariski closures) with this
property. It will be called an \textbf{${\mathcal O}(F)$-divisor} if $E$
is defined over $F$ and
$\pi(E(F))$ meets ${\mathcal O}(F)$ in a dense subset of $\ooo(F)$
as p-adic manifolds.
Let $E$ be an ${\mathcal O}(F)$-divisor.

First, we record an observation about
the fibration $\pi:E\rightarrow \overline{\ooo}$.
As a corollary to \cite[ch II, Prop. 3.1]{DG80}, we have

\begin{proclaim}{Lemma 5.10}
(a) The morphism $\pi:E\cap Z_S \rightarrow \overline{\ooo}$ is
smooth of relative
dimension $2r$, provided $\ooo$ is $r$-regular.

(b) $\pi(E\cap Z_S)\subset \ooo$.

(c) Let
\[
\{1\}\subset \overline{\mathcal O}_1={\mathcal O}_1\cup \{1\}
\subset ...\subset \overline{\mathcal O}={\mathcal O}\cup {\mathcal O}_m,
\]
where each ${\mathcal O}_i$ is a unipotent orbit of $G$,
and let $E_i=\pi^{-1}({\mathcal O}_i)\cap E$. Then $\pi|_{E}:E_i
\rightarrow {\mathcal O}_i$ is faithfully flat.
\end{proclaim}

We may apply Igusa's expansion to
\[
F(\lambda)=
\int_{\phi^{-1}(\lambda)(F)}f(y)|\omega_\lambda|,
\tag{5.11}
\]
where $f$ is any function on $Y_C(F)$ for which (a) $f$ is compactly
supported on the p-adic manifold $Y_C(F)$ and (b) $\phi$, $|\omega_\lambda|$,
$Y_C$, and $f$ satisfy the conditions of Igusa data in \cite{LS83}.
Since $f$ is compactly supported, so is $f|_{\ooo}$. On each fiber
$\tilde{E}(u)$ the pull-back $\pi^*(f)$ takes a constant value $f(u)$.
By \cite[Ch I]{Ha92a}, the regular function $\kappa(t_*)$ extends
from $E(u)$ to a regular function on $\tilde{E}(u)$. By the lemma
above, we have
\[
\int_{\overline{\ooo}(F)}f \cdot
PV\int_{\tilde{E}(u)(F)}\kappa(\tilde{t}_*)
|\omega_{\tilde{E}(u)}||\omega_{\ovo}|=
PV\int_{\tilde{E}(F)}\tilde{\pi}^*(f) \kappa(\tilde{t}_*)
|\omega_{\tilde{E}}|,
\tag{5.12}
\]
provided everything in sight satisfies the conditions (5.1-4).
Recall that, by definition, there are disjoint patches
$U_i\subset \tilde{Y}_1(F)$,
each analytically isomorphic to a finite box in affine space, such that
\[
PV\int_{\tilde{E}(F)}\tilde{\pi}^*(f) \kappa(\tilde{t}_*)
|\omega_{\tilde{E}}|=
\sum_iPV\int_{\tilde{E}(F)\cap U_i}\tilde{\pi}^*(f) \kappa(\tilde{t}_*)
|\omega_{\tilde{E}}|,
\tag{5.13}
\]
provided that the conditions (5.1-4) hold.
(By the way, the computations in \S 4 show that in all
the cases looked at in this paper,
$b_{{E}(F)}>0$. In particular, $b_{{E}(F)}\not= 0$.)

\begin{proclaim}{Lemma 5.14}
Let $G=G_2$.
Let $E$ be a base ${\mathcal O}(F)$-divisor of the
variety $Y_C=\tilde{Y}_1$. If $E$ is $F$-rational, $E(F)\not= \emptyset$,
then $e(\theta_E,\beta_E)=1$. In particular, spurious divisors
do not contribute to the Igusa expansion.
\end{proclaim}

\begin{proof} The proof is case-by-case using local coordinates in each case.
Assume $e(\theta_E,\beta_E)\geq 2$. By \cite{La83}, the
principal value integral contributing to
$F_1(\theta_E,\beta_E,f_\kappa)$ is
\[
\int_{\ooo(F)}f\int_{D_u(F)}\kappa(\tilde{t}_*)
|\omega_{D_u}||\omega_{\ooo}|,
\]
where $D=\cap_{i=1}^{e(\theta_E,\beta_E)}E_i$,
where $\theta_E=\theta_i$ and $\beta_E=\beta_i$,
and $D$ is $F$-rational (though the $E_i$ themselves need not be
$F$-rational).

Claim: The homogeneity degree of $|\omega_{D_u}|$, as a function
of the tangent direction $X$, is $dim(D)-dim(\ooo)<dim(E)-dim(\ooo)=2r$.

To prove this, we look at each divisor in the inverse image of
$\{\lambda =0\}$.

We may omit the regular and 1-regular cases since they were already
dealt with in \cite{Ha92a}, \cite{LS90}.

Let ${\mathcal O}={\mathcal O}_2$ denote a 2-regular orbit and let
$E=E_{\emptyset, \alpha,\delta}$.
First suppose $e(\theta_E,\beta_E)=2$.
From \S 7 below, the form on the $F$-rational fiber
$E_{\emptyset,\alpha,\delta}(u)$, $u\in {\mathcal O}$, is given
on $Y''_{\emptyset,\alpha,\delta}$ by
\[
\omega_{E_{\emptyset,\alpha,\delta}(u)}
=gdz(\beta)dx(\zeta)dv(\alpha)dv(\zeta),
\tag{5.15}
\]
where $g$ is a regular function.
% Claude: VARIANT. In g2new5d.tex the next parenthetical sentence reads:
%   (Though these coordinates may not be $F$-rational, we
%   may ``average'' them (or a suitable ${\Bbb Z}$-linear combination of them)
%   over the Galois group using the calculations of \S 2
%   and then we make a change
%   of variables in the form to get a form in $F$-rational coordinates.)
(Though these coordinates may not be $F$-rational, we
may ``average'' them over the image of the Galois group in $\Omega$
using the calculations of \S 2 and then we make a change
of variables in the form to get a form in $F$-rational coordinates
satisfying (5.1-4).)
By (1.22), this form is homogeneous
of degree 2 as a function of $X$. Let $E'$ denote another
$F$-rational ${\mathcal O}_2$-rational divisor with
$\theta_E=\theta_{E'}$, $\beta_E =\beta_{E'}$. By Hironaka's theorem,
$E$ intersects $E'$ normally in $Y_C$. The intersection $D=E\cap E'$
maps onto ${\mathcal O}$ under $\pi$ by \cite[chap II, Prop 3.1]{DG80}.
The variety $D$ must have codimension 1 in $E$ and will have
an Igusa form on
a fiber $D(u)$ which satisfies
\[
\omega_{D(u)}=g_Ddx dv_1dv_2,
\]
where $x$ is a coordinate on $Z$, $v_1,v_2$ are coordinates on the
$N_+$ components, and $g_D$ is regular. Therefore, $\omega_{D(u)}$
is homogeneous of degree $<2$ as a function of $X$.

Now suppose $e(\theta_E,\beta_E)=3$. The argument is similar,
hence omitted, in this case. The claim follows in the
2-regular case.

Let ${\mathcal O}={\mathcal O}_3$ denote a 3-regular orbit and let
$E=E_{\emptyset, \beta,\zeta}$.
First suppose $e(\theta_E,\beta_E)=2$.
From \S 8, the form on the $F$-rational fiber
$E_{\emptyset,\beta,\zeta}(u)$, $u\in {\mathcal O}$, is given
on $Y''_{\emptyset,\beta,\zeta}$ by
\[
\omega_{E_{\emptyset,\beta,\zeta}(u)}
=h dw(\epsilon)dw(\zeta)dx(\zeta)dv(\beta)dv(\epsilon)dv(\zeta),
\tag{5.16}
\]
where $h$ is a regular function.
(Again, these coordinates may be ``averaged to obtain $F$-rational
coordinates as in the 2-regular case.)
By (1.22), this form is homogeneous
of degree 3 as a function of $X$. Let $E'$ denote another
$F$-rational ${\mathcal O}_3$-rational divisor with
$\theta_E=\theta_{E'}$, $\beta_E =\beta_{E'}$. By Hironaka's theorem,
$E$ intersects $E'$ normally in $Y_C$. The intersection $D=E\cap E'$
maps onto ${\mathcal O}$ under $\pi$ by \cite[chap II, Prop 3.1]{DG80}.
The variety $D$ must have codimension 1 in $E$ and will have Igusa form on
a fiber $D(u)$ which satisfies
\[
\omega_{D(u)}=h_Ddx dy dv_1dv_2dv_3,
\]
where $x,y$ are coordinates on $Z$, $v_1,v_2,v_3$ are coordinates on the
$N_+$ components, and $h_D$ is regular. Therefore, $\omega_{D(u)}$
is homogeneous of degree $<3$ as a function of $X$.

Now suppose $e(\theta_E,\beta_E)=3$. The argument is similar,
hence omitted, in this case. The claim follows in the
3-regular case.
This proves the claim for all the orbits.

However, the claim
contradicts the degree of homogeneity of the $\ooo$-germ.
\end{proof}

\begin{proclaim}{Lemma 5.17} Let $G=G_2$.
Let $E$ be a base ${\mathcal O}(F)$-divisor of the
variety $Y_1$, $\tilde{E}$ its preimage under
$Y_C\rightarrow Y_1$.
In this case $F_1(\theta_E,\beta_E,f_\kappa)$ is given
(up to a constant) by
\[
PV\int_{\tilde{E}(F)}\tilde{\pi}^*(f_\kappa) |\omega_{\tilde{E}}|
=\int_{\ovo(F)_{st}}f
\int_{\tilde{E}_u(F)}\kappa(\tilde{t}_*)
|\omega_{\tilde{E}_u}||\omega_{\ovo}|,
\]
where $\omega_{\ovo_{st}}$
is a sum of forms of maximal degree on the
equidimensional pieces ${\mathcal O}_i$ in Lemma 5.10 of the stable
orbit $\ovo(F)_{st}=\cup_i{\mathcal O}_i(F)$.
\end{proclaim}

\begin{proof} Because of the previous lemma, (0.1), and formula 0.6,
it suffices to check that the conditions of \cite{LS83} are met. This
has essentially already been done in \cite{La83},
\cite[ch I, proof of Prop. 5.2]{Ha92a},
but we fill in some details regarding the function $\kappa(\tilde{t}_*)$
and the cocycle-valued function $t_*$.

Some information about $t_*$ is given following formula 0.6 in \S 0.
More explicit information about
$t_\sigma :Y_1(B_+,B_-)\rightarrow T_G(\overline{F})$, where
$Y_1(B_+,B_-)$ is defined in \S 1, is given in
\cite[ch I, Corollary 5.4]{Ha92a}. Since the
only $F$-form of $G_2$ over a p-adic field is the split form \cite{Ti65},
we have $G=G_2=G_{in}=G_{qs}$, $T_G=T_{in}, \ \underline{T}=T_{qs}$, and
$\underline{B}=B_{qs}$, the notation being as in \cite{Ha92a}. Let
$\alpha\in \Delta$ denote a simple root,
If $\sigma=\sigma_\alpha\in \Gamma$ acts on
$T_G(\overline{F})$ by conjugation by the
reflection $s_\alpha\in \Omega$, then
we choose $g$, $h$ as in
\cite[ch I, \S 5]{Ha92a}. Taking
$A_\sigma=w_\sigma^{-1}=s_\alpha \in N_G(T_G)$.
in \cite[ch I, Corollary 5.7]{Ha92a}, we find there is a ``constant''
$t_1\in T_G(\overline{F})$ independent of the star such that at a point
\[
(t^g,((B(W))_{W\in{\mathcal W}},(z(\alpha))_{\alpha\in\Delta},n)
\in\chi^{-1}(C)\times S_1(B_+,B_-)\times N_+,
\]
the value $t_\sigma$ is determined by
\[
s_\alpha \epsilon_{-\alpha}(-\xi)n^{-1}w_\sigma^{-1}s_\alpha^{-1}\in
N_+^- N_+ (t_1t_\sigma)^{w_\sigma^{-1}s_\alpha^{-1}}.
\tag{5.18}
\]
Take the transpose of both sides of Equation 5.5 in
\cite[ch I]{Ha92a}. The resulting equation,
along with (5.10) implies that
\[
t_1t_\sigma = \alpha^\vee(\xi),
\]
where $\alpha^{\vee}$ is its coroot, and $\xi$ satisfies (5.2)
as a function of the coordinates of $Y''(B_+,B_-)\times N_+$. Choose
$B_-, B_+$ such that these coordinates are $F$-rational.
From the appendix, we know that there is a natural bijection
\[
H^1(\Gamma,T_G)\rightarrow \zzz/ m\zzz,
\]
where $m\in \{2,3\}$. Since each such $H^1(\Gamma,T_G)$ is cyclic,
$\kappa(\tilde{t}_*)$ is determined by its value on a generator. This,
in turn, determines a character of $\zzz/m\zzz$ and
hence by local class field theory a
character $\theta$ of $F^\times/NE_{T_G}^\times$ (where $E_{T_G}$
denotes the Galois closure of the splitting field of $T_G$. One may
construct a (non-canonical) map
$\phi:H^1(\Gamma,T_G) \rightarrow F^\times$
such that $\kappa(\tilde{t}_*)=\theta(\phi(\tilde{t}_*))$
is a monomial in the coordinates of $Y''(B_+,B_-)\times N_+$.
This allows us to conclude that the
conditions of (5.1-4) have been met.
\end{proof}

We close this section by recalling an important reduction lemma
due to T. Hales. First, observe that for each $r$-regular
orbit ${\mathcal O}$, $r>0$,
every function $f\in C_c^\infty(G(F))$ may be written in the form $f=f_1+f_2$
with $f_1$ supported off the Zariski closure of ${\mathcal O}$ in $G$
and $f_2$ whose support does not meet the orbits outside the
Zariski closure of ${\mathcal O}$ in $G$.

\begin{proclaim}{Reduction Lemma 5.19 (Hales)} Suppose that for each $(T,\kappa)$
associated to $H$
the germs $\Gamma_{{\mathcal O}}^{(T,\kappa)}(\gamma)$ have the form
\[
\Gamma_{{\mathcal O}}^{(T,\kappa)}(exp(\lambda^2 X))=
|\lambda|^{2r}m_{{\mathcal O}}^{(T,\kappa)}(X),
\]
for $\lambda$ small, where $m_{{\mathcal O}}^{(T,\kappa)}(X)$ satisfies
\[
m_{{\mathcal O}}^{(T,\kappa)}(X)=m_{{\mathcal O}}^{(T,\kappa)}(X+Y),
\]
for $X,Y\in {\mathfrak t}_G$ sufficiently small, $X$ regular,
$X+Y$ regular, and $Y$ belonging to the Lie algebra of the center
of an endoscopic group $H\subset G$ satisfying $Cent(Y,G)=H$.
Then $(G,H)$ admits local
$\Delta$-transfer at the identity for all functions $f\in C_c^\infty(G(F))$
whose support does not meet the orbits outside the
Zariski closure of ${\mathcal O}$ in $G$.
\end{proclaim}

\begin{proof} Follows \cite[Lemma 6.1]{Ha92b}. (See also \cite[Proposition
2.11]{Ha89}.)
\end{proof}

This lemma implicitly uses \cite{LS90}.
% ======================================================================
% Claude: Section 6, source file g2new6.tex (\pageno 43). Here \off = \overline{F}.
% ======================================================================
\section*{\S\,6 The matching for the $1$-regular germ}

Let $G=G_2$ throughout this section.

The $\kappa$-stable subregular germ for $G_2$ should match with the
stable subregular germ of $PGL(3)$ or $(SL(2)\times SL(2))/\{\pm 1\}$,
for homogeniety reasons. The idea to verify this is simply to write down
the Igusa expansions of the two orbital integrals and match the $\beta =2$
terms. Furthermore, for the 1-regular comparisons, we can and do
assume that $f$ is supported off the 0-regular orbit.

Let $u=exp(X_\alpha+X_{\alpha+\beta})$. Then the subregular orbit
is given by ${\mathcal O}_1=u^{G_2}$. This is not a stable orbit,
but it is Richardson.
See the appendix for more on its $F$-classes.
It will turn out that both the biquadratic and the cubic Cartans
can have non-zero 1-regular $\kappa$-stable germs.

From \S\S 4,5, there are three base divisors with $\beta =2$,
and from \cite[ch. VI, \S 2]{Ha92a}, we know that all fundamental divisors
of $G_2$ have $\theta =1$.
They are described as follows:

$E_\alpha$ is given locally as $\{\lambda =0\}$
(with multiplicity one) on $Z_\alpha$
and $Z_\alpha$ is regular with generators of its local coordinate ring
being $\lambda$, $x(\alpha)$, $z(\beta)$, $x(\gamma)$,
$x(\delta)$, $x(\epsilon)$, $x(\zeta)$.

$E_\beta$ is given locally as $\{\lambda =0\}$ (with multiplicity one)
on $Z_\beta$
and $Z_\beta$ is regular with generators of its local coordinate ring
being $\lambda$, $z(\alpha)$, $x(\beta)$, $x(\gamma)$,
$x(\delta)$, $x(\epsilon)$, $x(\zeta)$.

$E_{\emptyset,\alpha,\gamma}$ is given locally
by $\{z(\alpha) =0\}$ (with multiplicity one)
on $Z_{\emptyset,\alpha,\gamma}$ and $Z_{\emptyset,\alpha,\gamma}$
is regular with generators of its local coordinate ring
being $z(\alpha)$, $z(\beta)$, $x(\gamma)$, $w(\gamma)$,
$x(\delta)$, $x(\epsilon)$, $x(\zeta)$. This divisor may be eliminated
from consideration by \cite[Ch V, Lemma 1.1]{Ha92a}.

As a corollary to \cite[ch 6, Lemma 1.1]{Ha92a}, we know that
$E_\alpha$ and $E_\beta$ are $F$-rational. Note that the action of the Galois
group $\Gamma$ on $T_G$ and on $Z$ factors through a finite group
$\Gamma_{T}\subset \Omega$.
To satisfy the hypotheses (5.1-4), one may use
the calculations of \S 2 to ``average'' the coordinates of $Z_\alpha$ or
$Z_\beta$ over the Galois group $\Gamma_T$
to get $F$-rational coordinates on these divisors.

Moreover, if we choose the Cartan $T_G$ such that $\Gamma$ acts
on $T_G$ via a product of two cyclic subgroups of order two of the Weyl group
$\Omega$ then the $\kappa$-stable subregular germ
of $G_2$ should match with the stable subregular germ of
$(SL(2)\times SL(2))/\{\pm 1\}$
(and by the appendix, $\kappa$ is quadratic)
if we choose the Cartan $T_G$ such that $\Gamma$ acts on $T_G$
via a cyclic subgroup of $\Omega$ of order three
then the $\kappa$-stable subregular germ
of $G_2$ should match with the stable subregular germ of
$PGL(3)$ (and by the appendix, $\kappa$ is a cubic character).

\begin{proclaim}{Theorem 6.1} Suppose $T_G$ is a non-split Cartan subgroup
of $G_2$ which splits over a
Galois extension for which the Galois group action on $T_G$
is either a direct product of two cyclic subgroups of $\Omega$ order $2$
or a cyclic subgroup of order $3$.
Assume that $f\in C_c^\infty(G(F))$ satisfies $supp(f)\cap {\mathcal O}_1
\not= \emptyset$, where ${\mathcal O}_1$ denotes a fixed $1$-regular
$G(F)$-orbit, and the support of $f$ does not meet any unipotent orbit
of $G(F)$ which is not $1$-regular. Then there is an $f^H$, smooth on
$H(F)$, for which, for all regular $\gamma_G\in T_G(F)$
in a sufficiently small neighborhood of $1$, we have
\[
\Delta_{G/H}^{T_G,\kappa}(\gamma_H,\gamma_G)
\Phi^{T_G,\kappa}(\gamma_G,f)=
\Phi^{T_H,st}(\gamma_H,f^H),
\]
where $\kappa$, $T_G$ satisfy
$H_{T_G,\kappa}=H_i$. Here $i=1$, so $H=PGL(3)$, if
the Galois group action on $T_G$
is a cyclic subgroup of order $3$, and
$i=2$, so $H=(SL(2)\times SL(2))/\{\pm 1\}$, if
the Galois group action on $T_G$ is
a direct product of two cyclic subgroups of $\Omega$ order $2$.
\end{proclaim}

This will be proven below.

\subsection*{Proof of Theorem 6.1}

By the work of Langlands and Shelstad, \cite{LS87}, \cite{LS90}, we need only
prove the corresponding matching of germs conjecture stated in
\S 0 above.

We discuss the Igusa expansion of the $G_2$ side more carefully:
Since only the $\theta=1$ term of formula 0.6(a)
contributes by \cite[ch 6, \S 2]{Ha92a},
the $\beta=2$ term of the Igusa expansion takes the form,
$|\lambda|\kappa(inv(\gamma_G^0,\gamma_{G,\lambda}))
D_H(\gamma_{H,\lambda})^{-1}$
times
\[
F_1(1,2,f_\kappa)=
PV\int_{U\cap E(F)}A_1(M)h_1|\nu_1|=
PV\int_{E(F)\cap U}\pi^*(f)\kappa_1(t_*)|\omega_E|
\tag{6.2}
\]
where $E$ is $E_\alpha$ or $E_\beta$, depending on $T_G$, and
where $A_1(M)$, $h_1$, $\nu_1$, ..., are as in \cite{La83},
and $f_\kappa$ includes the $\kappa$-factor. The factor
$m_\kappa(e)/\kappa_1(\lambda)|_{E(u)}$ is given by
\cite[Ch VI, \S 2 and Proposition 3.2]{Ha92a}, provided $E$ is the divisor
associated to a line of the Dynkin curve associated to the long root
of $G_2$ which is fixed by the Galois action.
By the homogeniety of germs, formula 0.6 and lemma 0.8,
this factor is equal to a multiple of
\[
\sum_{{\mathcal O}\subset {\mathcal O}_1^{st}}
\mu_{\mathcal O} (f)\Gamma^{T_G,\kappa}_{\mathcal O}(\gamma_\lambda),
\]
where ${\mathcal O}_1^{st}$ denotes the stable 1-regular unipotent orbit.
M. Assem has computed the number of $F$-rational classes in the
stable subregular class when $p\not= 2,3$ (see the appendix below).

We discussed the geometry of the fibering $E_\alpha\rightarrow
\overline{{\mathcal O}}_1$, where $\overline{{\mathcal O}}_1$ denotes the
Zariski closure, in \S 5.
We assume that the function $f$ meets ${\mathcal O}_1^{st}$ and no other
unipotent orbit.
By formula 0.6 and (6.2), we have
\[
PV\int_{U\cap E(F)}A_1(M)h_1|\nu_1|
=\int_{\pi(U)\cap {\mathcal O}_1(F)}f\int_{E_u(F)}
\kappa_1(t_*)|\omega_{E_\alpha(u)}||\omega_{{\mathcal O}_1}|,
\tag{6.3}
\]
where $E$ denotes either $E_\alpha$ or $E_\beta$.
We may restrict ourselves to
the case where $T_G$ is either biquadratic or cubic
as in Theorem 6.1. The dependence of (6.3) on
% Claude: In g2new6.tex \off is \overline{F}, so the next formula prints
% exp(\overline{F} X). Elsewhere (Sections 3-5, 7-9) \off is O_F; the
% intended meaning here may have been O_F.
$C\subset exp(\ovF\cdot X)$, where $X\in {\mathfrak t}$ is a tangent vector
for $C$, is via the dependence of $\kappa(t_*)$ and $\omega_{E_\alpha(u)}$
on $X$, where ${\mathfrak t}$ denotes the Lie algebra of $T_G$.

Recall \cite[Lemma 2.12]{La83} that the measure $\omega_Y$ already incorporates
the Weyl discriminant $\prod_{\sigma >0}(1-\sigma(\gamma)^{-1})$.
Also, from \S 4, recall
\[
|\omega_{Z_\alpha}|=
\frac{\lambda}{z(\beta)^2}d\lambda dx(\alpha) dz(\beta)dx(\gamma)
dx(\delta)dx(\epsilon)dx(\zeta),
\]
and
\[
|\omega_{Z_\beta}|=
\frac{\lambda}{z(\alpha)^2}d\lambda dz(\alpha) dx(\beta)dx(\gamma)
dx(\delta)dx(\epsilon)dx(\zeta).
\]
Note that these coordinates are not all $F$-rational if $c_\sigma$ is a
simple reflection by (2.16), (2.17).

On ${\mathcal O}_1$ the roots ``outside $P_\alpha$'' (or the
roots ``outside $P_\beta$'') may be used to
parameterize the orbit - see the lemma below.

The digression below is motivated by \cite[Remark V.2.2]{Ha92a}.

Let ${\mathcal O}$ be the Richardson orbit associated to a parabolic $P$
as in \cite[Theorem 2, \S 3.9]{St74}. Write $P=MN_P$, so that $N_P^G$
contains ${\mathcal O}$ as a dense subset, ${\mathcal O}\cap N_P$ is dense in
$N_P$ and ${\mathcal O}\cap N_P$ forms a single class under $P$. We have
dim${\mathcal O}$ = dim$G$ - dim$M$. Let $u\in {\mathcal O}\cap N_P$ and recall
$N_P^- P$ is dense in $G$ by the Bruhat decomposition. The $P$-orbit
$\{u^p\ |\ p\in P\}$ is equal to ${\mathcal O}\cap N_P$ and is dense in $N_P$.
Therefore, the coordinates of $N_P$ yield coordinates for the $P$-orbit
of $u\in {\mathcal O}\cap N_P$. The $G$-orbit of this $u$, that is
${\mathcal O}$ itself, contains the $N_P^- P$-orbit of $u$ as an open
subvariety. We may regard the $N_P^- P$-orbit of $u$ as the
$N_P^-$-orbit of ${\mathcal O}\cap N_P$. By Corollary 1.6, for each
$\overline{n}\in N_P^-$, $B^{\overline{n}}$ is opposite to
$B^-$ and $B^{\overline{n}}\not= B$ unless $\overline{n}=1$.
Therefore, for our $u$ as above, $u^{\overline{n}}\notin {\mathcal O}\cap N_P$.
From this it follows that the $N_P^-$-orbit of ${\mathcal O}\cap N_P$
is isomorphic to $N_P^-$.

\begin{proclaim}{Lemma 6.4 (Hales)}
If ${\mathcal O}$ is the Richardson orbit associated to a
parabolic $P$ then the coordinates on $N_P^- \times N_P$ serve as local
coordinates on ${\mathcal O}$.
\end{proclaim}

Returning to the proof of Theorem 6.1,
associated to the roots outside $P_\alpha$, we have the invariant measure
\[
\omega_{{\mathcal O}_1}=
\lim_{\lambda\rightarrow 0}
\frac{dx(\beta)...dx(\zeta)dv(\beta)...dv(\zeta)}{\lambda},
\]
Recall that the invariant measure $\omega_{{\mathcal O}_1}$ is
unique up to a constant \cite{Ra72}. On the fiber $E_\beta(u)$, we use
\[
\omega_{E_\beta(u)}=\frac{1}{z(\alpha)^2}dv(\alpha)dz(\alpha),
\]
by
\S 4 above, \cite[Lemma 2.12]{La83}. We have then
\[
\omega_{E_\beta(u)}=
\frac{(1-\alpha (\gamma_\lambda)^{-1})}{\lambda}\Big|_{\lambda=0}
dv(\alpha)\frac{dz_\alpha}{z_\alpha^2}
={\alpha(X)}dv(\alpha)\frac{dz_\alpha}{z_\alpha^2},
\]
where $\alpha(X):=\frac{1-\alpha^{-1}(\gamma_\lambda)}{\lambda}|_{\lambda=0}$
denotes the inner product of $\alpha$ with the tangent vector $X$ of $C$.
Similarly, associated to the roots outside $P_\beta$, we have
\[
\omega_{E_\alpha(u)}=\frac{1}{z(\beta)^2}dv(\beta)dz(\beta)
={\beta(X)}dv(\beta)\frac{dz_\beta}{z_\beta^2}.
\]
Therefore, the $\beta =2$ term is equal to
\[
\begin{split}
&|\lambda|\kappa(inv(\gamma_G^0,\gamma_{G,\lambda}))
D_H(\gamma_{H,\lambda})^{-1}
PV\int_{U\cap (E_\alpha)(F)}A_1(M)h_1|\nu_1|\\
&=
|\lambda|\kappa(inv(\gamma_G^0,\gamma_{G,\lambda}))
D_H(\gamma_{H,\lambda})^{-1}
|\beta(X)|\int_{{\mathcal O}_1}f\int_{E_\alpha(u)\cap (U\cap \pi^{-1}(u))}
\kappa(t_*)|dv'(\beta)
\frac{dz_\beta}{z_\beta^2}||\omega_{{\mathcal O}_1}|,
\end{split}
\tag{6.5a}
\]
or
\[
\begin{split}
&|\lambda|\kappa(inv(\gamma_G^0,\gamma_{G,\lambda}))
D_H(\gamma_{H,\lambda})^{-1}
PV\int_{U\cap (E_\beta)(F)}A_1(M)h_1|\nu_1|\\
&=|\alpha(X)|\int_{{\mathcal O}_1}f\int_{E_\beta(u)\cap (U\cap \pi^{-1}(u))}
\kappa(t_*)|dv'(\alpha)
\frac{dz_\alpha}{z_\alpha^2}||\omega_{{\mathcal O}_1}|,
\end{split}
\tag{6.5b}
\]
From these expressions and the remarks in \S 0 above,
one may verify the local translation invariance
property of Hales' lemma 5.19.
From this, the transfer follows in the $1$-regular case,
for compactly supported functions whose support meets the 1-regular
orbit and no other orbit.

% Claude: UNTYPESET. The following text follows \enddocument in g2new6.tex
% and was never typeset. It contains the source's tags (6.6) and (6.8)
% (there is no (6.7)) and a "Lemma 9.1" that does not belong to Section 9.
\iffalse
For the sake of comparison, we look at
the corresponding expression for the subregular
term for $PGL(3)$, $(SL(2)\times SL(2))/\{\pm 1\}$.

The stable subregular germ of $(SL(2)\times SL(2))/\{\pm 1\}$:
The stable subregular
unipotent orbit of $(SL(2)\times SL(2))/\{\pm 1\}$ is the union of
${\mathcal O}_{tr}\times {\mathcal O}_{reg}^{st}$ and
${\mathcal O}_{reg}\times {\mathcal O}_{tr}$. The stable regular germ on $SL(2)$
may be identified with the regular germ of $PGL(2)$ (they differ
at most by some non-zero constant of normalization -
here one may use another result of Repka or the fact that
$\sum_{{\mathcal O}\subset {\mathcal O}_H^{st}}C_{\mathcal O}(\gamma_\lambda)$
is a constant, in the notation of \cite[Lemma 2.4]{SS83},
where $p\not= 2$ is used) and
can be found, for example, in \cite{SS83}; to get the stable
subregular germ of $(SL(2)\times SL(2))/\{\pm 1\}$ we take the product of
the regular germ of $PGL(2)$ with the stable trivial germ of $SL(2)$.
On each fixed
elliptic Cartan of $PGL(2)$ the regular germ is a non-zero constant
times $|D(\gamma_\lambda)|^{-1/2}$, where $D$ is the usual Jacobian.
Denoting the simple positive roots on $(SL(2)\times SL(2))/\{\pm 1\}$
by $\alpha_1$, $\alpha_2$, we find that
\[
|D_{H_0}(\gamma_{\lambda,H})|^{1/2}=
|\lambda||\alpha_1(X_{H_0})|, H_0=SL(2),
\tag{6.6}
\]
for $\gamma_{\lambda,H}=exp(\lambda X_H)\in T^H(F)\rightarrow
SL(2,F)\times SL(2,F)$ elliptic regular. We want this function of $X_H$
to be equal, up to a constant factor, to the function of $X$ occuring
in $c_{\kappa, T,C}$. Recall that $\alpha_1$ corresponds to $\alpha$ and
$\alpha_2$ corresponds to $3\alpha +2\beta$.

The stable subregular germ of $PGL(3)$: The stable subregular
unipotent orbit of
$PGL(3)$ consists of only one $F$-class. The subregular germ on $GL(3)$
can be found in \cite{Re84}. On  regular elements of an
elliptic Cartan $E^{\times}$ of $GL(3)$, it is a non-zero constant
times $|\lambda (3\alpha +2\beta)(X)||D_H(\gamma_\lambda)|^{-1/2}$
(at least for $p\not= 3$). The matching of the Igusa transfer factor
$p_{{\mathcal O},T}$ with the LS transfer factor predicts something slightly
different: a non-zero constant times
$|\lambda \beta(X)||D_H(\gamma_\lambda)|^{-1/2}$. However, these are
actually the same quantity, the discrepancy being accounted for by
the fact that for such elliptic Cartans of $PGL(3)$,
$|1-\beta(\gamma_G)|=|1-(3\alpha +2\beta)(\gamma_G)|$
for all regular $\gamma_H\in T_H(F)$.
(I have also used that
$dist(t,F)=|t-\frac{1}{3}tr(t)|_E$, in the notation of \cite{Re84},
which may not be valid if $p=3$. The modification needed for
$p=3$ is minor and does not affect the final form of our results.)

Under the isomorphism between $T_G$ and $T_H$ (see \S A.3 of the appendix),
on the $H=(SL(2)\times SL(2))/\{\pm 1\},PGL(3)$ side,
the subregular term has the form, up to a factor of $|\lambda|^{-1}$,
\[
\begin{cases}
(|(3\alpha+2\beta)(X)|)^{-1}d_{T_H}\cdot
\mu_{{\mathcal O}_1(F)}(f^H),&\text{ for }
H=A_1\times A_1,\\
|(3\alpha+\beta)(X)(3\alpha+2\beta)(X)|^{-1}d'_{T_H}
\cdot \mu'_{{\mathcal O}_1(F)}(f^H),&\text{ for }H=A_2,
\end{cases}
\tag{6.8}
\]
provided $T$ is either biquadratic or cubic (respectively) and
$X$ is regular.
Observe that, as a function of $X$, the $G_2$-side differs
from the $H$-side by $c_{\kappa, T_G,\alpha(X)}$ times
$D_H(\gamma_{\lambda,H})$,
in the cubic case, and by $c_{\kappa, T_G,\beta(X)}$
times $D_H(\gamma_{\lambda,H})$,
in the biquadratic case.

\begin{proclaim}{Lemma 9.1} Let $s:=\sigma_\alpha$ and
$r:=\sigma_\alpha\sigma_\beta$.
On $E_\alpha$, $E_{\beta}$, the generators
of the (respective) local coordinates rings on
$Z_\alpha$, $Z_\beta$ are all $\sigma$-invariant, where
$\sigma=s$ or $\sigma=r^2$.
The local coordinate $x(\delta)$, restricted to either of these
divisors, is not $\sigma$-invariant if
\[
\sigma=\sigma_\beta, \sigma_\alpha\sigma_\beta,
\sigma_\beta\sigma_\alpha, \sigma_\beta\sigma_\alpha\sigma_\beta .
\]
(This last statement does not appear to be consistent
with \cite[Lemma 1.1, \S 1]{Ha92a} but is not needed anyway.)
\end{proclaim}
\fi
% ======================================================================
% Claude: Section 7, source file g2new7.tex (\pageno 47). Here \off = {\Cal O}_F
% (not used in this section).
% ======================================================================
\section*{\S\,7 On the 2-regular germ for $G_2$}

Let $G=G_2$ throughtout this section.

The only essential base divisor which can contribute to the 2-regular
germ for $G_2$ via Igusa theory is $E_{\emptyset,\alpha,\delta}
$,
which has $\beta(E_{\emptyset,\alpha,\delta})=3$.
It must be $F$-rational since for every $\sigma\in\Gamma$
$\sigma(E_{\emptyset,\alpha,\delta})$
is also a 2-regular base divisor (see observation (I) in \S 4),
so $\sigma(E_{\emptyset,\alpha,\delta})=E_{\emptyset,\alpha,\delta}$.
Recall that the action of the Galois
group $\Gamma$ on $T_G$ and on $Z$ factors through a finite group
$\Gamma_T$, say.
To satisfy the hypotheses (5.1-4), one may use
the calculations of \S 2 to ``average'' the coordinates of
$Z_{\emptyset,\alpha,\delta}$
over the Galois group $\Gamma_T$
to get $F$-rational coordinates on this divisor if needed.

The branch $E_{\emptyset,\alpha,\delta}$ is defined by
$\lambda =0$ on the regular subset $Z_{\emptyset,\alpha,\delta}$:
$\{x(\delta)\not= 0,w(\delta)\not= 0,z(\beta)\not= 0\}\subset Z$.
The patch $Z_{\emptyset,\alpha,\delta}$ is regular with generators
$z(\alpha)$, $z(\beta)$, $w(\gamma)$, $w(\delta)$, $x(\delta)$,
$x(\epsilon)$, and $x(\zeta)$. The divisor $\lambda =0$ occurs
on this patch as $z(\alpha)=0$ with multiplicity 2.
The local coordinate relations yield that on $E_{\emptyset,\alpha,\delta}$,
we have $\lambda=x(\alpha)=x(\beta)=z(\alpha)=x(\gamma)=w(\eps)=w(\zeta)=0$.

By \cite[ch 6, \S 2]{Ha92a}, only the term $\theta=1$
contributes to the sum over $\theta$ in formula 0.6(a).
The contribution of the 2-regular germ to the $\kappa$-orbital
integral is then
$\kappa(inv(\gamma_G^0,\gamma_{G,\lambda}))
D_H(\gamma_{H,\lambda})^{-1}$ times
\[
|\lambda|^2
PV\int_{E_{\emptyset,\alpha,\delta}(F)}\pi^*(f)
\kappa_1(t_*)
|\omega_{E_{\emptyset,\alpha,\delta}}|
=|\lambda|^2
PV\int_{{\mathcal O}_2}f\int_{E_{\emptyset,\alpha,\delta}(u)(F)}
\kappa_1(t_*)
|\omega_{E_{\emptyset,\alpha,\delta}(u)}|
.
\tag{7.1}
\]

\begin{proclaim}{Theorem 7.2} Suppose $T_G$ is a non-split Cartan subgroup
of $G_2$ which splits over a
Galois extension for which the Galois group action on $T_G$
is a direct product of two cyclic groups of order $2$.
Assume that $f\in C_c^\infty(G(F))$ satisfies $supp(f)\cap {\mathcal O}_2
\not= \emptyset$, where ${\mathcal O}_2$ is a fixed $2$-regular
$G(F)$-orbit,
and the support of $f$ does not meet unipotent orbit
of $G(F)$ which is not $2$-regular. Then there is an $f^H$, smooth on
$H(F)$, for which, for all regular $\gamma_G\in T_G(F)$
in a sufficiently small neighborhood of $1$, we have
\[
\Delta_{G/H}^{T_G,\kappa}(\gamma_G)
\Phi^{T_G,\kappa}(\gamma_G,f)=
\Phi^{T_H,st}(\gamma_H,f^H),
\]
where $\kappa$, $T_G$ satisfy
$H_{T_G,\kappa}=(SL(2)\times SL(2))/\{\pm 1\}$.
(In particular, $\kappa$ is quadratic and $T_G$ is biquadratic.)
\end{proclaim}

This will be proven below.

\subsection*{Proof of Theorem 7.2}

By \cite{Sh72} or \cite{Ro80} the germ on the
right-hand side is a non-zero constant. The theorem will follow provided
we can show that both sides are equal, as a function of the tangent
direction $X$, up to a constant factor. Indeed, such a constant
factor can be absorbed into
$f^H$ provided we choose $f$ in the beginning so that its support meets the
2-regular orbit and no other unipotent orbit. We fix such a choice
for the remainder of this section.

We explicitly construct
local coordinates for the unipotent orbit ${\mathcal O}_2^{st}$.
Let $u=\theta_{\alpha+\beta }(1)\in G_2$, and compute
the orbit $ad({\mathfrak g}_2)(X_{\alpha+\beta })$.
From this, we conclude that $u^{G_2}$
is 8-dimensional, so ${\mathcal O}_2^{st}=u^{G_2}$.
Analogously to Hales' lemma 6.4,
we find that the coordinates of the subvariety
$\prod_{\sigma\in \Sigma_G-\Sigma_{H_2}}U_\sigma$
of $G_2$ yields local coordinates for the $2$-regular orbit.
Here $\Sigma_{H_2}^+=\{\alpha,\zeta\}$.
There is only one $F$-orbit of $u$ contained in the stable orbit of $u$
(see the appendix), so the $2$-regular orbit of $G_2$ is stable.
Local coordinates for the fiber $E_{\emptyset,\alpha,\delta}(u)$
may be choosen from the remaining coordinates on
$E_{\emptyset,\alpha,\delta}$:  $z(\beta)$, $x(\zeta)$,
$v(\alpha)$, and $v(\zeta)$.

We use the form on ${\mathcal O}_2$,
\[
\begin{split}
\omega_{{\mathcal O}_2}
&:=\frac{|\bigwedge_{\sigma\in \Sigma_{{\mathcal O}_2}}dx(\sigma)\wedge dv(\sigma)|}{|\lambda|^{3/2}}\\
&=\frac{|x(\delta)|^{1/2}}{|w(\delta)|^{3/2}|z(\beta)|^{3/2}}
dz(\alpha)dw(\gamma)dw(\delta)dv(\alpha)dv(\gamma)dv(\delta),
\end{split}
\tag{7.3}
\]
where $\Sigma_{{\mathcal O}_3}=\Sigma_{G}^+ -\Sigma_{H_2}^+$.
On $E_{\emptyset,\alpha,\delta}(u)$,
\[
\omega_{E_{\emptyset,\alpha,\delta}(u)}:=
\frac{{|x(\delta)|^{5/2}}|w(\delta)|^{7/2}}{{|z(\beta)|^{1/2}}}
dz(\beta)dx(\zeta)dv(\alpha)dv(\zeta).
\tag{7.4}
\]
We need to rewrite this last form as an expression
\[
\omega_{E_{\emptyset,\alpha,\delta}(u)}=
|p_{{\mathcal O}_2,T_G}(X)|\omega'_{E_{\emptyset,\alpha,\delta}(u)},
\]
where $p_{{\mathcal O}_2,T_G}(X)$ is a homogeneous polynomial in the coordinates
of tangent vector $X$ (expressed in a basis of the Lie algebra
${\mathfrak t}$ of $T_G$) and $\omega'_{E_{\emptyset,\alpha,\delta}(u)}$
is independent of the tangent direction.
Using (1.22), we find that $p_{{\mathcal O}_2,T_G}(X)$ is a product of the form
\[
p_{{\mathcal O}_2,T_G}(X)=\prod_{\sigma \in \Sigma_{H_2}^+}\sigma(X)^{f_\sigma},
\tag{7.5}
\]
where (by homogeneity of germs)
$\sum_{\sigma >0}f_\sigma =2$. To verify the hypotheses to Hales'
lemma 5.19 one must verify the translation invariance property
of (7.1). The translation invariance of this factor (7.5)
follows from the above expression for
$|p_{{\mathcal O}_2,T_G}(X)|$ as a product over roots in $H_2$.

By continuity, $\kappa(t_*)$ satisfies the
local translation invariance property as well for sufficiently small $Y$.
Therefore, the hypotheses to Hales' reduction lemma hold and the
theorem follows in this case from lemma 5.19.

% ======================================================================
% Claude: Section 8, source file g2new8.tex (\pageno 49). Here \off = {\Cal O}_F
% (not used in this section).
% ======================================================================
\section*{\S\,8 On the 3-regular germ for $G_2$}

Let $G=G_2$ throughtout this section.

The only essential
base divisor which can contribute to the 3-regular
germ for $G_2$ via Igusa theory is $E_{\emptyset,\beta,\zeta}$,
which has $\beta(E_{\emptyset,\beta,\zeta})=4$.
It must be $F$-rational since $\sigma(E_{\emptyset,\beta,\zeta})$
is also a $3$-regular base divisor, by the observation (I) of \S 4
and the case $G=G_2$ of conjecture 3.4 (which we proved in \S 3
above).
Recall that the action of the Galois
group $\Gamma$ on $T_G$ and on $Z$ factors through a finite group
$\Gamma_T$.
To satisfy the hypotheses (5.1-4), one may use
the calculations of \S 2 to ``average'' the coordinates of
$Z_{\emptyset,\beta,\zeta}$
(or a suitable ${\mathbb Z}$-linear combination of them)
over $\Gamma_T$ to get $F$-rational coordinates on this divisor.

The branch $E_{\emptyset,\beta,\zeta}$ is defined by
$\lambda =0$ on the regular subset $Z_{\emptyset,\beta,\zeta}$:
$\{x(\zeta)\not= 0,w(\zeta)\not= 0,z(\alpha)\not= 0\}\subset Z$.
If we use the local coordinates
$z(\alpha)$, $z(\beta)$, $w(\gamma)$, $w(\delta)$,
$w(\eps)$, $w(\zeta)$, and $x(\zeta)$ on
$Z_{\emptyset,\beta,\zeta}$ to compute the form
$\omega_{E}$ on $E=E_{\emptyset,\beta,\zeta}$, we get
\[
\begin{split}
\omega_{E_{\emptyset,\beta,\zeta}}
&=\frac{\omega_{{Z}_1}}{\lambda^\beta \frac{dz(\beta)}{z(\beta)}}\\
&=\frac{{d\lambda \wedge dx(\alpha) \wedge ...\wedge dx(\zeta)}}{\lambda^4 \frac{dz(\beta)}{z(\beta)}}\\
&= \frac{{x(\zeta)^2w(\zeta)^4}}{w(\delta)^7z(\alpha)^2}
dz(\alpha)\wedge dw(\gamma)\wedge ...\wedge dw(\zeta)\wedge dx(\zeta).
\end{split}
\tag{8.1}
\]
The divisor $\lambda =0$ occurs on
$Z_{\emptyset,\beta,\zeta}$ as $z(\beta)=0$ with multiplicity 2.
On $E_{\emptyset,\beta,\zeta}$, $\lambda =z(\beta)=x(\epsilon)=x(\delta)=
x(\gamma)=x(\beta)=x(\alpha)=0$.

By \S\S 4-5 and formula 0.6, the contribution of the
3-regular germ to the $\kappa$-orbital
integral is therefore,
$\kappa(inv(\gamma_G^0,\gamma_{G,\lambda}))D_H(\gamma_{H,\lambda})^{-1}$ times
\[
|\lambda|^3\sum_{\theta}\theta(\lambda)
F_1(\theta,4,\pi^*(f)\kappa_\theta(t_*))=
|\lambda|^3\sum_{\theta}\theta(\lambda)
\int_{E_{\emptyset,\beta,\zeta}(F)}\pi^*(f)\kappa_\theta(t_*)
|\omega_{E_{\emptyset,\beta,\zeta}}|.
\]
Note that only $\theta =1$ contributes to the sum by \cite[ch 6, \S 2]{Ha92a}.
Recall that, on $Z_{\emptyset,\beta,\zeta}$,
\[
|\omega_{E_{\emptyset,\beta,\zeta}}|
=\frac{|x(\zeta)|^6|z(\alpha)|^{10}}{{|w(\delta)|^7}}
dz(\alpha)dw(\gamma)dw(\delta)dw(\eps)dw(\zeta)dx(\zeta)
dv(\alpha)...dv(\zeta).
\]
We may assume that the support of $f$ meets ${\mathcal O}_3^{st}$ and no other
unipotent orbit. In this case, by Formula 0.6(b), we have
\[
\int_{E_{\emptyset,\beta,\zeta}(F)}\pi^*(f)\kappa_\theta(t_*)
|\omega_{E_{\emptyset,\beta,\zeta}}|
=\int_{{\mathcal O}_3(F)}f
\int_{E_{\emptyset,\beta,\zeta}(u)}\kappa_1(t_*)
|\omega_{E_{\emptyset,\beta,\zeta}(u)}||\omega_{{\mathcal O}_3}|,
\tag{8.2}
\]
which is a multiple of the $\kappa$-stable unipotent
orbital integral of $f$ associated to the $3-$regular orbit
of $G_2$. It is shown in the appendix that the $3$-regular orbit is stable.

\begin{proclaim}{Theorem 8.3} Suppose $T_G$ is a non-split Cartan subgroup
of $G_2$ which splits over a
Galois extension for which the Galois group action on $T_G$
is realized by a cyclic subgroup of $\Omega$ of order three.
Assume that $f\in C_c^\infty(G(F))$ satisfies $supp(f)\cap {\mathcal O}_3
\not= \emptyset$ and the support of $f$ does not meet other unipotent orbit
of $G(F)$. Then there is an $f^H$, smooth on
$H(F)$, for which, for all regular $\gamma_G\in T_G(F)$
in a sufficiently small neighborhood of $1$, we have
\[
\Delta_{G/H}^{T_G,\kappa}(\gamma_H,\gamma_G)
\Phi^{T_G,\kappa}(\gamma_G,f)=
\Phi^{T_H,st}(\gamma_H,f^H),
\]
where $\kappa$, $T_G$ satisfy $H_{T_G,\kappa}=PGL(3)$.
(In particular, $\kappa$ is a cubic character and $T_G$ is a cyclic
``cubic'' elliptic Cartan.)
\end{proclaim}

This will be proven below.

\subsection*{Proof of Theorem 8.3}

Note that by \cite{Ro80}, the germ on the
right-hand side, $\Gamma_{{\mathcal O}_H^{st}}^{T_H,st}(\gamma_H)$, is a
non-zero constant. Therefore, the theorem will follow if we can show that,
as a function of the tangent vector $X$ of $C$,
the two sides are equal up to a constant.

Let $u=\theta_{\beta}(1)\in G_2$.
One may explicitly compute the orbit $ad({\mathfrak g}_2)(X_\beta)$.
From this, we conclude that $u^{G_2}$
is 6-dimensional. Therefore ${\mathcal O}_3^{st}=u^{G_2}$.
There is only one $F$-orbit of $u$ contained in the stable orbit of $u$
(as is verified in the appendix), so the $3$-regular orbit
of $G_2$ is stable.

To get the local coordinates on the fiber
$E_{\emptyset,\beta,\zeta}(u)\subset E_{\emptyset,\beta,\zeta}$ over
$u$ in the 3-regular orbit ${\mathcal O}_3$ take the local
coordinates in the ``right'' divisor $E_{\emptyset,\beta,\zeta}$,
then divide $\omega_{E_{\emptyset,\beta,\zeta}}$ by a Haar measure
$\omega_{{\mathcal O}_3}$ on ${\mathcal O}_3$.
Analogously to Hales' lemma 6.4,
the subvariety $\prod_{\sigma\in \Sigma_G-\Sigma_{H_1}}U_\sigma$
of $G_2$ yields local coordinates for the $3$-regular orbit,
where $\Sigma_{H_1}^+=\{\beta,\epsilon,\zeta\}$.
We use the form on ${\mathcal O}_3$,
\[
\begin{split}
|\omega_{{\mathcal O}_3}|
&:=\frac{|\bigwedge_{\sigma\in \Sigma_{{\mathcal O}_3}}dx(\sigma)\wedge dv(\sigma)|}{|\lambda|^2}\\
&=\frac{|x(\zeta)|}{|w(\zeta)||z(\alpha)|}
dz(\alpha)dw(\gamma)dw(\delta)dv(\alpha)dv(\gamma)dv(\delta),
\end{split}
\tag{8.4}
\]
where $\Sigma_{{\mathcal O}_3}=\Sigma_{G}^+ -\Sigma_{H_1}^+$.
On the generic fiber $E_{\emptyset,\beta,\zeta}(u)$, $u\in {\mathcal O}_3$,
\[
|\omega_{E_{\emptyset,\beta,\zeta}(u)}|=
\frac{{|x(\zeta)||w(\zeta)|^5}}{{|z(\alpha)||w(\delta)|^7}}
dw(\epsilon)dw(\zeta)dx(\zeta)dv(\beta)dv(\epsilon)dv(\zeta).
\tag{8.5}
\]
We need to rewrite this last form as an expression
\[
\omega_{E_{\emptyset,\beta,\zeta}(u)}=
|p_{{\mathcal O}_3,T_G}(X)|\omega'_{E_{\emptyset,\beta,\zeta}(u)},
\]
where $p_{{\mathcal O}_3,T_G}(X)$ is a homogeneous polynomial in the coordinates
of tangent vector $X$ to be computed below
and $\omega'_{E_{\emptyset,\beta,\zeta}(u)}$
is independent of the tangent direction.

By an easy calculation given in the appendix,
our assumption on $T_G$ implies that $\kappa$ is a character
of order $3$.

Once we do this it will follow that
\[
\begin{split}
&D_H(\gamma_{\lambda,H})\sum_{{\mathcal O}\subset {\mathcal O}_3^{st}}
\mu_{\mathcal O}(f)\Gamma_{\mathcal O}^{T_G,\kappa}(\gamma_\lambda)
=|\lambda|^2\kappa(inv(\gamma_G^0,\gamma_\lambda))
\int_{E_{\emptyset,\beta,\zeta}}
\pi^*(f)\kappa_1(t_*)|\omega_{{\mathcal E}_{\emptyset,\beta,\zeta}(F)}|\\
&=|\lambda|^2\kappa(inv(\gamma_G^0,\gamma_\lambda))
|p_{{\mathcal O}_3,T_G}(X)|
\sum_{{\mathcal O}\subset {\mathcal O}_3^{st}}\int_{{\mathcal O}_3(F)}f
\int_{E_{\emptyset,\beta,\zeta}(u)(F)}\kappa(t_*)
|\omega'_{E_{\emptyset,\beta,\zeta}(u)}||\omega_{\mathcal O}|,
\end{split}
\tag{8.6}
\]
where we normalize $\omega_{{\mathcal O}}$ consistent with the
Langlands-Shelstad transfer factor as in \cite{Ha89}.

To determine the general form of
$p_{{\mathcal O}_3,T_G}(X)$ we use the coordinate relations in (1.22).
We find that $p_{{\mathcal O}_3,T_G}(X)$ is a product of the form
\[
p_{{\mathcal O}_3,T_G}(X)=\prod_{\sigma \in\Sigma_H^+}\sigma(X)^{e_\sigma},
\tag{8.7}
\]
where $\sum_{\sigma \in\Sigma_H^+}e_\sigma=3$.
One interesting thing to note is that the product runs over roots
``in $H$''.

To complete the matching, it remains only to
verify the local translation invariance property in the hypothesis
to Hales' lemma 5.19 which allows us to apply the
descent argument.
By continuity, we see that the factor
$\kappa(inv(\gamma_G^0,\gamma_{G,\lambda}))\kappa(t_*)$ satisfies the
local translation invariance property of lemma 5.19 above, for
sufficiently small $Y$. One checks that the factor (8.7)
satisfies the translation invariance property using that
it is a product over roots in $H$. The theorem now follows from
lemma 5.19.

% ======================================================================
% Claude: Section 9, source file g2new9.tex (\pageno 52). Here \off = {\Cal O}_F
% (not used in g2new9.tex; used in the VARIANT g2new9d.tex below).
% ======================================================================
\section*{\S\,9  A conjecture}

We conclude with a conjecture, which seems to be indicated by the
proof above.

\begin{proclaim}{Conjecture 9.1} Suppose $T_G$ is an elliptic Cartan subgroup
of a connected, simply connected, $F$-quasi-split group $G$. Let
${\mathcal O}$ denote an $r$-regular unipotent orbit of $G(F)$,
let $\kappa$ denote a character of $H^1(\Gamma,T_G)$, and let
$H=H_{\kappa,T_G}$ denote an endoscopic group of $G$ \cite{La79}.
Assume

(1) there is an essential, admissible $F$-rational ${\mathcal O}$-divisor
$E=E_{S_1,S_2,\sigma}$ such that, for each $u\in {\mathcal O}$, the
Igusa form on the fiber satisfies
\[
\omega_{E(u)}=p_{{\mathcal O},T_G}(X)\omega'_{E(u)},
\]
where $\omega'_{E(u)}$ is independent of $X$ and where
\[
p_{{\mathcal O},T_G}(X)=\prod_{\sigma\in\Sigma_{\mathcal O}^+}\sigma(X),
\]
for some $\Sigma_{\mathcal O}^+\subset \Sigma_G^+$,

(2) the function $p_{{\mathcal O},T_G}$ in (1) satisfies the hypotheses to
Hales' reduction lemma,

(3) $b_{E,i}\not= 0$ in the $F$-rational coordinates of $Y''(B_+)$,

(4) $e(\theta,r+1)\leq 1$ and $e(\theta,r+1)=0$ unless $\theta =1$,

(5) there is only one $r$-regular unipotent orbit ${\mathcal O}$
of $G(\overline{F})$, and
it is either Richardson or endoscopically Richardson,

(6) there is only one $r$-regular orbit ${\mathcal O}_H$ of $H(\overline{F})$.

Then the matching of germs conjecture stated in the introduction
holds between these $r$-regular orbits. Moreover,
for each $f\in C_c^\infty (G(F))$ such that
$supp(f)\cap {\mathcal O}'=\emptyset$ if ${\mathcal O}'$
is not stably conjugate to ${\mathcal O}$, we have,
for all regular $\gamma_G\in T_G(F)$
in a sufficiently small neighborhood of $1$ and for any
function $f$,
\[
\Delta_{G/H}^{T_G,\kappa}(\gamma_H,\gamma_G)
\Phi^{T_G,\kappa}(\gamma_G,f)=
\Phi^{T_H,st}(\gamma_H,f^H),
\]
for some $f^H\in C_c^\infty (H(F))$ such that
$supp(f^H)\cap {\mathcal O}_H'=\emptyset$ if ${\mathcal O}_H'$
is not stably conjugate in $H$ to ${\mathcal O}_H$.
\end{proclaim}

Finally, we remark that there are a several low-dimensional cases which
appear to
satisfy the hypotheses of this conjecture. We mention one:

$G=SL(4)$, ${\mathcal O}={\mathcal O}_{(2,2)}$ the $2$-regular orbit associated
to the partition $4=2+2$, $H=GL(2,E)^1$ where $E/F$ is quadratic
and the superscript one denotes the kernel of the map
$Nm_{E/F}\circ \det :GL(2,E)\rightarrow F^{\times}$, and
${\mathcal O}_H$ is the orbit of $H(F)$ associated to the
trivial $1$-regular orbit of $GL(2,E)$ by restriction of scalars.

% Claude: VARIANT. g2new9d.tex has the same \S 9 heading and the same opening
% sentence. Its Conjecture 9.1 and closing paragraphs, transcribed below,
% replace the corresponding text above. In g2new9d.tex \off = {\Cal O}_F,
% so hypotheses (5) and (6) print G(O_F) and H(O_F); compare G(\overline{F})
% in g2new9.tex. The labels [A] and [H4] do not occur in the bibliography of
% the source and are typeset literally.
\iffalse
\begin{proclaim}{Conjecture 9.1} Suppose $T_G$ is an elliptic Cartan subgroup
of a connected, simply connected, $F$-quasi-split group $G$. Let
${\mathcal O}$ denote an $r$-regular unipotent orbit of $G(F)$,
let $\kappa$ denote a character of $H^1(\Gamma,T_G)$, and let
$H=H_{\kappa,T_G}$ denote an endoscopic group of $G$ \cite{La79}.
Assume

(1) there is an essential, admissible $F$-rational ${\mathcal O}$-divisor
$E=E_{S_1,S_2,\sigma}$ such that, for each $u\in {\mathcal O}$, the
Igusa form on the fiber satisfies
\[
\omega_{E(u)}=p_{{\mathcal O},T_G}(X)\omega'_{E(u)},
\]
where $\omega'_{E(u)}$ is independent of $X$ and where
\[
p_{{\mathcal O},T_G}(X)=\prod_{\sigma\in\Sigma_{\mathcal O}^+}\sigma(X),
\]
for some $\Sigma_{\mathcal O}^+\subset \Sigma_G^+$,

(2) the function $p_{{\mathcal O},T_G}$ in (1) satisfies the hypotheses to
Hales' reduction lemma,

(3) $b_{E,i}\not= 0$ in the $F$-rational coordinates of $Y''(B_+)$,

(4) $e(\theta,r+1)\leq 1$,

(5) there is only one $r$-regular unipotent orbit of $G(\OF)$, and
it is either Richardson or endoscopically Richardson,

(6) the orbit ${\mathcal O}$ of $G(\OF)$ corresponds to an $s$-regular orbit
${\mathcal O}_H$ of $H(\OF)$ under the endoscopic correspondence
of Assem-Hales-Spaltenstein [A], [H4], \cite{Spa82}, which
is the only $s$-regular orbit of $H$,

Then, for each $f\in C_c^\infty (G(F))$ such that
$supp(f)\cap {\mathcal O}'=\emptyset$ if ${\mathcal O}'$
is not stably conjugate to ${\mathcal O}$, we have,
for all regular $\gamma_G\in T_G(F)$
in a sufficiently small neighborhood of $1$ and for any
function $f$,
\[
\Delta_{G/H}^{T_G,\kappa}(\gamma_H,\gamma_G)
\Phi^{T_G,\kappa}(\gamma_G,f)=
\Phi^{T_H,st}(\gamma_H,f^H),
\]
for some $f^H\in C_c^\infty (H(F))$ such that
$supp(f^H)\cap {\mathcal O}_H'=\emptyset$ if ${\mathcal O}_H'$
is not stably conjugate in $H$ to ${\mathcal O}_H$.
\end{proclaim}

Finally, we remark that there are a several low-dimensional cases which
appear to
satisfy the hypotheses of this conjecture. We mention one:

$G=SL(4)$, ${\mathcal O}={\mathcal O}_{(2,2)}$ the $2$-regular orbit associated
to the partition $4=2+2$, $H=GL(2,E)^1$ where $E/F$ is quadratic
and the superscript one denotes the kernel of the map
$Nm_{E/F}\circ \det :GL(2,E)\rightarrow F^{\times}$, and
${\mathcal O}_H$ is the orbit of $H(F)$ associated to the
trivial $1$-regular orbit of $GL(2,E)$ by restriction of scalars.

Details will appear later.
\fi
% ======================================================================
% Claude: Appendix, source file g2newapp.tex (\pageno 53).
% ======================================================================
\section*{Appendix: Background on $G_2$}

The first four sections of this appendix have been taken from an
unpublished joint paper with D. Bump.

The purpose of this appendix is to collect some foundational material on
the exceptional group $G_2$, then to use this material to
compute the cuspidal endoscopic groups of
$G_2$, the groups $H^1(\Gamma, T)$, and the $F$-classes in the
stable 2- and 3-regular unipotent orbit.

In the first few sections, we shall be concerned with the phenomenon
known as ``triality'',
which yields an embedding of $G_2$ into the double cover Spin(8) of
the split form of SO(8). We shall
determine the triality map explicitly at the root level, and also use it to
give an explicit Chevalley basis for the Lie algebra ${\mathfrak g}_2$,
which should be useful in other calculations.

\noindent
Acknowledgement: We thank Becci Davies for coding the root diagrams
(in AMS--TEX).

\subsection*{Background on G$_2$ and SO(8)}

\subsection*{A.1. SO(8)}

Let J denote the $8\times 8$ matrix with 1's on the skew--diagonal and
0's elsewhere:
\[
J := \begin{pmatrix} 0 & & 1 \\
   & \iddots & \\
   1 & & 0
   \end{pmatrix}.
\]
Define the $8\times 8$ \textbf{special orthogonal group}
by
\[
SO(8) := \{A \in GL(8) \ |\ A\cdot J\cdot {}^{t}A = J\}^o,
\]
so that $SO(8)$ is semi--simple, connected, neither simply connected
nor adjoint, non--compact, and split over $\qqq$. Its maximal $\qqq$-split
torus consists of all matrices of the form
\[
\begin{pmatrix} A&0\\0&B \end{pmatrix},\ \text{with }
A:=\begin{pmatrix} a_1&&&\\
&a_2&&\\
&&a_3&\\
&&&a_4 \end{pmatrix},\
B:= \begin{pmatrix} a_4^{-1}&&&\\
&a_3^{-1}&&\\
&&a_2^{-1}&\\
&&&a_1^{-1} \end{pmatrix}.\tag{A.1}
\]
The compact real form of $SO(8)$ is obtained, of course, by replacing
J by the $8\times 8$ identity matrix.

The simply--connected covering group of $SO(8)$ is
$SO(8)_{sc} =$ Spin(8), which has center $\zzz /2\zzz \times \zzz /2\zzz$,
and since $SO(8)$ has center $\zzz /2\zzz$, Spin(8) is a 2-fold cover
of $SO(8)$ and a 4-fold cover of $PSO(8)$. The outer automorphisms
of Spin(8) are in one--to--one correspondence with the automorphisms
of the Dynkin diagram for $SO(8)$:
% Claude: The source draws the four nodes as bullets in an AMS-TeX matrix,
% without edges. The same arrangement is reproduced with TikZ.
\[
\begin{tikzpicture}[baseline=(c.base)]
\node (c) at (0,0) {$\bullet$};
\node at (-1.6,0) {$\bullet$};
\node at (1.1,1.0) {$\bullet$};
\node at (1.1,-1.0) {$\bullet$};
\end{tikzpicture}
\]
The Lie algebra
\[
so(8) := \{A \in M_8\ |\ AJ + J{}^{t}A = 0\}
\]
may be expressed as the Lie algebra of $8\times 8$ matrices of the
form
\[
A =  {\begin{pmatrix} t_{1}&a_{12}&a_{13}&a_{14}&a_{15}&a_{16}&a_{17}&0\\
a_{21}&t_2&a_{23}&a_{24}&a_{25}&a_{26}&0&-a_{17}\\
a_{31}&a_{32}&t_3&a_{34}&a_{35}&0&-a_{26}&-a_{16}\\
\vdots&\vdots &\vdots &\vdots &\vdots &\vdots &\vdots &\vdots\\
0&-a_{71}&-a_{61}&-a_{51}&-a_{41}&-a_{31}&-a_{21}&-t_1\\
\end{pmatrix}}.\tag{A.2}
\]
If $i \not= j$ then define $E_{ij}$ to be the matrix A as above with
$a_{ij} = 1$ and all other entries = 0. Define $E_{ii}$ to be the
matrix A as above with $t_i = 1$ and all other entries = 0. Then
\[
\begin{split}
&\big\{ E_{11},\,  E_{22},\, E_{33},\, E_{44},\\
&\quad E_{12},\, E_{13},\dots ,E_{17},
       E_{23},\dots ,E_{26},
       E_{34},\, E_{35},\\
&\quad {{}^{t}E_{12}},\, {{}^{t}E_{13}},\dots ,{{}^{t}E_{17}},
       {{}^{t}E_{23}},\dots ,{{}^{t}E_{26}},
       {{}^{t}E_{34}},{{}^{t}E_{35}} \big\}
\end{split}\tag{A.3}
\]
forms a basis
for the 28-dimensional vector space $so(8)$.

\subsection*{A.2. The Cayley algebra and $G_2$}

The \textbf{split Cayley algebra} is the $2\times 2$ matrix algebra
\[
C:=\Big\{\begin{pmatrix} \alpha&a\\b&\beta \end{pmatrix}\ | \
\alpha ,\beta \in \qqq,\ a,b \in V \Big\},\tag{A.4}
\]
where $V \cong \qqq^3$ is the cross product algebra
\[
\begin{split}
V:= \{ xi+yj+zk\ | \ x,y,z \in \qqq,
&\ i\times i = j\times j = k\times k = 0,\\
&\quad\quad
\ i\times j = k,\ j\times k = i,\ k\times j = i,\ j\times i = -k,
\ \text{etc.}\}.
\end{split}
\]
Define
\[
G_2(F) :=\text{Aut}_F(C\otimes F),\tag{A.5}
\]
where F is any field containing $\qqq$. Since Der(C) $\cong$ Lie(Aut C),
it follows that the Lie algebra ${\mathfrak g}_2$ may be
identified with Der(C). For $G_2$, there are 2 real forms:  one
being the non-compact $\qqq$--split from above, the other being a
compact form associated to a division algebra (see\cite[pp. 142-147]{Ja71}
or \cite{Ti65}). We have
\[
\text{dim}_{{\qqq}}C = 8,\ \text{dim}_{{\qqq}}{\mathfrak g}_2 = 14,
\]
and there is a trace bilinear form which induces an embedding
\[
G_2 = \text{Aut}(C) \hookrightarrow SO(8).
\]
Here SO(8) is regarded as the special orthogonal group associated to
the trace form $trace(\beta)$ on C, where
\[
\beta(x,y)=\frac{1}{2}(x\cdot y^*+y^*\cdot x),
\ \ \ \ \beta(x,x)=N(x),
\]
and where $N$ denotes the norm on $C$. Regarding $G_2$ as a subgroup of
$SO(8)$, we have the following

\begin{proclaim}{Triality theorem A.6(\cite{Ja71})} Given any orthogonal map
\[
A_0 : C \rightarrow C,
\]
there are orthogonal maps $A_1,\ A_2$ such that, for all $x,\,y \in C$,
\[
A_0(xy) = A_1(x)A_2(y).
\]
Moreover, the only other orthogonal maps satisfying this property are
$-A_1,\ -A_2$, so
\[
A_0 \longmapsto (A_1,A_2)
\]
is well--defined, up to signs.
\end{proclaim}

\begin{proof} We will verify that our form of triality is the global
equivalent of Jacobson's form \cite[p. 8]{Ja71}. Let
\[
\begin{split}
&A_0 = 1+\delta B_0,\\
&A_1 = 1+\delta B_1,\\
&A_2 = 1+\delta B_2,
\end{split}
\]
in $SO(8,\qqq )$ (regarded as a group of orthogonal maps with respect to the
trace form on C), where $\delta^2 = 0$ and $B_i \in so(8,\qqq )$. Thus
$A_0,\ A_1,\ A_2$ may be regarded as lying in the tangent space to SO(8)
at the identity (cf. \cite[chapter 0]{Bor91}). Then
$A_0(xy) = A_1(x)A_2(y)$ is equivalent to
\[
B_0(xy) = B_1(x)y+ xB_2(y).
\]
This is Jacobson's expression.
\end{proof}

\begin{proclaim}{Corollary A.7} The map
\[
\begin{split}
SO(8,F) &\rightarrow PSO(8,F),\\
\quad A_0 &\longmapsto \quad \pm A_1,
\end{split}
\]
defines a projective representation of SO(8), i.e., belongs to
PSO(8). In particular, it induces an outer automorphism of degree 3
(the ``triality'' automorphism)
corresponding to a rotation of the Dynkin diagram of $PSO(8,F)$.
\end{proclaim}

\noindent
\textit{Remarks} (A.8) If $A_0 \in G_2 = \text{Aut}(C)$ then triality sends
$A_0 \longmapsto A_0$. Conversely, if triality sends $A_0 \longmapsto A_0$
then $A_0 \in G_2$, so $G_2$ is the subgroup fixed by triality. In
general, triality sends (up to signs)
\[
A_0 \longmapsto A_1,\ A_1 \longmapsto A_2,\ A_2 \longmapsto A_0.
\]
This way we can explicitly construct an outer automorphism of $SO(8)$.

(A.9) Triality can be described explicitly on the standard Borel
subgroup B of PSO(8): triality cyclically permutes the ``root groups''
$U_{\alpha}$, defined below, in the same way as it permutes the simple roots
via its action on the Dynkin diagram; it also permutes the embeddings
$PSL(2) \hookrightarrow PSO(8)$ associated to the simple roots. The
images of the maximal split torus $T \subset PSL(2)$ embedded in $PSO(8)$
generate the maximal split torus in $PSO(8)$ and, of course, the simple
root groups generate the unipotent radical of B. Putting these all together
explicitly gives the action of triality on B.

\subsection*{A.3. Root systems for $G_2$ and $SO(8)$}

Since the symmetric group $S_3$ acts on the Dynkin diagram of $SO(8)$
by triality,
it acts on its 24 roots. Of these, 6 roots remain fixed under
this action, and 18 are permuted cyclically in groups of three. In
particular, the roots corresponding to
\[
\{E_{17}\},\ \ \{E_{16}\},\ \ \{E_{23}\},
\]
are fixed and the triples of roots corresponding to
\[
\{E_{14},\ E_{15},\ E_{26}\},\ \{E_{13},\ E_{24},\ E_{25}\},
\ \{E_{12},\ E_{34},\ E_{35}\},
\]
are each permuted by the action of $S_3$. More explicitly, let us
represent a basis of $so(8)$ by
\[
\eps_1 - \eps_2,\  \eps_2 - \eps_3,\ \eps_3 - \eps_4,\ \eps_3 + \eps_4,
\]
where
\[
\eps_i : \begin{pmatrix} A&0\\0&B\end{pmatrix}
\longmapsto a_i, \tag{A.10}
\]
in the notation of (A.1), $1 \leq i \leq 4$.
Here $E_{12}$ corresponds to $\eps_1 - \eps_2,\ E_{13}$ to
$\eps_1 - \eps_2 + \eps_2 - \eps_3 = \eps_1 - \eps_3$, etc. Since
$\eps_2 - \eps_3$ is connected to the others (since its Killing form inner
product with them is non--zero), it corresponds to the center of the
Dynkin diagram
\[
\begin{tikzpicture}[baseline=(c.base)]
\node (c) at (0,0) {$\bullet$};
\node[right] at (0.1,0) {$\eps_2-\eps_3$};
\node at (-2.0,0) {$\bullet$};
\node[left] at (-2.1,0) {$\eps_3-\eps_4$};
\node at (1.4,1.0) {$\bullet$};
\node[right] at (1.5,1.0) {$\eps_1-\eps_2$};
\node at (1.4,-1.0) {$\bullet$};
\node[right] at (1.5,-1.0) {$\eps_3+\eps_4$};
\end{tikzpicture}
\]
Therefore, under the action of $S_3$, the triality automorphism of
so(8) induces a linear map $\phi$ on the roots such that
\[
\phi (\eps_1 - \eps_2) = \eps_3 - \eps_4,
\ \phi (\eps_3 - \eps_4) = \eps_3 + \eps_4,
\ \phi (\eps_3 + \eps_4) = \eps_1 - \eps_2,
\ \phi (\eps_2 - \eps_3) = \eps_2 - \eps_3.\tag{A.11}
\]
Thus
\[
\begin{split}
&\phi (\eps_1) = \frac{1}{{2}}(\eps_1 + \eps_2 + \eps_3 - \eps_4),\\
&\phi (\eps_2) = \frac{1}{{2}}(\eps_1 + \eps_2 - \eps_3 + \eps_4),\\
&\phi (\eps_3) = \frac{1}{{2}}(\eps_1 - \eps_2 + \eps_3 + \eps_4),\\
&\phi (\eps_4) = \frac{1}{{2}}(\eps_1 - \eps_2 - \eps_3 - \eps_4).
\end{split}
\]
Call the induced automorphism on the Lie algebra so(8), $\phi^*$. As
we've already observed, the Lie subalgebra
\[
\{X \in so(8)\ |\ \phi^*(X) = X\}
\]
is the Lie algebra ${\mathfrak g}_2$. Identifying its Lie
group with a subgroup of $SO(8)$, we find that its maximal
$\qqq-$split torus is of the form
\[
\begin{split}
T':= \Big\{\begin{pmatrix} A&0\\0&B \end{pmatrix},\ |\
&A:={\left(\begin{smallmatrix} a_1&&&\\
&a_2&&\\
&&a_3&\\
&&&1 \end{smallmatrix}\right)},\
B:= {\left(\begin{smallmatrix} 1&&&\\
&a_3^{-1}&&\\
&&a_2^{-1}&\\
&&&a_1^{-1} \end{smallmatrix}\right)}\\
&\quad\quad\quad\quad\text{and } a_1 = a_2a_3\Big\}.
\end{split}\tag{A.12}
\]
We will identify $T'$ with the maximal split torus of
$SL(3)$ by means of the map
\[
\begin{pmatrix} A&0\\0&B \end{pmatrix} \in T' \longmapsto
\text{diag}(a_2,a_3,a_1^{-1}),
\]
with the maximal split torus of
$PGL(3)$ by means of the map
\[
\begin{pmatrix} A&0\\0&B \end{pmatrix} \in T' \longmapsto
\text{diag}(a_1a_2,a_1a_3,1),\tag{A.13}
\]
with the maximal split torus of
$PGL(2)\times SL(2)$ by means of the map
\[
\begin{pmatrix} A&0\\0&B \end{pmatrix} \in T' \longmapsto
\text{diag}(a_3,1)\times diag(a_1a_2,1/a_1a_2),
\tag{A.13$'$}
\]
and with the maximal split torus of
$(SL(2)\times SL(2))/\{\pm 1\}$ by means of the map
\[
\begin{pmatrix} A&0\\0&B \end{pmatrix} \in T' \longmapsto
\text{diag}(a_1,1/a_1)\times diag(a_2,1/a_2).
\tag{A.13$''$}
\]

\subsection*{A.4. The Chevalley basis}

The embedding of ${\mathfrak g}_2$ into so(8) gives a convenient framework
for doing calculations with explicit matrices. We now write down
the Chevalley basis of $G_2$. The notation for the root system is given
by the following diagram:
% Claude: The source draws the twelve roots as bullets placed in an AMS-TeX
% matrix (only the positive roots are labeled). The same picture is
% reproduced with TikZ at the actual root positions, with no lines.
\[
\begin{tikzpicture}[x=1.6cm,y=1.6cm]
\foreach \p/\l in {{(0,1.732)}/{3\alpha+2\beta}, {(-1.5,0.866)}/{\beta},
  {(-0.5,0.866)}/{\alpha+\beta}, {(0.5,0.866)}/{2\alpha+\beta},
  {(1.5,0.866)}/{3\alpha+\beta}, {(1,0)}/{\alpha}}
  {\node at \p {$\bullet$}; \node[above] at \p {$\scriptstyle \l$};}
\foreach \p in {(-1,0),(-1.5,-0.866),(-0.5,-0.866),(0.5,-0.866),(1.5,-0.866),(0,-1.732)}
  {\node at \p {$\bullet$};}
\node at (0,0) {$\cdot$};
\end{tikzpicture}
\]
\centerline{the root system for $G_2$}

\medskip
\[
X_{\alpha} :=
\begin{pmatrix} 0&1& & & & & & \\
          &0& & & & & & \\
          & &0&1&1& & & \\
          & & &0&0&-1& & \\
          & & & &0&-1& & \\
          & & & & &0& & \\
          & & & & & &0&-1\\
          & & & & & & &0
         \end{pmatrix},\ \
h_{\alpha} := \begin{pmatrix} 1& & & & & & & \\
                          &-1& & & & & & \\
                          & &2& & & & & \\
                          & & &0& & & & \\
                          & & & &0& & & \\
                          & & & & &-2& & \\
                          & & & & & &1& \\
                          & & & & & & &-1
                         \end{pmatrix}.
\]
(Note that $X_{\alpha} = E_{12} + E_{34} + E_{35}$ and
$\alpha\begin{pmatrix} A& \\ &B\end{pmatrix} =a_3$, where $A,B$ are as
in (A.12).)
\[
X_{\beta} :=
\begin{pmatrix} 0& & & & & & & \\
          &0&1& & & & & \\
          & &0& & & & & \\
          & & &0& & & & \\
          & & & &0& & & \\
          & & & & &0&-1& \\
          & & & & & &0&\\
          & & & & & & &0
         \end{pmatrix},\ \
h_{\beta} := \begin{pmatrix} 0& & & & & & & \\
                          &1& & & & & & \\
                          & &-1& & & & & \\
                          & & &0& & & & \\
                          & & & &0& & & \\
                          & & & & &1& & \\
                          & & & & & &-1& \\
                          & & & & & & &0
                         \end{pmatrix}.
\]
(Note that $X_{\beta} = E_{23}$ and
$\beta\begin{pmatrix} A& \\ &B\end{pmatrix} =a_2a_3^{-1}$, where $A,B$ are as
in (A.12).) For $\sigma = \alpha+\beta$:
\[
X_{\sigma} :=
\begin{pmatrix} 0& &-1& & & & & \\
          &0& &1&1& & & \\
          & &0& & & & & \\
          & & &0& & &-1& \\
          & & & &0& &-1& \\
          & & & & &0& &1\\
          & & & & & &0& \\
          & & & & & & &0
         \end{pmatrix},\ \
h_{\sigma} := \begin{pmatrix} 1& & & & & & & \\
                          &2& & & & & & \\
                          & &-1& & & & & \\
                          & & &0& & & & \\
                          & & & &0& & & \\
                          & & & & &1& & \\
                          & & & & & &-2& \\
                          & & & & & & &-1
                         \end{pmatrix}.
\]
(Note that $X_{\alpha+\beta} = E_{13}+E_{24}+E_{25}.$) For
$\sigma = 2\alpha+\beta$:
\[
X_{\sigma} :=
\begin{pmatrix} 0& & &1&1& & & \\
          &0& & & &1& & \\
          & &0& & & &-1& \\
          & & &0& & & &-1\\
          & & & &0& & &-1\\
          & & & & &0& & \\
          & & & & & &0& \\
          & & & & & & &0
         \end{pmatrix},\ \
h_{\sigma} := \begin{pmatrix} 2& & & & & & & \\
                          &1& & & & & & \\
                          & &1& & & & & \\
                          & & &0& & & & \\
                          & & & &0& & & \\
                          & & & & &-1& & \\
                          & & & & & &-1& \\
                          & & & & & & &-2
                         \end{pmatrix}.
\]
(Note that $X_{2\alpha + \beta} = E_{14}+E_{15}+E_{26}$.) For
$\sigma = 3\alpha+\beta$:
\[
X_{\sigma} :=
\begin{pmatrix} 0& & & & &1& & \\
          &0& & & &0& & \\
          & &0& & &0&0&-1\\
          & & &0& & & & \\
          & & & &0& & & \\
          & & & & &0& & \\
          & & & & & &0& \\
          & & & & & & &0
         \end{pmatrix},\ \
h_{\sigma} := \begin{pmatrix} 1& & & & & & & \\
                          &0& & & & & & \\
                          & &1& & & & & \\
                          & & &0& & & & \\
                          & & & &0& & & \\
                          & & & & &-1& & \\
                          & & & & & &0& \\
                          & & & & & & &-1
                         \end{pmatrix}.
\]
(Note that $X_{3\alpha + \beta} = E_{16}$.) For $\sigma =
3\alpha+2\beta$:
\[
X_{\sigma} :=
\begin{pmatrix} 0& & & & & &1& \\
          &0& & & & &0&-1\\
          & &0& & & & & \\
          & & &0& & & & \\
          & & & &0& & & \\
          & & & & &0& & \\
          & & & & & &0& \\
          & & & & & & &0
         \end{pmatrix},\ \
h_{\sigma} := \begin{pmatrix} 1& & & & & & & \\
                          &1& & & & & & \\
                          & &0& & & & & \\
                          & & &0& & & & \\
                          & & & &0& & & \\
                          & & & & &0& & \\
                          & & & & & &-1& \\
                          & & & & & & &-1
                         \end{pmatrix}.
\]
(Note that $X_{3\alpha+2\beta} = E_{17}$.)

There is a canonical embedding $G_2 \hookrightarrow SO(8)$, given by
the natural faithful representation associated to the Cayley algebra.
In fact, each one--parameter unipotent subgroup $U_{\sigma}$ associated to a
short root $\sigma$ of $G_2$ is contained in a three dimensional unipotent
subgroup of SO(8) generated by three (uniquely determined) roots of
$D_4$ which are permuted by $S_3$.
The possibilities for this embedding of unipotents can be tabulated
as follows:
\[
\begin{split}
&U_{\alpha} \subset G_2 \longleftrightarrow
U := \exp[E_{12}F+E_{34}F+E_{35}F],\\
&U_{2\alpha +\beta} \subset G_2 \longleftrightarrow
U := \exp[E_{14}F+ E_{15}F+E_{26}F],\\
&U_{\alpha+\beta} \subset G_2 \longleftrightarrow
U := \exp[E_{13}F+E_{24}F+E_{25}F].
\end{split}\tag{A.14}
\]
For the long roots of $G_2$, we have
\[
\begin{split}
&U_{\beta} \subset G_2 \longleftrightarrow
U := \exp[E_{23}F],\\
&U_{3\alpha +2\beta} \subset G_2 \longleftrightarrow
U := \exp[E_{17}F],\\
&U_{3\alpha+\beta} \subset G_2 \longleftrightarrow
U := \exp[E_{16}F].
\end{split}\tag{A.15}
\]
In other words, under the action of $S_3$ on the roots of
$D_4$, the six stable roots correspond to the long roots of $G_2$
and the other 18, permuted in orbits of 3, correspond to the short
roots of $ G_2$.

\subsection*{A.5. Endoscopic groups of $G_2$}

The calculation of the cuspidal endoscopic groups of $G_2$ is as
follows. Motivated by facts
about root systems of centralizers of semi-simple elements
\cite[Ch III, \S 12]{Spa82}, let $\Delta:=\{\alpha,\beta\}$,
$\Delta_1:=\{\alpha,3\alpha+2\beta\}$,
$\Delta_2:=\{\beta,3\alpha+\beta\}$,
$\Delta^{\vee}:=\{\alpha^{\vee},\beta^{\vee}\}$,
$\Delta_1^{\vee}:=\{\alpha^{\vee},(3\alpha+2\beta)^{\vee}\}$,
and $\Delta_2^{\vee}:=\{\beta^{\vee},(3\alpha+\beta)^{\vee}\}$,
where $\sigma^{\vee}=2\sigma/(\sigma,\sigma)$.
Representing the root $\alpha$ by the vector
$(1,0)$ and $\beta$ by $(-\frac{3}{2},\frac{\sqrt{3}}{2})$,
we find that the root lattice $Q$ is given by
\[
Q={\mathbb Z}[\Delta]={\mathbb Z}[(1,0),(-\frac{3}{2},\frac{\sqrt{3}}{2})].
\]
We have $\alpha^{\vee}=(2,0)$ and $\beta^{\vee}=(-1,\frac{1}{\sqrt{3}})$.
Since $G_2$ is adjoint, its character lattice $X$ is equal to its
root lattice: $X=Q$. $G_2$ is also simply connected, so the weight lattice
\[
\begin{split}
P
&=\{x\in {\mathbb R}^2\ |\ (x,\Sigma^{\vee})\subset {\mathbb Z}\}\\
&=\{(a,b)\ |\ (a,b)\cdot (2,0)\in {\mathbb Z},\
(a,b)\cdot (-1,\frac{1}{\sqrt{3}})\in {\mathbb Z}\}\\
&={\mathbb Z}[(1/2,0),(0,\frac{\sqrt{3}}{2})]
\end{split}
\]
is also equal to the character lattice. Similarly,
$Q^{\vee}={\mathbb Z}[(2,0),\allowbreak(-1,\frac{1}{\sqrt{3}})]\allowbreak=P^{\vee}$.
We have $Q_1={\mathbb Z}[(1,0),(0,\sqrt{3})]$,
$P_1={\mathbb Z}[(1/2,0),(0,\sqrt{3}/2)]$,
$Q_1^{\vee}={\mathbb Z}[(2,0),(0,2/\sqrt{3})]$,
and $P_1^{\vee}={\mathbb Z}[(1,0),(0,1/\sqrt{3})]$.
Fix a maximal split torus $^LT^o$ of $^LG_2^o$ shared by $^LH^o$
and identify $X$ with the lattice of characters of $^LT^o$. Take
$^LH^o$ to be of type $A_1\times A_1$, with root system generated by
$\beta^{\vee}$ (which is now short) and $(3\alpha+2\beta)^{\vee}$
(which is not equal to $3\beta^{\vee}+2\alpha^{\vee}$).
The character lattice is equal to $X_1^{\vee}=X^{\vee}=Q^{\vee}$
but its root lattice $Q_1^{\vee}$ is strictly contained in $Q^{\vee}$ and
$Q^{\vee}=X_1^{\vee}$ strictly contains $P^{\vee}$. Therefore, $^LH^0$
is neither
simply connected nor adjoint if $^LH^o$ is of type $A_1\times A_1$.
This forces either $H=PGL(2)\times SL(2)$ or
$H=(SL(2)\times SL(2))/\{\pm 1\}$. Following a suggestion of T. Hales, we now
show that the former possibility may be discarded. The center
$Z_{^LH^o}$ is of order two: let $z\in Z_{^LH^o}$ and note that
$\alpha^{\vee}(z)=\zeta^{\vee}(z)=1$ (the dual roots acting on the maximal
split torus of the dual group); these imply that
$\beta^{\vee}(z)=\pm 1$, which means that
\[
Z_{^LH^o}({\mathbb C})=\{diag(\epsilon,\epsilon,1,1,1,1,\epsilon,\epsilon)\ |\
\text{ where }\epsilon =\pm 1\}.
\]
The image of this in $PGL(2)\times SL(2)$, under the embedding in
(A.13) is trivial, so $H=(SL(2)\times SL(2))/\{\pm 1\}$.
Similarly, using $\Delta_2$ in place of $\Delta_1$,
one finds that $^LH^o\ \subset\ ^LG_2^o$
is simply connected if $^LH^o$ is of type $A_2$. (This fact
can also be found in \cite[ch II, 5.3-5.5]{SS86}.)
Therefore $H=PGL(3)$ and $H=(SL(2)\times SL(2))/\{\pm 1\}$, as desired.

\subsection*{A.6. Tori of $G_2$}

The stable conjugacy classes of Cartans in a quasi-split group $G$
are in 1-1 correspondence with the cohomology classes in
$H^1(\Gamma, W)$, where $W$ denotes the absolute Weyl group of
a fixed Cartan $T$ of $G$ defined over $F$ \cite[\S 4, Lemma 3]{Ro80}.
One may also determine the Cartans from the $\Gamma$-module
structure on their lattice of characters \cite{Sa71}.

Let $T$ be a Cartan of $G_2$ defined over $F$ which splits over a Galois
extension $K$. The Galois
group $\Gamma_K:=Gal(K/F)$ acts on the character lattice
$X=X(T)$, as does the Weyl group $W$. Recall
that $T$ is determined (up to $F$-isomorphism)
from the $\Gamma_K$-module $X$. The $\Gamma_K$-invariant submodules
\[
X^{\Gamma_K}:=\{\chi\in X\ |\ \chi^\sigma =\chi,\ \forall\sigma\in\Gamma_K\},
\]
and
\[
X_{\Gamma_K}:=\{\chi\in X\ |\ \sum_{\sigma\in\Gamma_K}\chi^\sigma =0\},
\]
satisfy the property that $X/X^{\Gamma_K}$ and $X/X_{\Gamma_K}$ are
torsion free. Let $T_0$ denote the annihilator of
$X^{\Gamma_K}$ in $T$ and let $A$ denote the annihilator of
$X_{\Gamma_K}$ in $T$.
Then $T_0\cap A$ is finite, $T=T_0A$, $A$ is the largest $F$-split
subtorus of $T$, and $T_0$ is the largest $F$-anisotropic subtorus of $T$.
Let $\Sigma_{\Gamma_K}\subset \Sigma$ denote the closed
subsystem defined by $\Sigma_{\Gamma_K}:=X_{\Gamma_K}\cap \Sigma$, let
$\pi:X\rightarrow X/X_{\Gamma_K} \cong X(A)$, and let $\overline{\Sigma}:=
\pi( \Sigma-\Sigma_{\Gamma_K})$ denote the set of restricted roots
relative to $A$. By Tate-Nakayama \cite[p. 702]{La79},
the character $\kappa$ of ${\mathfrak D}(T/F)\subset H^1(\Gamma_K,T)$
may be regarded as a character of the ``cocycle sublattice''
$X_{\Gamma_K}^{\vee}$, trivial on the ``coboundary sublattice'',
\[
\partial X^{\vee}:=\{\chi^{\vee}\in X^{\vee}\ |\
\chi^{\vee}=\sum_{\sigma\in\Gamma_K}\mu_\sigma^\sigma-\mu_\sigma,
\text{ for some }\mu_\sigma\in X^{\vee}\}.
\]

Can the action of $\Gamma$ on $T$ be realized as
an inner automorphism (defined over $\overline{F}$)? By a theorem of
Chevalley, any automorphism
$\sigma\in\Gamma$ may be factored into an inner automorphism and
an outer automorphism of $G_2$ contained in
$Aut(X,\Delta)$, where $\Delta$ denotes the set of simple positive
roots of $G_2$ with respect to $T/\overline{F}$. Since $G_2$ has no outer
automorphisms, it follows that the answer to our question is yes.

By the previous paragraph, the action of $\Gamma$ on $T$ must factor
through the absolute Weyl group $W$. The only Cartan which had
a $\kappa$-stable germ which did not vanish identically when restricted
to $T$ was the following case:

\noindent
Suppose that the non-trivial action of $\Gamma$ on $T$ is realized
by (a) simple reflections $<s_\alpha,s_\zeta>$, (b) the cyclic group
of order $6$, $<s_\beta s_\alpha>$, or (c) the cyclic group
of order $3$, $<(s_\beta s_\alpha)^2>$.
In the first case, $T$ splits over a
biquadratic extension $E/F$ and $Gal(E/F)$ acts on the character lattice
by $<s_\alpha,s_\zeta>$. Moreover, $T=T_G$ is isomorphic over $F$
to $T_H=(T_1\times T_2)/\{\pm 1\}
\hookrightarrow (SL(2,F)\times SL(2,F))/\{\pm 1\}$.
In this case,
\[
H^1(\Gamma,T)=(H^1(\Gamma,T_1)\times H^1(\Gamma,T_2))/\{\pm 1\}
\cong {\mathbb Z}/2{\mathbb Z},
\]
so the character $\kappa$ must be quadratic.

In the second case, note $s_\beta s_\alpha$ acts by a counterclockwise
rotation on the lattice $X_\Gamma^{\vee}$:
\[
s_\beta s_\alpha (m\alpha^{\vee}+n\beta^{\vee})=
(2m-n)\alpha^{\vee}+(m-n)\beta^{\vee}.
\]
Since
\[
\sum_{i=0}^5(s_\beta s_\alpha)^i (m\alpha^{\vee}+n\beta^{\vee})=0,
\]
we have $X_\Gamma^{\vee}=X^{\vee}$ (this also remains true if we replace
$s_\beta s_\alpha$ by $(s_\beta s_\alpha)^2$). Since
\[
s_\beta s_\alpha (-2\alpha^{\vee}-\beta^{\vee})-(-2\alpha^{\vee}-\beta^{\vee})
=\alpha^{\vee},
\]
\[
s_\beta s_\alpha (-3\alpha^{\vee}-2\beta^{\vee})-
(-3\alpha^{\vee}-2\beta^{\vee})=
\beta^{\vee},
\]
we have $\partial X^{\vee}=X^{\vee}$. Therefore, $H^1(\Gamma,T)$ is trivial by
Tate-Nakayama duality.

In the cubic case,
\[
s_\beta s_\alpha (m\alpha^{\vee}+n\beta^{\vee})=
(-3n)\alpha^{\vee}+(m-3n)\beta^{\vee}.
\]
Therefore, in this case,
\[
H^1(\Gamma,T)\cong {\mathbb Z}/3{\mathbb Z},
\]
and $\kappa$ must be a character of order $3$.

All other non-split Cartans of $G_2$ do not contribute to the support
of the $2$-regular and $3$-regular $\kappa$-stable germs of $G_2$.

\subsection*{A.7. The stability of the $2$-regular and $3$-regular
unipotent orbits for $G_2$}

I thank M. Assem for explaining to me the technique of this section.

Using the appendices above, we find that
\[
{\mathfrak m}:=cent_{{\mathfrak g}_2}(h_{\alpha +\beta})
={\mathfrak t}+F\cdot X_{3\alpha +\beta}+F\cdot X_{-3\alpha -\beta},
\]
where ${\mathfrak t}$ denotes the Cartan of ${\mathfrak g}_2$. Let
$M:=Cent_{G_2}(h_{\alpha +\beta})$ and let
\[
{\mathfrak g}_2(i):=\{X\in{\mathfrak g}_2\ |\ [h_{\alpha +\beta},X]=iX\}.
\]
We compute
\[
{\mathfrak g}_2(2)=F\cdot X_{\alpha +\beta}.
\]
The map $m\longmapsto Ad(m)X_{\alpha +\beta}$ defines a submersion
$M\rightarrow {\mathfrak g}_2(2)$ \cite[ch III, \S 4.14]{SS86}. By a lemma of
Rao \cite{Ra72}, the $G(F)$-orbit of $X_{\alpha +\beta}$ is
$Ad(K(V(X_{\alpha +\beta})+{\mathfrak n}(2)))$, where
\[
V(X_{\alpha +\beta})=\{ad(X)X_{\alpha +\beta}\ |\ X\in {\mathfrak m}\},
\]
$K$ is a maximal compact subgroup of $G_2(F)$, and
\[
{\mathfrak n}(2):=\sum_{i >2}{\mathfrak g}_2(i).
\]
The fact that $dim\ {\mathfrak g}_2(2)=1$ gives us that orbit
of $M$ in ${\mathfrak g}_2(2)$ is onto,
so by Rao's result the $2$-regular unipotent orbit is stable.

A similar calculation shows that the $3$-regular orbit is stable.
(In fact, replacing $h_{\alpha +\beta}$ and $X_{\alpha +\beta}$
by $h_{\beta}$ and $X_{\beta}$, we compute
\[
{\mathfrak g}_2(2)=F\cdot X_{\beta},
\]
and the identical argument holds.)

A Lie triple for the subregular orbit is
\[
X_+=X_\beta+X_{2\alpha +\beta},
X_-=X_{-\beta}+X_{-2\alpha -\beta},
H=h_{\beta}+h_{2\alpha +\beta},
\]
and
\[
{\mathfrak g}_2(2)=
F\cdot X_{\alpha +\beta}+F\cdot X_{2\alpha +\beta}
+F\cdot X_{3\alpha +\beta}+F\cdot X_{\beta}.
\]
Direct matrix calculations show that the subregular orbit is
not stable. M. Assem has shown that there are 7 $F$-classes in
the stable class, $p\not= 2,3$.

% Claude: UNTYPESET. The following text follows \enddocument in g2newapp.tex
% and was never typeset. Its heading repeats the number A.6. The label [D]
% does not occur in the bibliography of the source and is typeset literally.
\iffalse
\subsection*{A.6. Weighted Dynkin diagrams for $G_2$}

Let $e_+,\ e_-, \ f$ be elements of ${\mathfrak g}_2$ satisfying
$[e_+,e_-]=f,[f,e_-]=-2e_-$, and $[f,e_+]=2e_+$. Of course,
this triple gives rise to a 3-dimensional subalgebra ${\mathfrak g}$
of ${\mathfrak g}_2$. Call two such subalgebras ${\mathfrak g}$,
${\mathfrak g}'$ \textbf{equivalent} if there is an automorphism
$\phi\in Aut({\mathfrak g}_2)$ for which ${\mathfrak g}'=\phi({\mathfrak g})$.
Following [D, \S 8], we may attach to
such a subalgebra ${\mathfrak g}$ a weighted diagram depending only on the
equivalence class of ${\mathfrak g}$. They are constructed as follows:
let
\[
f:= \begin{pmatrix} 5/3&&&&&&& \\
&4/3&&&&&& \\
&&1/3&&&&& \\
&&&0&&&& \\
&&&&0&&& \\
&&&&&-1/3&& \\
&&&&&&-4/3& \\
&&&&&&&-5/3
\end{pmatrix}.
\]
Then $f=\frac{5}{3}h_\alpha +3h_\beta$. In general, if
$e_+=c_\alpha X_\alpha +c_\beta X_\beta$ then
$e_-=c_\alpha X_{-\alpha} +c_\beta X_{-\beta}$ and
$f=2c_\alpha^2 h_\alpha +2c_\beta^2 h_\beta$. From this one can compute
$e_+$ and $e_-$. The weighted diagram has two nodes connected by
a triple bond, the $\alpha$ node having weight $(f,h_\alpha)=2$ and the
$\beta$ node having weight $(f,h_\beta)=2$.
It corresponds to the regular unipotent class.

Let
\[
f:= \begin{pmatrix} 2/3&&&&&&& \\
&1/3&&&&&& \\
&&1/3&&&&& \\
&&&0&&&& \\
&&&&0&&& \\
&&&&&-1/3&& \\
&&&&&&-1/3& \\
&&&&&&&-2/3
\end{pmatrix}.
\]
Then $f=\frac{2}{3}h_\alpha +h_\beta$. From this one can compute
$e_+$ and $e_-$. The weighted diagram has two nodes connected by
a triple bond, the $\alpha$ node having weight $2$ and the
$\beta$ node having weight $0$. It corresponds to the subregular unipotent
class.

Let
\[
f:= \begin{pmatrix} 1/2&&&&&&& \\
&1/2&&&&&& \\
&&0&&&&& \\
&&&0&&&& \\
&&&&0&&& \\
&&&&&0&& \\
&&&&&&-1/2& \\
&&&&&&&-1/2
\end{pmatrix}.
\]
Then $f=\frac{1}{2}h_\alpha +h_\beta$. From this one can compute
$e_+$ and $e_-$. The weighted diagram has two nodes connected by
a triple bond, the $\alpha$ node having weight $0$ and the
$\beta$ node having weight $1$. It corresponds to the 2-regular unipotent
class.

Let
\[
f:= \begin{pmatrix} 1/3&&&&&&& \\
&1/6&&&&&& \\
&&1/6&&&&& \\
&&&0&&&& \\
&&&&0&&& \\
&&&&&-1/6&& \\
&&&&&&-1/6& \\
&&&&&&&-1/3
\end{pmatrix}.
\]
Then $f=\frac{1}{3}h_\alpha +\frac{1}{2}h_\beta$. From this one can compute
$e_+$ and $e_-$. The weighted diagram has two nodes connected by
a triple bond, the $\alpha$ node having weight $1$ and the
$\beta$ node having weight $0$. It corresponds to the 3-regular unipotent
class.
\fi
% ======================================================================
% Claude: References, source file g2newref.tex (\pageno 64).
% Printed labels are those of the source. Underlined titles in the source
% (\underbar) are set in italics here. Entries are in the source's order.
% Claude: Labels cited in the text but missing here: [Sp3] (Section 1),
% [Sp] (Section 3), [A], [H4] (g2new9d.tex), [D] (untypeset appendix text).
% Claude: [Ho], [Rep'] and [Sp1] are never cited; [SS] and [Rep] are cited
% only in untypeset text.
% ======================================================================
\begin{thebibliography}{SpSt}

\bibitem[Bor]{Bor91} A. Borel, \textit{Linear Algebraic Groups}, 2nd ed,
Springer-Verlag, 1991

\bibitem[Car]{Ca89} R. Carter, ``Simple groups of Lie type'',
Wiley-Interscience, 1989

\bibitem[DG]{DG80} M. Demazure, P. Gabriel, \textit{Introduction
to Algebraic Geometry and Algebraic Groups},
North-Holland Math. Studies, vol. 39, 1980

\bibitem[H1]{Ha92a} T. Hales, \textit{The Subregular Germ of Orbital Integrals},
Memoirs AMS (1992)

\bibitem[H2]{Ha89} ------, ``Shalika germs on $GSp(4)$'', in
\textit{Orbites Unipotentes et Repr\'esentations}, Ast\'erisque
\textit{171-172}(1989)

\bibitem[H3]{Ha92b} ------, ``Orbital integrals on $U(3)$'', in
\textit{The Zeta Functions of Picard Modular Surfaces},
ed. R. Langlands, D. Ramakrishnan, Centre de Recherches
Math., Univ. of Montr\'eal, 1992

\bibitem[Har]{Har77} R. Hartshorne, \textit{Algebraic Geometry},
Springer-Verlag

\bibitem[Ho]{Ho72} M. Hochster, ``Rings of invariants of
tori, Cohen-Macauley rings generated by monomials, and polytopes'',
Ann. Math. \textit{96}(1972)318-337

\bibitem[I]{Ig78} J.-I. Igusa, \textit{Lectures on Forms of Higher Degree},
Tata Institute, Bombay, 1978

\bibitem[Jac]{Ja71} N. Jacobson, \textit{Exceptional Lie Algebras},
Marcel-Dekker, 1971

\bibitem[K]{Ke72} G. Kempf et al, \textit{Toroidal Embeddings, I},
Springer LNM \textit{339}, 1972

\bibitem[KP]{KP79} H. Kraft, C. Procesi, ``Closures of conjugacy
classes of matrices are normal'', Inv. Math.
\textit{53}(1979)227-247

\bibitem[L]{La83} R. Langlands, ``Orbital integrals on forms of
$SL(3)$, I'', Amer J Math \textit{105}(1983)

\bibitem[L$'$]{La79} ------, ``Stable conjugacy: definitions and lemmas'',
Can. J. Math. \textit{31}(1979)700-725

\bibitem[LS0]{LS83} ------ and D. Shelstad, ``On principal values on
p-adic manifolds'' in \textit{Lie Group Rep., II},
SLN Math \textit{1041}(1983)

\bibitem[LS1]{LS87} ------, ``On the definition of transfer factors'',
Math. Ann. \textit{278}(1987)219-271

\bibitem[LS2]{LS90} ------, ``Descent for transfer factors'', in
\textit{Grothendieck Festschrift}, vol. II,
Birkhauser, 1990

\bibitem[LS3]{LS89} ------, ``Orbital integrals on forms of
$SL(3)$, II'', Can. J. of Math. \textit{41}(1989)480-507

\bibitem[M]{Mu88} D. Mumford,
\textit{The Red Book of Varieties and Schemes}, Springer LNM
\textit{1358}(1988)

\bibitem[R]{Ra72} R. Ranga Rao, ``Orbital integrals in reductive groups'',
Ann. Math. \textit{96}(1972)503-510

\bibitem[Rep]{Re84} J. Repka, ``Shalika's germs for p-adic $GL(n)$, II:
the subregular term'', Pac. J. Math. \textit{113}(1984)173-182

\bibitem[Rep$'$]{Re85} ------, ``Germs associated to regular unipotent
classes in p-adic $SL(n)$'', Can. Math. Bull. \textit{28}(1985)257-266

\bibitem[Ro]{Ro80} J. Rogawski, ``Applications of the building to orbital
integrals'', Thesis, Princeton U., 1980

\bibitem[SS]{SS83} P. Sally and J. Shalika, ``The Fourier transform of orbital
integrals on $SL(2)$ over a p-adic field'', in
\textit{Lie Group Representations II},
ed. R. Herb et al, Springer LNM \textit{1041}(1983)303-340

\bibitem[Sat]{Sa71} I. Satake, \textit{Classification Theory
of Semi-simple Algebraic Groups}, Dekker, 1971

\bibitem[Sh]{Sh72} J. Shalika, ``A theorem on semi-simple p-adic groups'',
Ann. Math. \textit{95}(1972)226-242

\bibitem[S1]{She82} D. Shelstad, ``L-indistinguishability for real
groups'', Math. Ann. \textit{259}(1982)385-430

\bibitem[S2]{She89} ------, ``A formula for regular germs'',
in \textit{Orbites Unipotentes et Repr\'esentations},
Ast\'erisque, \textit{171-172}(1989)

\bibitem[Slo]{Sl80} P. Slodowy, \textit{Simple singularities
and simple algebraic groups}, Springer LNM 815, 1980

\bibitem[Spal]{Spa82} N. Spaltenstein, \textit{Classes
Unipotentes et Sous-groupes de Borel},
Springer LNM 946, 1982

\bibitem[Sp1]{Sp69} T. A. Springer, ``The unipotent variety of a semi-simple
group'', International Coll. on Alg. Geometry, Bombay, 1968,
Oxford Univ Press, 1969

\bibitem[Sp2]{Sp79} -------, ``Reductive groups'', in
\textit{Automomorphic Forms, Represenattions, and L-functions},
Part I, Proc. Symp. Pure Math. vol 33,
AMS, 1979

\bibitem[SpSt]{SS86} T. A. Springer and R. Steinberg, ``Conjugacy Classes'',
in \textit{Seminar on Algebraic Groups}, Springer-Verlag LNM 131, 1986(2nd ed.)

\bibitem[St]{St74} R. Steinberg, \textit{Conjugacy Classes
in Algebraic Groups},
Springer LNM 366, 1974

\bibitem[T]{Ti65} Tits, ``Classification of algebraic semi-simple groups'',
in \textit{Algebraic Groups and Discontinuous Groups}
(Boulder conf., 1965), Proc. Symp. Pure Math., vol. 9, AMS

\end{thebibliography}

\medskip
\noindent
D. Joyner, Math Dept, USNA, Annapolis, MD 21402

\end{document}
